\documentclass[10pt]{report}
\usepackage[slovak]{babel}
\usepackage{lmodern}
\usepackage{graphicx}
\usepackage{varwidth}
\usepackage{enumitem}
\usepackage{amsmath}
\usepackage{amssymb}
\usepackage{mathtools}
\usepackage{pifont}
\usepackage{arydshln}
\usepackage{tikz}
\usepackage{lscape}
\usepackage{bm}
\usepackage{nicefrac}
\usepackage{multicol}
\usepackage{vwcol}
\usepackage{physics}
\usepackage{fontsize}
\usepackage[landscape]{geometry}
\geometry{a4paper, total={296mm,210mm}, left=-5mm, top=0mm}
\renewcommand{\labelenumi}{\bfseries(\alph{enumi})\phantom{x}}
\newcommand\omicron{o}
\hfuzz=50pt
\setlist[enumerate]{leftmargin=0pt,itemindent=34pt}
\pagenumbering{gobble}
\setlength{\tabcolsep}{0pt}
\begin{document}
\thispagestyle{empty}
\begin{tabular}{c:c}
\begin{minipage}[c][104.5mm][t]{0.5\linewidth}
\begin{center}
\vspace{7mm}
{\huge \LARGE Funkcia dvoch premenných, skupina \textit{Alpha $\alpha$} -\romannumeral1}\\[5mm]
\begin{minipage}{0.95\linewidth}
\begin{center}
{\small \begin{multicols}{2} \raggedright
\textbf{(a)} Vypočítajte zadanú parciálnu deriváciu funkcie \(f(x,y)\). \\[1mm]
\textbf{(b)} Určte gradient \(f(x,y)\) v bode \(A\). \\[1mm]
\textbf{(c)} Napíšte rovnicu dotykovej roviny ku grafu funkcie \(g(x,y)\) v dotykovom bode \(D\). \par \columnbreak
\textbf{(d)} Pomocou lineárnej aproximácie odhadnite funkčnú hodnotu funkcie \(g(x,y)\) v bode, ktorý je posunutý o vektor \((\Delta x, \Delta y)\) od bodu \(D_0\). \\[1mm]
\textbf{(e)} Vypočítajte \(h(x,y)\) v jej lokálnom extréme. \\[1mm]
\textbf{(f)} Nájdite bod viazeného lokálneho extrému funkcie \(h(x,y)\) na priamke. \\[1mm]
\end{multicols}
}
\end{center}
\end{minipage}
\\[4mm]
\begin{minipage}{0.79\linewidth}
\begin{center}
\begin{varwidth}{\linewidth}
\begin{enumerate}
\normalsizerr
\item $f(x,y) = \sqrt[3]{8x^2y + 9xy^2}\enspace , \enspace \pdv{f}{y}$\quad \dotfill\; ???\;\dotfill \quad $\frac{8x^2+18xy}{3\,\sqrt[3]{(8x^2y+9xy^2)^{2}}}$
\item $A = [1, 1]$\quad \dotfill\; ???\;\dotfill \quad $\left( \frac{50}{3\,\sqrt[3]{17^2}} \;,\; \frac{26}{3\,\sqrt[3]{17^2}} \right)$
\item $g(x,y) = \sin\frac{x}{y} \enspace , \enspace D = [0, -3, ?]$\quad \dotfill\; ???\;\dotfill \quad $x+3z=0$
\item $\Delta x = {0.02} \enspace,\enspace \Delta y = {-0.3}$\quad \dotfill\; ???\;\dotfill \quad $-0.06$
\item $h(x,y) = 3x^2+5y^2-4xy+x$\quad \dotfill\; ???\;\dotfill \quad $\nicefrac{-5}{76}$
\item $y = x-7$\quad \dotfill\; ???\;\dotfill \quad $\left[\frac{41}{8}\;,\;\frac{-15}{8}\right]$
\end{enumerate}
\end{varwidth}
\end{center}
\end{minipage}
\begin{minipage}{0.20\linewidth}
\begin{center}
{\Huge\bfseries 1.} \\[2mm]
\includegraphics[height=37mm]{images/braille.png}
{\small Braillove písmeno}
\end{center}
\end{minipage}
\end{center}
\end{minipage}
&
\begin{minipage}[c][104.5mm][t]{0.5\linewidth}
\begin{center}
\vspace{7mm}
{\huge \LARGE Funkcia dvoch premenných, skupina \textit{Alpha $\alpha$} -\romannumeral2}\\[5mm]
\begin{minipage}{0.95\linewidth}
\begin{center}
{\small \begin{multicols}{2} \raggedright
\textbf{(a)} Vypočítajte zadanú parciálnu deriváciu funkcie \(f(x,y)\). \\[1mm]
\textbf{(b)} Určte gradient \(f(x,y)\) v bode \(A\). \\[1mm]
\textbf{(c)} Napíšte rovnicu dotykovej roviny ku grafu funkcie \(g(x,y)\) v dotykovom bode \(D\). \par \columnbreak
\textbf{(d)} Pomocou lineárnej aproximácie odhadnite funkčnú hodnotu funkcie \(g(x,y)\) v bode, ktorý je posunutý o vektor \((\Delta x, \Delta y)\) od bodu \(D_0\). \\[1mm]
\textbf{(e)} Vypočítajte \(h(x,y)\) v jej lokálnom extréme. \\[1mm]
\textbf{(f)} Nájdite bod viazeného lokálneho extrému funkcie \(h(x,y)\) na priamke. \\[1mm]
\end{multicols}
}
\end{center}
\end{minipage}
\\[4mm]
\begin{minipage}{0.79\linewidth}
\begin{center}
\begin{varwidth}{\linewidth}
\begin{enumerate}
\normalsizerr
\item $f(x,y) = \sqrt[4]{2x^2y + 9xy^2}\enspace , \enspace \pdv{f}{x}$\quad \dotfill\; ???\;\dotfill \quad $\frac{4xy+9y^2}{4\,\sqrt[4]{(2x^2y+9xy^2)^{3}}}$
\item $A = [3, 3]$\quad \dotfill\; ???\;\dotfill \quad $\left( \frac{117}{4\,\sqrt[4]{297^3}} \;,\; \frac{360}{4\,\sqrt[4]{297^3}} \right)$
\item $g(x,y) = \ln\frac{x}{y} \enspace , \enspace D=[e,1,?]$\quad \dotfill\; ???\;\dotfill \quad $\frac{x}{e} + y + z + 1 = 0$
\item $\Delta x = {-0.01} \enspace,\enspace \Delta y = {0.02}$\quad \dotfill\; ???\;\dotfill \quad $1.006$
\item $h(x,y) = -x^2-4y^2-3xy-x$\quad \dotfill\; ???\;\dotfill \quad $\nicefrac{4}{25}$
\item $y = -2x-3$\quad \dotfill\; ???\;\dotfill \quad $\left[\frac{28}{11}\;,\;\frac{7}{11}\right]$
\end{enumerate}
\end{varwidth}
\end{center}
\end{minipage}
\begin{minipage}{0.20\linewidth}
\begin{center}
{\Huge\bfseries 2.} \\[2mm]
\includegraphics[height=37mm]{images/braille.png}
{\small Braillove písmeno}
\end{center}
\end{minipage}
\end{center}
\end{minipage}
\\ \hdashline
\begin{minipage}[c][104.5mm][t]{0.5\linewidth}
\begin{center}
\vspace{7mm}
{\huge \LARGE Funkcia dvoch premenných, skupina \textit{Alpha $\alpha$} -\romannumeral3}\\[5mm]
\begin{minipage}{0.95\linewidth}
\begin{center}
{\small \begin{multicols}{2} \raggedright
\textbf{(a)} Vypočítajte zadanú parciálnu deriváciu funkcie \(f(x,y)\). \\[1mm]
\textbf{(b)} Určte gradient \(f(x,y)\) v bode \(A\). \\[1mm]
\textbf{(c)} Napíšte rovnicu dotykovej roviny ku grafu funkcie \(g(x,y)\) v dotykovom bode \(D\). \par \columnbreak
\textbf{(d)} Pomocou lineárnej aproximácie odhadnite funkčnú hodnotu funkcie \(g(x,y)\) v bode, ktorý je posunutý o vektor \((\Delta x, \Delta y)\) od bodu \(D_0\). \\[1mm]
\textbf{(e)} Vypočítajte \(h(x,y)\) v jej lokálnom extréme. \\[1mm]
\textbf{(f)} Nájdite bod viazeného lokálneho extrému funkcie \(h(x,y)\) na priamke. \\[1mm]
\end{multicols}
}
\end{center}
\end{minipage}
\\[4mm]
\begin{minipage}{0.79\linewidth}
\begin{center}
\begin{varwidth}{\linewidth}
\begin{enumerate}
\normalsizerr
\item $f(x,y) = \sqrt[4]{4x^2y + 4xy^2}\enspace , \enspace \pdv{f}{x}$\quad \dotfill\; ???\;\dotfill \quad $\frac{8xy+4y^2}{4\,\sqrt[4]{(4x^2y+4xy^2)^{3}}}$
\item $A = [2, 2]$\quad \dotfill\; ???\;\dotfill \quad $\left( \frac{48}{4\,\sqrt[4]{64^3}} \;,\; \frac{48}{4\,\sqrt[4]{64^3}} \right)$
\item $g(x,y) = \sin\frac{x}{y} \enspace , \enspace D = [0, -1, ?]$\quad \dotfill\; ???\;\dotfill \quad $x+z=0$
\item $\Delta x = {0.3} \enspace,\enspace \Delta y = {-0.3}$\quad \dotfill\; ???\;\dotfill \quad $0.3$
\item $h(x,y) = 5x^2+2y^2+5xy+4x$\quad \dotfill\; ???\;\dotfill \quad $\nicefrac{-31}{15}$
\item $y = -2x-1$\quad \dotfill\; ???\;\dotfill \quad $\left[\frac{-7}{6}\;,\;\frac{4}{3}\right]$
\end{enumerate}
\end{varwidth}
\end{center}
\end{minipage}
\begin{minipage}{0.20\linewidth}
\begin{center}
{\Huge\bfseries 3.} \\[2mm]
\includegraphics[height=37mm]{images/braille.png}
{\small Braillove písmeno}
\end{center}
\end{minipage}
\end{center}
\end{minipage}
&
\begin{minipage}[c][104.5mm][t]{0.5\linewidth}
\begin{center}
\vspace{7mm}
{\huge \LARGE Funkcia dvoch premenných, skupina \textit{Alpha $\alpha$} -\romannumeral4}\\[5mm]
\begin{minipage}{0.95\linewidth}
\begin{center}
{\small \begin{multicols}{2} \raggedright
\textbf{(a)} Vypočítajte zadanú parciálnu deriváciu funkcie \(f(x,y)\). \\[1mm]
\textbf{(b)} Určte gradient \(f(x,y)\) v bode \(A\). \\[1mm]
\textbf{(c)} Napíšte rovnicu dotykovej roviny ku grafu funkcie \(g(x,y)\) v dotykovom bode \(D\). \par \columnbreak
\textbf{(d)} Pomocou lineárnej aproximácie odhadnite funkčnú hodnotu funkcie \(g(x,y)\) v bode, ktorý je posunutý o vektor \((\Delta x, \Delta y)\) od bodu \(D_0\). \\[1mm]
\textbf{(e)} Vypočítajte \(h(x,y)\) v jej lokálnom extréme. \\[1mm]
\textbf{(f)} Nájdite bod viazeného lokálneho extrému funkcie \(h(x,y)\) na priamke. \\[1mm]
\end{multicols}
}
\end{center}
\end{minipage}
\\[4mm]
\begin{minipage}{0.79\linewidth}
\begin{center}
\begin{varwidth}{\linewidth}
\begin{enumerate}
\normalsizerr
\item $f(x,y) = \sqrt[3]{7x^2y + 3xy^2}\enspace , \enspace \pdv{f}{y}$\quad \dotfill\; ???\;\dotfill \quad $\frac{7x^2+6xy}{3\,\sqrt[3]{(7x^2y+3xy^2)^{2}}}$
\item $A = [4, 1]$\quad \dotfill\; ???\;\dotfill \quad $\left( \frac{59}{3\,\sqrt[3]{124^2}} \;,\; \frac{272}{3\,\sqrt[3]{124^2}} \right)$
\item $g(x,y) = \ln\frac{x}{y} \enspace , \enspace D=[e^{3},1,?]$\quad \dotfill\; ???\;\dotfill \quad $\frac{x}{e^3} + y + z + 3 = 0$
\item $\Delta x = {-0.3} \enspace,\enspace \Delta y = {0.3}$\quad \dotfill\; ???\;\dotfill \quad $3.285$
\item $h(x,y) = -5x^2-4y^2+3xy-x$\quad \dotfill\; ???\;\dotfill \quad $\nicefrac{4}{89}$
\item $y = -x+3$\quad \dotfill\; ???\;\dotfill \quad $\left[\frac{4}{3}\;,\;\frac{10}{3}\right]$
\end{enumerate}
\end{varwidth}
\end{center}
\end{minipage}
\begin{minipage}{0.20\linewidth}
\begin{center}
{\Huge\bfseries 4.} \\[2mm]
\includegraphics[height=37mm]{images/braille.png}
{\small Braillove písmeno}
\end{center}
\end{minipage}
\end{center}
\end{minipage}
%
\end{tabular}
\newpage
\thispagestyle{empty}
\begin{tabular}{c:c}
\begin{minipage}[c][104.5mm][t]{0.5\linewidth}
\begin{center}
\vspace{7mm}
{\huge \LARGE Funkcia dvoch premenných, skupina \textit{Beta $\beta$} -\romannumeral1}\\[5mm]
\begin{minipage}{0.95\linewidth}
\begin{center}
{\small \begin{multicols}{2} \raggedright
\textbf{(a)} Vypočítajte zadanú parciálnu deriváciu funkcie \(f(x,y)\). \\[1mm]
\textbf{(b)} Určte gradient \(f(x,y)\) v bode \(A\). \\[1mm]
\textbf{(c)} Napíšte rovnicu dotykovej roviny ku grafu funkcie \(g(x,y)\) v dotykovom bode \(D\). \par \columnbreak
\textbf{(d)} Pomocou lineárnej aproximácie odhadnite funkčnú hodnotu funkcie \(g(x,y)\) v bode, ktorý je posunutý o vektor \((\Delta x, \Delta y)\) od bodu \(D_0\). \\[1mm]
\textbf{(e)} Vypočítajte \(h(x,y)\) v jej lokálnom extréme. \\[1mm]
\textbf{(f)} Nájdite bod viazeného lokálneho extrému funkcie \(h(x,y)\) na priamke. \\[1mm]
\end{multicols}
}
\end{center}
\end{minipage}
\\[4mm]
\begin{minipage}{0.79\linewidth}
\begin{center}
\begin{varwidth}{\linewidth}
\begin{enumerate}
\normalsizerr
\item $f(x,y) = \sqrt[4]{x^2y + 5xy^2}\enspace , \enspace \pdv{f}{y}$\quad \dotfill\; ???\;\dotfill \quad $\frac{x^2+10xy}{4\,\sqrt[4]{(x^2y+5xy^2)^{3}}}$
\item $A = [4, 3]$\quad \dotfill\; ???\;\dotfill \quad $\left( \frac{69}{4\,\sqrt[4]{228^3}} \;,\; \frac{136}{4\,\sqrt[4]{228^3}} \right)$
\item $g(x,y) = \ln\frac{x}{y} \enspace , \enspace D=[e^{3},1,?]$\quad \dotfill\; ???\;\dotfill \quad $\frac{x}{e^3} - y - z + 3 = 0$
\item $\Delta x = {0.2} \enspace,\enspace \Delta y = {-0.01}$\quad \dotfill\; ???\;\dotfill \quad $3.02$
\item $h(x,y) = -4x^2-4y^2+xy+x$\quad \dotfill\; ???\;\dotfill \quad $\nicefrac{5}{63}$
\item $y = x-1$\quad \dotfill\; ???\;\dotfill \quad $\left[\frac{4}{7}\;,\;\frac{-1}{3}\right]$
\end{enumerate}
\end{varwidth}
\end{center}
\end{minipage}
\begin{minipage}{0.20\linewidth}
\begin{center}
{\Huge\bfseries 1.} \\[2mm]
\includegraphics[height=37mm]{images/braille.png}
{\small Braillove písmeno}
\end{center}
\end{minipage}
\end{center}
\end{minipage}
&
\begin{minipage}[c][104.5mm][t]{0.5\linewidth}
\begin{center}
\vspace{7mm}
{\huge \LARGE Funkcia dvoch premenných, skupina \textit{Beta $\beta$} -\romannumeral2}\\[5mm]
\begin{minipage}{0.95\linewidth}
\begin{center}
{\small \begin{multicols}{2} \raggedright
\textbf{(a)} Vypočítajte zadanú parciálnu deriváciu funkcie \(f(x,y)\). \\[1mm]
\textbf{(b)} Určte gradient \(f(x,y)\) v bode \(A\). \\[1mm]
\textbf{(c)} Napíšte rovnicu dotykovej roviny ku grafu funkcie \(g(x,y)\) v dotykovom bode \(D\). \par \columnbreak
\textbf{(d)} Pomocou lineárnej aproximácie odhadnite funkčnú hodnotu funkcie \(g(x,y)\) v bode, ktorý je posunutý o vektor \((\Delta x, \Delta y)\) od bodu \(D_0\). \\[1mm]
\textbf{(e)} Vypočítajte \(h(x,y)\) v jej lokálnom extréme. \\[1mm]
\textbf{(f)} Nájdite bod viazeného lokálneho extrému funkcie \(h(x,y)\) na priamke. \\[1mm]
\end{multicols}
}
\end{center}
\end{minipage}
\\[4mm]
\begin{minipage}{0.79\linewidth}
\begin{center}
\begin{varwidth}{\linewidth}
\begin{enumerate}
\normalsizerr
\item $f(x,y) = \sqrt[3]{7x^2y + 2xy^2}\enspace , \enspace \pdv{f}{x}$\quad \dotfill\; ???\;\dotfill \quad $\frac{14xy+2y^2}{3\,\sqrt[3]{(7x^2y+2xy^2)^{2}}}$
\item $A = [2, 3]$\quad \dotfill\; ???\;\dotfill \quad $\left( \frac{204}{3\,\sqrt[3]{120^2}} \;,\; \frac{52}{3\,\sqrt[3]{120^2}} \right)$
\item $g(x,y) = \sin\frac{x}{y} \enspace , \enspace D = [0, -1, ?]$\quad \dotfill\; ???\;\dotfill \quad $x-z=0$
\item $\Delta x = {0.03} \enspace,\enspace \Delta y = {0.03}$\quad \dotfill\; ???\;\dotfill \quad $-0.03$
\item $h(x,y) = -x^2-3y^2-2xy+5x$\quad \dotfill\; ???\;\dotfill \quad $\nicefrac{75}{16}$
\item $y = 4x+2$\quad \dotfill\; ???\;\dotfill \quad $\left[\frac{-47}{114}\;,\;\frac{20}{41}\right]$
\end{enumerate}
\end{varwidth}
\end{center}
\end{minipage}
\begin{minipage}{0.20\linewidth}
\begin{center}
{\Huge\bfseries 2.} \\[2mm]
\includegraphics[height=37mm]{images/braille.png}
{\small Braillove písmeno}
\end{center}
\end{minipage}
\end{center}
\end{minipage}
\\ \hdashline
\begin{minipage}[c][104.5mm][t]{0.5\linewidth}
\begin{center}
\vspace{7mm}
{\huge \LARGE Funkcia dvoch premenných, skupina \textit{Beta $\beta$} -\romannumeral3}\\[5mm]
\begin{minipage}{0.95\linewidth}
\begin{center}
{\small \begin{multicols}{2} \raggedright
\textbf{(a)} Vypočítajte zadanú parciálnu deriváciu funkcie \(f(x,y)\). \\[1mm]
\textbf{(b)} Určte gradient \(f(x,y)\) v bode \(A\). \\[1mm]
\textbf{(c)} Napíšte rovnicu dotykovej roviny ku grafu funkcie \(g(x,y)\) v dotykovom bode \(D\). \par \columnbreak
\textbf{(d)} Pomocou lineárnej aproximácie odhadnite funkčnú hodnotu funkcie \(g(x,y)\) v bode, ktorý je posunutý o vektor \((\Delta x, \Delta y)\) od bodu \(D_0\). \\[1mm]
\textbf{(e)} Vypočítajte \(h(x,y)\) v jej lokálnom extréme. \\[1mm]
\textbf{(f)} Nájdite bod viazeného lokálneho extrému funkcie \(h(x,y)\) na priamke. \\[1mm]
\end{multicols}
}
\end{center}
\end{minipage}
\\[4mm]
\begin{minipage}{0.79\linewidth}
\begin{center}
\begin{varwidth}{\linewidth}
\begin{enumerate}
\normalsizerr
\item $f(x,y) = \sqrt{5x^2y + xy^2}\enspace , \enspace \pdv{f}{y}$\quad \dotfill\; ???\;\dotfill \quad $\frac{5x^2+2xy}{2\,\sqrt{5x^2y+xy^2}}$
\item $A = [1, 1]$\quad \dotfill\; ???\;\dotfill \quad $\left( \frac{11}{2\,\sqrt{6}} \;,\; \frac{7}{2\,\sqrt{6}} \right)$
\item $g(x,y) = \ln\frac{x}{y} \enspace , \enspace D=[e^{3},1,?]$\quad \dotfill\; ???\;\dotfill \quad $\frac{x}{e^3} + y + z + 3 = 0$
\item $\Delta x = {0.1} \enspace,\enspace \Delta y = {-0.1}$\quad \dotfill\; ???\;\dotfill \quad $2.905$
\item $h(x,y) = 4x^2+3y^2-2xy+x$\quad \dotfill\; ???\;\dotfill \quad $\nicefrac{-3}{44}$
\item $y = -3x-3$\quad \dotfill\; ???\;\dotfill \quad $\left[\frac{47}{74}\;,\;\frac{-39}{74}\right]$
\end{enumerate}
\end{varwidth}
\end{center}
\end{minipage}
\begin{minipage}{0.20\linewidth}
\begin{center}
{\Huge\bfseries 3.} \\[2mm]
\includegraphics[height=37mm]{images/braille.png}
{\small Braillove písmeno}
\end{center}
\end{minipage}
\end{center}
\end{minipage}
&
\begin{minipage}[c][104.5mm][t]{0.5\linewidth}
\begin{center}
\vspace{7mm}
{\huge \LARGE Funkcia dvoch premenných, skupina \textit{Beta $\beta$} -\romannumeral4}\\[5mm]
\begin{minipage}{0.95\linewidth}
\begin{center}
{\small \begin{multicols}{2} \raggedright
\textbf{(a)} Vypočítajte zadanú parciálnu deriváciu funkcie \(f(x,y)\). \\[1mm]
\textbf{(b)} Určte gradient \(f(x,y)\) v bode \(A\). \\[1mm]
\textbf{(c)} Napíšte rovnicu dotykovej roviny ku grafu funkcie \(g(x,y)\) v dotykovom bode \(D\). \par \columnbreak
\textbf{(d)} Pomocou lineárnej aproximácie odhadnite funkčnú hodnotu funkcie \(g(x,y)\) v bode, ktorý je posunutý o vektor \((\Delta x, \Delta y)\) od bodu \(D_0\). \\[1mm]
\textbf{(e)} Vypočítajte \(h(x,y)\) v jej lokálnom extréme. \\[1mm]
\textbf{(f)} Nájdite bod viazeného lokálneho extrému funkcie \(h(x,y)\) na priamke. \\[1mm]
\end{multicols}
}
\end{center}
\end{minipage}
\\[4mm]
\begin{minipage}{0.79\linewidth}
\begin{center}
\begin{varwidth}{\linewidth}
\begin{enumerate}
\normalsizerr
\item $f(x,y) = \sqrt{6x^2y + 8xy^2}\enspace , \enspace \pdv{f}{y}$\quad \dotfill\; ???\;\dotfill \quad $\frac{6x^2+16xy}{2\,\sqrt{6x^2y+8xy^2}}$
\item $A = [3, 3]$\quad \dotfill\; ???\;\dotfill \quad $\left( \frac{180}{2\,\sqrt{378}} \;,\; \frac{396}{2\,\sqrt{378}} \right)$
\item $g(x,y) = \sin\frac{x}{y} \enspace , \enspace D = [0, 2, ?]$\quad \dotfill\; ???\;\dotfill \quad $x-2z=1$
\item $\Delta x = {0.3} \enspace,\enspace \Delta y = {0.3}$\quad \dotfill\; ???\;\dotfill \quad $0.3$
\item $h(x,y) = 3x^2+6y^2+5xy+x$\quad \dotfill\; ???\;\dotfill \quad $\nicefrac{-6}{97}$
\item $y = 2x+2$\quad \dotfill\; ???\;\dotfill \quad $\left[\frac{-59}{74}\;,\;\frac{15}{17}\right]$
\end{enumerate}
\end{varwidth}
\end{center}
\end{minipage}
\begin{minipage}{0.20\linewidth}
\begin{center}
{\Huge\bfseries 4.} \\[2mm]
\includegraphics[height=37mm]{images/braille.png}
{\small Braillove písmeno}
\end{center}
\end{minipage}
\end{center}
\end{minipage}
%
\end{tabular}
\newpage
\thispagestyle{empty}
\begin{tabular}{c:c}
\begin{minipage}[c][104.5mm][t]{0.5\linewidth}
\begin{center}
\vspace{7mm}
{\huge \LARGE Funkcia dvoch premenných, skupina \textit{Gamma $\gamma$} -\romannumeral1}\\[5mm]
\begin{minipage}{0.95\linewidth}
\begin{center}
{\small \begin{multicols}{2} \raggedright
\textbf{(a)} Vypočítajte zadanú parciálnu deriváciu funkcie \(f(x,y)\). \\[1mm]
\textbf{(b)} Určte gradient \(f(x,y)\) v bode \(A\). \\[1mm]
\textbf{(c)} Napíšte rovnicu dotykovej roviny ku grafu funkcie \(g(x,y)\) v dotykovom bode \(D\). \par \columnbreak
\textbf{(d)} Pomocou lineárnej aproximácie odhadnite funkčnú hodnotu funkcie \(g(x,y)\) v bode, ktorý je posunutý o vektor \((\Delta x, \Delta y)\) od bodu \(D_0\). \\[1mm]
\textbf{(e)} Vypočítajte \(h(x,y)\) v jej lokálnom extréme. \\[1mm]
\textbf{(f)} Nájdite bod viazeného lokálneho extrému funkcie \(h(x,y)\) na priamke. \\[1mm]
\end{multicols}
}
\end{center}
\end{minipage}
\\[4mm]
\begin{minipage}{0.79\linewidth}
\begin{center}
\begin{varwidth}{\linewidth}
\begin{enumerate}
\normalsizerr
\item $f(x,y) = \sqrt{7x^2y + 6xy^2}\enspace , \enspace \pdv{f}{x}$\quad \dotfill\; ???\;\dotfill \quad $\frac{14xy+6y^2}{2\,\sqrt{7x^2y+6xy^2}}$
\item $A = [2, 3]$\quad \dotfill\; ???\;\dotfill \quad $\left( \frac{138}{2\,\sqrt{192}} \;,\; \frac{100}{2\,\sqrt{192}} \right)$
\item $g(x,y) = \sin\frac{x}{y} \enspace , \enspace D = [0, 3, ?]$\quad \dotfill\; ???\;\dotfill \quad $x-3z=0$
\item $\Delta x = {-0.03} \enspace,\enspace \Delta y = {0.03}$\quad \dotfill\; ???\;\dotfill \quad $-0.01$
\item $h(x,y) = -x^2-4y^2+xy-4x$\quad \dotfill\; ???\;\dotfill \quad $\nicefrac{64}{17}$
\item $y = 2x-4$\quad \dotfill\; ???\;\dotfill \quad $\left[\frac{28}{15}\;,\;\frac{-4}{19}\right]$
\end{enumerate}
\end{varwidth}
\end{center}
\end{minipage}
\begin{minipage}{0.20\linewidth}
\begin{center}
{\Huge\bfseries 1.} \\[2mm]
\includegraphics[height=37mm]{images/braille.png}
{\small Braillove písmeno}
\end{center}
\end{minipage}
\end{center}
\end{minipage}
&
\begin{minipage}[c][104.5mm][t]{0.5\linewidth}
\begin{center}
\vspace{7mm}
{\huge \LARGE Funkcia dvoch premenných, skupina \textit{Gamma $\gamma$} -\romannumeral2}\\[5mm]
\begin{minipage}{0.95\linewidth}
\begin{center}
{\small \begin{multicols}{2} \raggedright
\textbf{(a)} Vypočítajte zadanú parciálnu deriváciu funkcie \(f(x,y)\). \\[1mm]
\textbf{(b)} Určte gradient \(f(x,y)\) v bode \(A\). \\[1mm]
\textbf{(c)} Napíšte rovnicu dotykovej roviny ku grafu funkcie \(g(x,y)\) v dotykovom bode \(D\). \par \columnbreak
\textbf{(d)} Pomocou lineárnej aproximácie odhadnite funkčnú hodnotu funkcie \(g(x,y)\) v bode, ktorý je posunutý o vektor \((\Delta x, \Delta y)\) od bodu \(D_0\). \\[1mm]
\textbf{(e)} Vypočítajte \(h(x,y)\) v jej lokálnom extréme. \\[1mm]
\textbf{(f)} Nájdite bod viazeného lokálneho extrému funkcie \(h(x,y)\) na priamke. \\[1mm]
\end{multicols}
}
\end{center}
\end{minipage}
\\[4mm]
\begin{minipage}{0.79\linewidth}
\begin{center}
\begin{varwidth}{\linewidth}
\begin{enumerate}
\normalsizerr
\item $f(x,y) = \sqrt[4]{4x^2y + xy^2}\enspace , \enspace \pdv{f}{y}$\quad \dotfill\; ???\;\dotfill \quad $\frac{4x^2+2xy}{4\,\sqrt[4]{(4x^2y+xy^2)^{3}}}$
\item $A = [3, 1]$\quad \dotfill\; ???\;\dotfill \quad $\left( \frac{50}{4\,\sqrt[4]{39^3}} \;,\; \frac{42}{4\,\sqrt[4]{39^3}} \right)$
\item $g(x,y) = \sin\frac{x}{y} \enspace , \enspace D = [0, 2, ?]$\quad \dotfill\; ???\;\dotfill \quad $x+2z=0$
\item $\Delta x = {-0.3} \enspace,\enspace \Delta y = {-0.02}$\quad \dotfill\; ???\;\dotfill \quad $-0.01$
\item $h(x,y) = 7x^2+y^2+5xy-5x$\quad \dotfill\; ???\;\dotfill \quad $-8$
\item $y = -3x-3$\quad \dotfill\; ???\;\dotfill \quad $\left[1\;,\;\frac{-6}{31}\right]$
\end{enumerate}
\end{varwidth}
\end{center}
\end{minipage}
\begin{minipage}{0.20\linewidth}
\begin{center}
{\Huge\bfseries 2.} \\[2mm]
\includegraphics[height=37mm]{images/braille.png}
{\small Braillove písmeno}
\end{center}
\end{minipage}
\end{center}
\end{minipage}
\\ \hdashline
\begin{minipage}[c][104.5mm][t]{0.5\linewidth}
\begin{center}
\vspace{7mm}
{\huge \LARGE Funkcia dvoch premenných, skupina \textit{Gamma $\gamma$} -\romannumeral3}\\[5mm]
\begin{minipage}{0.95\linewidth}
\begin{center}
{\small \begin{multicols}{2} \raggedright
\textbf{(a)} Vypočítajte zadanú parciálnu deriváciu funkcie \(f(x,y)\). \\[1mm]
\textbf{(b)} Určte gradient \(f(x,y)\) v bode \(A\). \\[1mm]
\textbf{(c)} Napíšte rovnicu dotykovej roviny ku grafu funkcie \(g(x,y)\) v dotykovom bode \(D\). \par \columnbreak
\textbf{(d)} Pomocou lineárnej aproximácie odhadnite funkčnú hodnotu funkcie \(g(x,y)\) v bode, ktorý je posunutý o vektor \((\Delta x, \Delta y)\) od bodu \(D_0\). \\[1mm]
\textbf{(e)} Vypočítajte \(h(x,y)\) v jej lokálnom extréme. \\[1mm]
\textbf{(f)} Nájdite bod viazeného lokálneho extrému funkcie \(h(x,y)\) na priamke. \\[1mm]
\end{multicols}
}
\end{center}
\end{minipage}
\\[4mm]
\begin{minipage}{0.79\linewidth}
\begin{center}
\begin{varwidth}{\linewidth}
\begin{enumerate}
\normalsizerr
\item $f(x,y) = \sqrt[4]{3x^2y + 5xy^2}\enspace , \enspace \pdv{f}{y}$\quad \dotfill\; ???\;\dotfill \quad $\frac{3x^2+10xy}{4\,\sqrt[4]{(3x^2y+5xy^2)^{3}}}$
\item $A = [3, 2]$\quad \dotfill\; ???\;\dotfill \quad $\left( \frac{56}{4\,\sqrt[4]{114^3}} \;,\; \frac{87}{4\,\sqrt[4]{114^3}} \right)$
\item $g(x,y) = \sin\frac{x}{y} \enspace , \enspace D = [0, -1, ?]$\quad \dotfill\; ???\;\dotfill \quad $x+z=1$
\item $\Delta x = {-0.03} \enspace,\enspace \Delta y = {-0.2}$\quad \dotfill\; ???\;\dotfill \quad $0.03$
\item $h(x,y) = -2x^2-5y^2+3xy+4x$\quad \dotfill\; ???\;\dotfill \quad $\nicefrac{80}{49}$
\item $y = 4x+1$\quad \dotfill\; ???\;\dotfill \quad $\left[\frac{-33}{140}\;,\;\frac{2}{35}\right]$
\end{enumerate}
\end{varwidth}
\end{center}
\end{minipage}
\begin{minipage}{0.20\linewidth}
\begin{center}
{\Huge\bfseries 3.} \\[2mm]
\includegraphics[height=37mm]{images/braille.png}
{\small Braillove písmeno}
\end{center}
\end{minipage}
\end{center}
\end{minipage}
&
\begin{minipage}[c][104.5mm][t]{0.5\linewidth}
\begin{center}
\vspace{7mm}
{\huge \LARGE Funkcia dvoch premenných, skupina \textit{Gamma $\gamma$} -\romannumeral4}\\[5mm]
\begin{minipage}{0.95\linewidth}
\begin{center}
{\small \begin{multicols}{2} \raggedright
\textbf{(a)} Vypočítajte zadanú parciálnu deriváciu funkcie \(f(x,y)\). \\[1mm]
\textbf{(b)} Určte gradient \(f(x,y)\) v bode \(A\). \\[1mm]
\textbf{(c)} Napíšte rovnicu dotykovej roviny ku grafu funkcie \(g(x,y)\) v dotykovom bode \(D\). \par \columnbreak
\textbf{(d)} Pomocou lineárnej aproximácie odhadnite funkčnú hodnotu funkcie \(g(x,y)\) v bode, ktorý je posunutý o vektor \((\Delta x, \Delta y)\) od bodu \(D_0\). \\[1mm]
\textbf{(e)} Vypočítajte \(h(x,y)\) v jej lokálnom extréme. \\[1mm]
\textbf{(f)} Nájdite bod viazeného lokálneho extrému funkcie \(h(x,y)\) na priamke. \\[1mm]
\end{multicols}
}
\end{center}
\end{minipage}
\\[4mm]
\begin{minipage}{0.79\linewidth}
\begin{center}
\begin{varwidth}{\linewidth}
\begin{enumerate}
\normalsizerr
\item $f(x,y) = \sqrt{3x^2y + 5xy^2}\enspace , \enspace \pdv{f}{y}$\quad \dotfill\; ???\;\dotfill \quad $\frac{3x^2+10xy}{2\,\sqrt{3x^2y+5xy^2}}$
\item $A = [4, 3]$\quad \dotfill\; ???\;\dotfill \quad $\left( \frac{117}{2\,\sqrt{324}} \;,\; \frac{336}{2\,\sqrt{324}} \right)$
\item $g(x,y) = \sin\frac{x}{y} \enspace , \enspace D = [0, 1, ?]$\quad \dotfill\; ???\;\dotfill \quad $x-z=1$
\item $\Delta x = {-0.1} \enspace,\enspace \Delta y = {0.2}$\quad \dotfill\; ???\;\dotfill \quad $0.2$
\item $h(x,y) = 3x^2+3y^2-2xy-4x$\quad \dotfill\; ???\;\dotfill \quad $\nicefrac{-47}{32}$
\item $y = -2x-2$\quad \dotfill\; ???\;\dotfill \quad $\left[\frac{12}{19}\;,\;\frac{-14}{19}\right]$
\end{enumerate}
\end{varwidth}
\end{center}
\end{minipage}
\begin{minipage}{0.20\linewidth}
\begin{center}
{\Huge\bfseries 4.} \\[2mm]
\includegraphics[height=37mm]{images/braille.png}
{\small Braillove písmeno}
\end{center}
\end{minipage}
\end{center}
\end{minipage}
%
\end{tabular}
\newpage
\thispagestyle{empty}
\begin{tabular}{c:c}
\begin{minipage}[c][104.5mm][t]{0.5\linewidth}
\begin{center}
\vspace{7mm}
{\huge \LARGE Funkcia dvoch premenných, skupina \textit{Delta $\delta$} -\romannumeral1}\\[5mm]
\begin{minipage}{0.95\linewidth}
\begin{center}
{\small \begin{multicols}{2} \raggedright
\textbf{(a)} Vypočítajte zadanú parciálnu deriváciu funkcie \(f(x,y)\). \\[1mm]
\textbf{(b)} Určte gradient \(f(x,y)\) v bode \(A\). \\[1mm]
\textbf{(c)} Napíšte rovnicu dotykovej roviny ku grafu funkcie \(g(x,y)\) v dotykovom bode \(D\). \par \columnbreak
\textbf{(d)} Pomocou lineárnej aproximácie odhadnite funkčnú hodnotu funkcie \(g(x,y)\) v bode, ktorý je posunutý o vektor \((\Delta x, \Delta y)\) od bodu \(D_0\). \\[1mm]
\textbf{(e)} Vypočítajte \(h(x,y)\) v jej lokálnom extréme. \\[1mm]
\textbf{(f)} Nájdite bod viazeného lokálneho extrému funkcie \(h(x,y)\) na priamke. \\[1mm]
\end{multicols}
}
\end{center}
\end{minipage}
\\[4mm]
\begin{minipage}{0.79\linewidth}
\begin{center}
\begin{varwidth}{\linewidth}
\begin{enumerate}
\normalsizerr
\item $f(x,y) = \sqrt{3x^2y + 4xy^2}\enspace , \enspace \pdv{f}{y}$\quad \dotfill\; ???\;\dotfill \quad $\frac{3x^2+8xy}{2\,\sqrt{3x^2y+4xy^2}}$
\item $A = [3, 1]$\quad \dotfill\; ???\;\dotfill \quad $\left( \frac{22}{2\,\sqrt{39}} \;,\; \frac{102}{2\,\sqrt{39}} \right)$
\item $g(x,y) = \sin\frac{x}{y} \enspace , \enspace D = [0, -1, ?]$\quad \dotfill\; ???\;\dotfill \quad $x+z=1$
\item $\Delta x = {0.03} \enspace,\enspace \Delta y = {0.2}$\quad \dotfill\; ???\;\dotfill \quad $-0.2$
\item $h(x,y) = 6x^2+3y^2-7xy-3x$\quad \dotfill\; ???\;\dotfill \quad $\nicefrac{-27}{23}$
\item $y = -x-4$\quad \dotfill\; ???\;\dotfill \quad $\left[\frac{-1}{32}\;,\;\frac{-79}{32}\right]$
\end{enumerate}
\end{varwidth}
\end{center}
\end{minipage}
\begin{minipage}{0.20\linewidth}
\begin{center}
{\Huge\bfseries 1.} \\[2mm]
\includegraphics[height=37mm]{images/braille.png}
{\small Braillove písmeno}
\end{center}
\end{minipage}
\end{center}
\end{minipage}
&
\begin{minipage}[c][104.5mm][t]{0.5\linewidth}
\begin{center}
\vspace{7mm}
{\huge \LARGE Funkcia dvoch premenných, skupina \textit{Delta $\delta$} -\romannumeral2}\\[5mm]
\begin{minipage}{0.95\linewidth}
\begin{center}
{\small \begin{multicols}{2} \raggedright
\textbf{(a)} Vypočítajte zadanú parciálnu deriváciu funkcie \(f(x,y)\). \\[1mm]
\textbf{(b)} Určte gradient \(f(x,y)\) v bode \(A\). \\[1mm]
\textbf{(c)} Napíšte rovnicu dotykovej roviny ku grafu funkcie \(g(x,y)\) v dotykovom bode \(D\). \par \columnbreak
\textbf{(d)} Pomocou lineárnej aproximácie odhadnite funkčnú hodnotu funkcie \(g(x,y)\) v bode, ktorý je posunutý o vektor \((\Delta x, \Delta y)\) od bodu \(D_0\). \\[1mm]
\textbf{(e)} Vypočítajte \(h(x,y)\) v jej lokálnom extréme. \\[1mm]
\textbf{(f)} Nájdite bod viazeného lokálneho extrému funkcie \(h(x,y)\) na priamke. \\[1mm]
\end{multicols}
}
\end{center}
\end{minipage}
\\[4mm]
\begin{minipage}{0.79\linewidth}
\begin{center}
\begin{varwidth}{\linewidth}
\begin{enumerate}
\normalsizerr
\item $f(x,y) = \sqrt{7x^2y + 7xy^2}\enspace , \enspace \pdv{f}{x}$\quad \dotfill\; ???\;\dotfill \quad $\frac{14xy+7y^2}{2\,\sqrt{7x^2y+7xy^2}}$
\item $A = [4, 1]$\quad \dotfill\; ???\;\dotfill \quad $\left( \frac{63}{2\,\sqrt{140}} \;,\; \frac{168}{2\,\sqrt{140}} \right)$
\item $g(x,y) = \ln\frac{x}{y} \enspace , \enspace D=[e,1,?]$\quad \dotfill\; ???\;\dotfill \quad $\frac{x}{e} - y - z + 1 = 0$
\item $\Delta x = {-0.3} \enspace,\enspace \Delta y = {-0.03}$\quad \dotfill\; ???\;\dotfill \quad $0.86$
\item $h(x,y) = 3x^2+y^2+xy-6x$\quad \dotfill\; ???\;\dotfill \quad $\nicefrac{-36}{13}$
\item $y = 3x-5$\quad \dotfill\; ???\;\dotfill \quad $\left[\frac{41}{30}\;,\;\frac{-9}{10}\right]$
\end{enumerate}
\end{varwidth}
\end{center}
\end{minipage}
\begin{minipage}{0.20\linewidth}
\begin{center}
{\Huge\bfseries 2.} \\[2mm]
\includegraphics[height=37mm]{images/braille.png}
{\small Braillove písmeno}
\end{center}
\end{minipage}
\end{center}
\end{minipage}
\\ \hdashline
\begin{minipage}[c][104.5mm][t]{0.5\linewidth}
\begin{center}
\vspace{7mm}
{\huge \LARGE Funkcia dvoch premenných, skupina \textit{Delta $\delta$} -\romannumeral3}\\[5mm]
\begin{minipage}{0.95\linewidth}
\begin{center}
{\small \begin{multicols}{2} \raggedright
\textbf{(a)} Vypočítajte zadanú parciálnu deriváciu funkcie \(f(x,y)\). \\[1mm]
\textbf{(b)} Určte gradient \(f(x,y)\) v bode \(A\). \\[1mm]
\textbf{(c)} Napíšte rovnicu dotykovej roviny ku grafu funkcie \(g(x,y)\) v dotykovom bode \(D\). \par \columnbreak
\textbf{(d)} Pomocou lineárnej aproximácie odhadnite funkčnú hodnotu funkcie \(g(x,y)\) v bode, ktorý je posunutý o vektor \((\Delta x, \Delta y)\) od bodu \(D_0\). \\[1mm]
\textbf{(e)} Vypočítajte \(h(x,y)\) v jej lokálnom extréme. \\[1mm]
\textbf{(f)} Nájdite bod viazeného lokálneho extrému funkcie \(h(x,y)\) na priamke. \\[1mm]
\end{multicols}
}
\end{center}
\end{minipage}
\\[4mm]
\begin{minipage}{0.79\linewidth}
\begin{center}
\begin{varwidth}{\linewidth}
\begin{enumerate}
\normalsizerr
\item $f(x,y) = \sqrt[3]{6x^2y + 5xy^2}\enspace , \enspace \pdv{f}{x}$\quad \dotfill\; ???\;\dotfill \quad $\frac{12xy+5y^2}{3\,\sqrt[3]{(6x^2y+5xy^2)^{2}}}$
\item $A = [4, 2]$\quad \dotfill\; ???\;\dotfill \quad $\left( \frac{116}{3\,\sqrt[3]{272^2}} \;,\; \frac{352}{3\,\sqrt[3]{272^2}} \right)$
\item $g(x,y) = \sin\frac{x}{y} \enspace , \enspace D = [0, -1, ?]$\quad \dotfill\; ???\;\dotfill \quad $x+z=0$
\item $\Delta x = {0.3} \enspace,\enspace \Delta y = {0.02}$\quad \dotfill\; ???\;\dotfill \quad $-0.02$
\item $h(x,y) = 2x^2+4y^2+xy-5x$\quad \dotfill\; ???\;\dotfill \quad $\nicefrac{-100}{31}$
\item $y = -2x+6$\quad \dotfill\; ???\;\dotfill \quad $\left[\frac{-97}{32}\;,\;\frac{1}{16}\right]$
\end{enumerate}
\end{varwidth}
\end{center}
\end{minipage}
\begin{minipage}{0.20\linewidth}
\begin{center}
{\Huge\bfseries 3.} \\[2mm]
\includegraphics[height=37mm]{images/braille.png}
{\small Braillove písmeno}
\end{center}
\end{minipage}
\end{center}
\end{minipage}
&
\begin{minipage}[c][104.5mm][t]{0.5\linewidth}
\begin{center}
\vspace{7mm}
{\huge \LARGE Funkcia dvoch premenných, skupina \textit{Delta $\delta$} -\romannumeral4}\\[5mm]
\begin{minipage}{0.95\linewidth}
\begin{center}
{\small \begin{multicols}{2} \raggedright
\textbf{(a)} Vypočítajte zadanú parciálnu deriváciu funkcie \(f(x,y)\). \\[1mm]
\textbf{(b)} Určte gradient \(f(x,y)\) v bode \(A\). \\[1mm]
\textbf{(c)} Napíšte rovnicu dotykovej roviny ku grafu funkcie \(g(x,y)\) v dotykovom bode \(D\). \par \columnbreak
\textbf{(d)} Pomocou lineárnej aproximácie odhadnite funkčnú hodnotu funkcie \(g(x,y)\) v bode, ktorý je posunutý o vektor \((\Delta x, \Delta y)\) od bodu \(D_0\). \\[1mm]
\textbf{(e)} Vypočítajte \(h(x,y)\) v jej lokálnom extréme. \\[1mm]
\textbf{(f)} Nájdite bod viazeného lokálneho extrému funkcie \(h(x,y)\) na priamke. \\[1mm]
\end{multicols}
}
\end{center}
\end{minipage}
\\[4mm]
\begin{minipage}{0.79\linewidth}
\begin{center}
\begin{varwidth}{\linewidth}
\begin{enumerate}
\normalsizerr
\item $f(x,y) = \sqrt{7x^2y + 2xy^2}\enspace , \enspace \pdv{f}{y}$\quad \dotfill\; ???\;\dotfill \quad $\frac{7x^2+2xy}{2\,\sqrt{7x^2y+2xy^2}}$
\item $A = [4, 2]$\quad \dotfill\; ???\;\dotfill \quad $\left( \frac{120}{2\,\sqrt{256}} \;,\; \frac{144}{2\,\sqrt{256}} \right)$
\item $g(x,y) = \ln\frac{x}{y} \enspace , \enspace D=[e^{4},1,?]$\quad \dotfill\; ???\;\dotfill \quad $\frac{x}{e^4} - y - z + 4 = 0$
\item $\Delta x = {-0.3} \enspace,\enspace \Delta y = {0.3}$\quad \dotfill\; ???\;\dotfill \quad $4.295$
\item $h(x,y) = -6x^2-4y^2-3xy-5x$\quad \dotfill\; ???\;\dotfill \quad $\nicefrac{20}{21}$
\item $y = x-4$\quad \dotfill\; ???\;\dotfill \quad $\left[\frac{3}{2}\;,\;\frac{-5}{2}\right]$
\end{enumerate}
\end{varwidth}
\end{center}
\end{minipage}
\begin{minipage}{0.20\linewidth}
\begin{center}
{\Huge\bfseries 4.} \\[2mm]
\includegraphics[height=37mm]{images/braille.png}
{\small Braillove písmeno}
\end{center}
\end{minipage}
\end{center}
\end{minipage}
%
\end{tabular}
\newpage
\thispagestyle{empty}
\begin{tabular}{c:c}
\begin{minipage}[c][104.5mm][t]{0.5\linewidth}
\begin{center}
\vspace{7mm}
{\huge \LARGE Funkcia dvoch premenných, skupina \textit{Epsilon $\epsilon$} -\romannumeral1}\\[5mm]
\begin{minipage}{0.95\linewidth}
\begin{center}
{\small \begin{multicols}{2} \raggedright
\textbf{(a)} Vypočítajte zadanú parciálnu deriváciu funkcie \(f(x,y)\). \\[1mm]
\textbf{(b)} Určte gradient \(f(x,y)\) v bode \(A\). \\[1mm]
\textbf{(c)} Napíšte rovnicu dotykovej roviny ku grafu funkcie \(g(x,y)\) v dotykovom bode \(D\). \par \columnbreak
\textbf{(d)} Pomocou lineárnej aproximácie odhadnite funkčnú hodnotu funkcie \(g(x,y)\) v bode, ktorý je posunutý o vektor \((\Delta x, \Delta y)\) od bodu \(D_0\). \\[1mm]
\textbf{(e)} Vypočítajte \(h(x,y)\) v jej lokálnom extréme. \\[1mm]
\textbf{(f)} Nájdite bod viazeného lokálneho extrému funkcie \(h(x,y)\) na priamke. \\[1mm]
\end{multicols}
}
\end{center}
\end{minipage}
\\[4mm]
\begin{minipage}{0.79\linewidth}
\begin{center}
\begin{varwidth}{\linewidth}
\begin{enumerate}
\normalsizerr
\item $f(x,y) = \sqrt[3]{9x^2y + 9xy^2}\enspace , \enspace \pdv{f}{y}$\quad \dotfill\; ???\;\dotfill \quad $\frac{9x^2+18xy}{3\,\sqrt[3]{(9x^2y+9xy^2)^{2}}}$
\item $A = [1, 2]$\quad \dotfill\; ???\;\dotfill \quad $\left( \frac{72}{3\,\sqrt[3]{54^2}} \;,\; \frac{90}{3\,\sqrt[3]{54^2}} \right)$
\item $g(x,y) = \ln\frac{x}{y} \enspace , \enspace D=[1,1,?]$\quad \dotfill\; ???\;\dotfill \quad $x - y - z + 1 = 0$
\item $\Delta x = {0.02} \enspace,\enspace \Delta y = {0.02}$\quad \dotfill\; ???\;\dotfill \quad $0.04$
\item $h(x,y) = 3x^2+7y^2-2xy+3x$\quad \dotfill\; ???\;\dotfill \quad $\nicefrac{-63}{80}$
\item $y = -3x-1$\quad \dotfill\; ???\;\dotfill \quad $\left[\frac{-47}{144}\;,\;\frac{-1}{48}\right]$
\end{enumerate}
\end{varwidth}
\end{center}
\end{minipage}
\begin{minipage}{0.20\linewidth}
\begin{center}
{\Huge\bfseries 1.} \\[2mm]
\includegraphics[height=37mm]{images/braille.png}
{\small Braillove písmeno}
\end{center}
\end{minipage}
\end{center}
\end{minipage}
&
\begin{minipage}[c][104.5mm][t]{0.5\linewidth}
\begin{center}
\vspace{7mm}
{\huge \LARGE Funkcia dvoch premenných, skupina \textit{Epsilon $\epsilon$} -\romannumeral2}\\[5mm]
\begin{minipage}{0.95\linewidth}
\begin{center}
{\small \begin{multicols}{2} \raggedright
\textbf{(a)} Vypočítajte zadanú parciálnu deriváciu funkcie \(f(x,y)\). \\[1mm]
\textbf{(b)} Určte gradient \(f(x,y)\) v bode \(A\). \\[1mm]
\textbf{(c)} Napíšte rovnicu dotykovej roviny ku grafu funkcie \(g(x,y)\) v dotykovom bode \(D\). \par \columnbreak
\textbf{(d)} Pomocou lineárnej aproximácie odhadnite funkčnú hodnotu funkcie \(g(x,y)\) v bode, ktorý je posunutý o vektor \((\Delta x, \Delta y)\) od bodu \(D_0\). \\[1mm]
\textbf{(e)} Vypočítajte \(h(x,y)\) v jej lokálnom extréme. \\[1mm]
\textbf{(f)} Nájdite bod viazeného lokálneho extrému funkcie \(h(x,y)\) na priamke. \\[1mm]
\end{multicols}
}
\end{center}
\end{minipage}
\\[4mm]
\begin{minipage}{0.79\linewidth}
\begin{center}
\begin{varwidth}{\linewidth}
\begin{enumerate}
\normalsizerr
\item $f(x,y) = \sqrt{x^2y + xy^2}\enspace , \enspace \pdv{f}{y}$\quad \dotfill\; ???\;\dotfill \quad $\frac{x^2+2xy}{2\,\sqrt{x^2y+xy^2}}$
\item $A = [2, 4]$\quad \dotfill\; ???\;\dotfill \quad $\left( \frac{32}{2\,\sqrt{48}} \;,\; \frac{20}{2\,\sqrt{48}} \right)$
\item $g(x,y) = \sin\frac{x}{y} \enspace , \enspace D = [0, -2, ?]$\quad \dotfill\; ???\;\dotfill \quad $x+2z=0$
\item $\Delta x = {0.03} \enspace,\enspace \Delta y = {-0.01}$\quad \dotfill\; ???\;\dotfill \quad $0.005$
\item $h(x,y) = x^2+4y^2+xy+3x$\quad \dotfill\; ???\;\dotfill \quad $\nicefrac{-12}{5}$
\item $y = x+4$\quad \dotfill\; ???\;\dotfill \quad $\left[\frac{-13}{4}\;,\;\frac{9}{8}\right]$
\end{enumerate}
\end{varwidth}
\end{center}
\end{minipage}
\begin{minipage}{0.20\linewidth}
\begin{center}
{\Huge\bfseries 2.} \\[2mm]
\includegraphics[height=37mm]{images/braille.png}
{\small Braillove písmeno}
\end{center}
\end{minipage}
\end{center}
\end{minipage}
\\ \hdashline
\begin{minipage}[c][104.5mm][t]{0.5\linewidth}
\begin{center}
\vspace{7mm}
{\huge \LARGE Funkcia dvoch premenných, skupina \textit{Epsilon $\epsilon$} -\romannumeral3}\\[5mm]
\begin{minipage}{0.95\linewidth}
\begin{center}
{\small \begin{multicols}{2} \raggedright
\textbf{(a)} Vypočítajte zadanú parciálnu deriváciu funkcie \(f(x,y)\). \\[1mm]
\textbf{(b)} Určte gradient \(f(x,y)\) v bode \(A\). \\[1mm]
\textbf{(c)} Napíšte rovnicu dotykovej roviny ku grafu funkcie \(g(x,y)\) v dotykovom bode \(D\). \par \columnbreak
\textbf{(d)} Pomocou lineárnej aproximácie odhadnite funkčnú hodnotu funkcie \(g(x,y)\) v bode, ktorý je posunutý o vektor \((\Delta x, \Delta y)\) od bodu \(D_0\). \\[1mm]
\textbf{(e)} Vypočítajte \(h(x,y)\) v jej lokálnom extréme. \\[1mm]
\textbf{(f)} Nájdite bod viazeného lokálneho extrému funkcie \(h(x,y)\) na priamke. \\[1mm]
\end{multicols}
}
\end{center}
\end{minipage}
\\[4mm]
\begin{minipage}{0.79\linewidth}
\begin{center}
\begin{varwidth}{\linewidth}
\begin{enumerate}
\normalsizerr
\item $f(x,y) = \sqrt[4]{x^2y + 4xy^2}\enspace , \enspace \pdv{f}{y}$\quad \dotfill\; ???\;\dotfill \quad $\frac{x^2+8xy}{4\,\sqrt[4]{(x^2y+4xy^2)^{3}}}$
\item $A = [3, 3]$\quad \dotfill\; ???\;\dotfill \quad $\left( \frac{108}{4\,\sqrt[4]{135^3}} \;,\; \frac{81}{4\,\sqrt[4]{135^3}} \right)$
\item $g(x,y) = \ln\frac{x}{y} \enspace , \enspace D=[e,1,?]$\quad \dotfill\; ???\;\dotfill \quad $\frac{x}{e} + y + z + 1 = 0$
\item $\Delta x = {0.01} \enspace,\enspace \Delta y = {0.02}$\quad \dotfill\; ???\;\dotfill \quad $1.024$
\item $h(x,y) = 2x^2+4y^2+2xy+x$\quad \dotfill\; ???\;\dotfill \quad $\nicefrac{-3}{28}$
\item $y = -x-6$\quad \dotfill\; ???\;\dotfill \quad $\left[\frac{59}{8}\;,\;\frac{-11}{8}\right]$
\end{enumerate}
\end{varwidth}
\end{center}
\end{minipage}
\begin{minipage}{0.20\linewidth}
\begin{center}
{\Huge\bfseries 3.} \\[2mm]
\includegraphics[height=37mm]{images/braille.png}
{\small Braillove písmeno}
\end{center}
\end{minipage}
\end{center}
\end{minipage}
&
\begin{minipage}[c][104.5mm][t]{0.5\linewidth}
\begin{center}
\vspace{7mm}
{\huge \LARGE Funkcia dvoch premenných, skupina \textit{Epsilon $\epsilon$} -\romannumeral4}\\[5mm]
\begin{minipage}{0.95\linewidth}
\begin{center}
{\small \begin{multicols}{2} \raggedright
\textbf{(a)} Vypočítajte zadanú parciálnu deriváciu funkcie \(f(x,y)\). \\[1mm]
\textbf{(b)} Určte gradient \(f(x,y)\) v bode \(A\). \\[1mm]
\textbf{(c)} Napíšte rovnicu dotykovej roviny ku grafu funkcie \(g(x,y)\) v dotykovom bode \(D\). \par \columnbreak
\textbf{(d)} Pomocou lineárnej aproximácie odhadnite funkčnú hodnotu funkcie \(g(x,y)\) v bode, ktorý je posunutý o vektor \((\Delta x, \Delta y)\) od bodu \(D_0\). \\[1mm]
\textbf{(e)} Vypočítajte \(h(x,y)\) v jej lokálnom extréme. \\[1mm]
\textbf{(f)} Nájdite bod viazeného lokálneho extrému funkcie \(h(x,y)\) na priamke. \\[1mm]
\end{multicols}
}
\end{center}
\end{minipage}
\\[4mm]
\begin{minipage}{0.79\linewidth}
\begin{center}
\begin{varwidth}{\linewidth}
\begin{enumerate}
\normalsizerr
\item $f(x,y) = \sqrt[3]{3x^2y + 4xy^2}\enspace , \enspace \pdv{f}{x}$\quad \dotfill\; ???\;\dotfill \quad $\frac{6xy+4y^2}{3\,\sqrt[3]{(3x^2y+4xy^2)^{2}}}$
\item $A = [1, 4]$\quad \dotfill\; ???\;\dotfill \quad $\left( \frac{176}{3\,\sqrt[3]{76^2}} \;,\; \frac{35}{3\,\sqrt[3]{76^2}} \right)$
\item $g(x,y) = \ln\frac{x}{y} \enspace , \enspace D=[e,1,?]$\quad \dotfill\; ???\;\dotfill \quad $\frac{x}{e} - y - z + 1 = 0$
\item $\Delta x = {-0.1} \enspace,\enspace \Delta y = {0.03}$\quad \dotfill\; ???\;\dotfill \quad $1.063$
\item $h(x,y) = 2x^2+3y^2+3xy+x$\quad \dotfill\; ???\;\dotfill \quad $\nicefrac{-1}{11}$
\item $y = 3x+1$\quad \dotfill\; ???\;\dotfill \quad $\left[\frac{7}{38}\;,\;\frac{5}{38}\right]$
\end{enumerate}
\end{varwidth}
\end{center}
\end{minipage}
\begin{minipage}{0.20\linewidth}
\begin{center}
{\Huge\bfseries 4.} \\[2mm]
\includegraphics[height=37mm]{images/braille.png}
{\small Braillove písmeno}
\end{center}
\end{minipage}
\end{center}
\end{minipage}
%
\end{tabular}
\newpage
\thispagestyle{empty}
\begin{tabular}{c:c}
\begin{minipage}[c][104.5mm][t]{0.5\linewidth}
\begin{center}
\vspace{7mm}
{\huge \LARGE Funkcia dvoch premenných, skupina \textit{Zeta $\zeta$} -\romannumeral1}\\[5mm]
\begin{minipage}{0.95\linewidth}
\begin{center}
{\small \begin{multicols}{2} \raggedright
\textbf{(a)} Vypočítajte zadanú parciálnu deriváciu funkcie \(f(x,y)\). \\[1mm]
\textbf{(b)} Určte gradient \(f(x,y)\) v bode \(A\). \\[1mm]
\textbf{(c)} Napíšte rovnicu dotykovej roviny ku grafu funkcie \(g(x,y)\) v dotykovom bode \(D\). \par \columnbreak
\textbf{(d)} Pomocou lineárnej aproximácie odhadnite funkčnú hodnotu funkcie \(g(x,y)\) v bode, ktorý je posunutý o vektor \((\Delta x, \Delta y)\) od bodu \(D_0\). \\[1mm]
\textbf{(e)} Vypočítajte \(h(x,y)\) v jej lokálnom extréme. \\[1mm]
\textbf{(f)} Nájdite bod viazeného lokálneho extrému funkcie \(h(x,y)\) na priamke. \\[1mm]
\end{multicols}
}
\end{center}
\end{minipage}
\\[4mm]
\begin{minipage}{0.79\linewidth}
\begin{center}
\begin{varwidth}{\linewidth}
\begin{enumerate}
\normalsizerr
\item $f(x,y) = \sqrt[4]{3x^2y + 7xy^2}\enspace , \enspace \pdv{f}{x}$\quad \dotfill\; ???\;\dotfill \quad $\frac{6xy+7y^2}{4\,\sqrt[4]{(3x^2y+7xy^2)^{3}}}$
\item $A = [4, 1]$\quad \dotfill\; ???\;\dotfill \quad $\left( \frac{31}{4\,\sqrt[4]{76^3}} \;,\; \frac{208}{4\,\sqrt[4]{76^3}} \right)$
\item $g(x,y) = \sin\frac{x}{y} \enspace , \enspace D = [0, -3, ?]$\quad \dotfill\; ???\;\dotfill \quad $x+3z=0$
\item $\Delta x = {0.3} \enspace,\enspace \Delta y = {0.02}$\quad \dotfill\; ???\;\dotfill \quad $-0.9$
\item $h(x,y) = -3x^2-y^2-3xy-2x$\quad \dotfill\; ???\;\dotfill \quad $\nicefrac{4}{21}$
\item $y = 4x+1$\quad \dotfill\; ???\;\dotfill \quad $\left[\frac{-13}{62}\;,\;\frac{5}{31}\right]$
\end{enumerate}
\end{varwidth}
\end{center}
\end{minipage}
\begin{minipage}{0.20\linewidth}
\begin{center}
{\Huge\bfseries 1.} \\[2mm]
\includegraphics[height=37mm]{images/braille.png}
{\small Braillove písmeno}
\end{center}
\end{minipage}
\end{center}
\end{minipage}
&
\begin{minipage}[c][104.5mm][t]{0.5\linewidth}
\begin{center}
\vspace{7mm}
{\huge \LARGE Funkcia dvoch premenných, skupina \textit{Zeta $\zeta$} -\romannumeral2}\\[5mm]
\begin{minipage}{0.95\linewidth}
\begin{center}
{\small \begin{multicols}{2} \raggedright
\textbf{(a)} Vypočítajte zadanú parciálnu deriváciu funkcie \(f(x,y)\). \\[1mm]
\textbf{(b)} Určte gradient \(f(x,y)\) v bode \(A\). \\[1mm]
\textbf{(c)} Napíšte rovnicu dotykovej roviny ku grafu funkcie \(g(x,y)\) v dotykovom bode \(D\). \par \columnbreak
\textbf{(d)} Pomocou lineárnej aproximácie odhadnite funkčnú hodnotu funkcie \(g(x,y)\) v bode, ktorý je posunutý o vektor \((\Delta x, \Delta y)\) od bodu \(D_0\). \\[1mm]
\textbf{(e)} Vypočítajte \(h(x,y)\) v jej lokálnom extréme. \\[1mm]
\textbf{(f)} Nájdite bod viazeného lokálneho extrému funkcie \(h(x,y)\) na priamke. \\[1mm]
\end{multicols}
}
\end{center}
\end{minipage}
\\[4mm]
\begin{minipage}{0.79\linewidth}
\begin{center}
\begin{varwidth}{\linewidth}
\begin{enumerate}
\normalsizerr
\item $f(x,y) = \sqrt{6x^2y + 2xy^2}\enspace , \enspace \pdv{f}{y}$\quad \dotfill\; ???\;\dotfill \quad $\frac{6x^2+4xy}{2\,\sqrt{6x^2y+2xy^2}}$
\item $A = [1, 1]$\quad \dotfill\; ???\;\dotfill \quad $\left( \frac{14}{2\,\sqrt{8}} \;,\; \frac{10}{2\,\sqrt{8}} \right)$
\item $g(x,y) = \sin\frac{x}{y} \enspace , \enspace D = [0, -1, ?]$\quad \dotfill\; ???\;\dotfill \quad $x+z=0$
\item $\Delta x = {0.02} \enspace,\enspace \Delta y = {-0.01}$\quad \dotfill\; ???\;\dotfill \quad $0.01$
\item $h(x,y) = 2x^2+3y^2+3xy-3x$\quad \dotfill\; ???\;\dotfill \quad $\nicefrac{-9}{5}$
\item $y = -x+1$\quad \dotfill\; ???\;\dotfill \quad $\left[\frac{3}{2}\;,\;\frac{-1}{8}\right]$
\end{enumerate}
\end{varwidth}
\end{center}
\end{minipage}
\begin{minipage}{0.20\linewidth}
\begin{center}
{\Huge\bfseries 2.} \\[2mm]
\includegraphics[height=37mm]{images/braille.png}
{\small Braillove písmeno}
\end{center}
\end{minipage}
\end{center}
\end{minipage}
\\ \hdashline
\begin{minipage}[c][104.5mm][t]{0.5\linewidth}
\begin{center}
\vspace{7mm}
{\huge \LARGE Funkcia dvoch premenných, skupina \textit{Zeta $\zeta$} -\romannumeral3}\\[5mm]
\begin{minipage}{0.95\linewidth}
\begin{center}
{\small \begin{multicols}{2} \raggedright
\textbf{(a)} Vypočítajte zadanú parciálnu deriváciu funkcie \(f(x,y)\). \\[1mm]
\textbf{(b)} Určte gradient \(f(x,y)\) v bode \(A\). \\[1mm]
\textbf{(c)} Napíšte rovnicu dotykovej roviny ku grafu funkcie \(g(x,y)\) v dotykovom bode \(D\). \par \columnbreak
\textbf{(d)} Pomocou lineárnej aproximácie odhadnite funkčnú hodnotu funkcie \(g(x,y)\) v bode, ktorý je posunutý o vektor \((\Delta x, \Delta y)\) od bodu \(D_0\). \\[1mm]
\textbf{(e)} Vypočítajte \(h(x,y)\) v jej lokálnom extréme. \\[1mm]
\textbf{(f)} Nájdite bod viazeného lokálneho extrému funkcie \(h(x,y)\) na priamke. \\[1mm]
\end{multicols}
}
\end{center}
\end{minipage}
\\[4mm]
\begin{minipage}{0.79\linewidth}
\begin{center}
\begin{varwidth}{\linewidth}
\begin{enumerate}
\normalsizerr
\item $f(x,y) = \sqrt[4]{4x^2y + 3xy^2}\enspace , \enspace \pdv{f}{y}$\quad \dotfill\; ???\;\dotfill \quad $\frac{4x^2+6xy}{4\,\sqrt[4]{(4x^2y+3xy^2)^{3}}}$
\item $A = [1, 1]$\quad \dotfill\; ???\;\dotfill \quad $\left( \frac{22}{4\,\sqrt[4]{7^3}} \;,\; \frac{10}{4\,\sqrt[4]{7^3}} \right)$
\item $g(x,y) = \sin\frac{x}{y} \enspace , \enspace D = [0, -3, ?]$\quad \dotfill\; ???\;\dotfill \quad $x+3z=1$
\item $\Delta x = {0.02} \enspace,\enspace \Delta y = {0.1}$\quad \dotfill\; ???\;\dotfill \quad $-0.06$
\item $h(x,y) = 2x^2+y^2-xy+2x$\quad \dotfill\; ???\;\dotfill \quad $\nicefrac{-3}{7}$
\item $y = x+8$\quad \dotfill\; ???\;\dotfill \quad $\left[\frac{-5}{2}\;,\;\frac{11}{4}\right]$
\end{enumerate}
\end{varwidth}
\end{center}
\end{minipage}
\begin{minipage}{0.20\linewidth}
\begin{center}
{\Huge\bfseries 3.} \\[2mm]
\includegraphics[height=37mm]{images/braille.png}
{\small Braillove písmeno}
\end{center}
\end{minipage}
\end{center}
\end{minipage}
&
\begin{minipage}[c][104.5mm][t]{0.5\linewidth}
\begin{center}
\vspace{7mm}
{\huge \LARGE Funkcia dvoch premenných, skupina \textit{Zeta $\zeta$} -\romannumeral4}\\[5mm]
\begin{minipage}{0.95\linewidth}
\begin{center}
{\small \begin{multicols}{2} \raggedright
\textbf{(a)} Vypočítajte zadanú parciálnu deriváciu funkcie \(f(x,y)\). \\[1mm]
\textbf{(b)} Určte gradient \(f(x,y)\) v bode \(A\). \\[1mm]
\textbf{(c)} Napíšte rovnicu dotykovej roviny ku grafu funkcie \(g(x,y)\) v dotykovom bode \(D\). \par \columnbreak
\textbf{(d)} Pomocou lineárnej aproximácie odhadnite funkčnú hodnotu funkcie \(g(x,y)\) v bode, ktorý je posunutý o vektor \((\Delta x, \Delta y)\) od bodu \(D_0\). \\[1mm]
\textbf{(e)} Vypočítajte \(h(x,y)\) v jej lokálnom extréme. \\[1mm]
\textbf{(f)} Nájdite bod viazeného lokálneho extrému funkcie \(h(x,y)\) na priamke. \\[1mm]
\end{multicols}
}
\end{center}
\end{minipage}
\\[4mm]
\begin{minipage}{0.79\linewidth}
\begin{center}
\begin{varwidth}{\linewidth}
\begin{enumerate}
\normalsizerr
\item $f(x,y) = \sqrt[4]{3x^2y + 6xy^2}\enspace , \enspace \pdv{f}{x}$\quad \dotfill\; ???\;\dotfill \quad $\frac{6xy+6y^2}{4\,\sqrt[4]{(3x^2y+6xy^2)^{3}}}$
\item $A = [4, 1]$\quad \dotfill\; ???\;\dotfill \quad $\left( \frac{30}{4\,\sqrt[4]{72^3}} \;,\; \frac{192}{4\,\sqrt[4]{72^3}} \right)$
\item $g(x,y) = \ln\frac{x}{y} \enspace , \enspace D=[e,1,?]$\quad \dotfill\; ???\;\dotfill \quad $\frac{x}{e} - y - z + 1 = 0$
\item $\Delta x = {0.1} \enspace,\enspace \Delta y = {-0.03}$\quad \dotfill\; ???\;\dotfill \quad $1.007$
\item $h(x,y) = 4x^2+2y^2-2xy-4x$\quad \dotfill\; ???\;\dotfill \quad $\nicefrac{-8}{9}$
\item $y = -x-1$\quad \dotfill\; ???\;\dotfill \quad $\left[\frac{3}{8}\;,\;\frac{-7}{8}\right]$
\end{enumerate}
\end{varwidth}
\end{center}
\end{minipage}
\begin{minipage}{0.20\linewidth}
\begin{center}
{\Huge\bfseries 4.} \\[2mm]
\includegraphics[height=37mm]{images/braille.png}
{\small Braillove písmeno}
\end{center}
\end{minipage}
\end{center}
\end{minipage}
%
\end{tabular}
\newpage
\thispagestyle{empty}
\begin{tabular}{c:c}
\begin{minipage}[c][104.5mm][t]{0.5\linewidth}
\begin{center}
\vspace{7mm}
{\huge \LARGE Funkcia dvoch premenných, skupina \textit{Eta $\eta$} -\romannumeral1}\\[5mm]
\begin{minipage}{0.95\linewidth}
\begin{center}
{\small \begin{multicols}{2} \raggedright
\textbf{(a)} Vypočítajte zadanú parciálnu deriváciu funkcie \(f(x,y)\). \\[1mm]
\textbf{(b)} Určte gradient \(f(x,y)\) v bode \(A\). \\[1mm]
\textbf{(c)} Napíšte rovnicu dotykovej roviny ku grafu funkcie \(g(x,y)\) v dotykovom bode \(D\). \par \columnbreak
\textbf{(d)} Pomocou lineárnej aproximácie odhadnite funkčnú hodnotu funkcie \(g(x,y)\) v bode, ktorý je posunutý o vektor \((\Delta x, \Delta y)\) od bodu \(D_0\). \\[1mm]
\textbf{(e)} Vypočítajte \(h(x,y)\) v jej lokálnom extréme. \\[1mm]
\textbf{(f)} Nájdite bod viazeného lokálneho extrému funkcie \(h(x,y)\) na priamke. \\[1mm]
\end{multicols}
}
\end{center}
\end{minipage}
\\[4mm]
\begin{minipage}{0.79\linewidth}
\begin{center}
\begin{varwidth}{\linewidth}
\begin{enumerate}
\normalsizerr
\item $f(x,y) = \sqrt{3x^2y + 6xy^2}\enspace , \enspace \pdv{f}{y}$\quad \dotfill\; ???\;\dotfill \quad $\frac{3x^2+12xy}{2\,\sqrt{3x^2y+6xy^2}}$
\item $A = [4, 4]$\quad \dotfill\; ???\;\dotfill \quad $\left( \frac{384}{2\,\sqrt{576}} \;,\; \frac{240}{2\,\sqrt{576}} \right)$
\item $g(x,y) = \ln\frac{x}{y} \enspace , \enspace D=[e^{4},1,?]$\quad \dotfill\; ???\;\dotfill \quad $\frac{x}{e^4} - y - z + 4 = 0$
\item $\Delta x = {-0.3} \enspace,\enspace \Delta y = {-0.2}$\quad \dotfill\; ???\;\dotfill \quad $4.295$
\item $h(x,y) = 6x^2+4y^2+2xy-4x$\quad \dotfill\; ???\;\dotfill \quad $\nicefrac{-16}{23}$
\item $y = x+2$\quad \dotfill\; ???\;\dotfill \quad $\left[\frac{-2}{3}\;,\;\frac{4}{3}\right]$
\end{enumerate}
\end{varwidth}
\end{center}
\end{minipage}
\begin{minipage}{0.20\linewidth}
\begin{center}
{\Huge\bfseries 1.} \\[2mm]
\includegraphics[height=37mm]{images/braille.png}
{\small Braillove písmeno}
\end{center}
\end{minipage}
\end{center}
\end{minipage}
&
\begin{minipage}[c][104.5mm][t]{0.5\linewidth}
\begin{center}
\vspace{7mm}
{\huge \LARGE Funkcia dvoch premenných, skupina \textit{Eta $\eta$} -\romannumeral2}\\[5mm]
\begin{minipage}{0.95\linewidth}
\begin{center}
{\small \begin{multicols}{2} \raggedright
\textbf{(a)} Vypočítajte zadanú parciálnu deriváciu funkcie \(f(x,y)\). \\[1mm]
\textbf{(b)} Určte gradient \(f(x,y)\) v bode \(A\). \\[1mm]
\textbf{(c)} Napíšte rovnicu dotykovej roviny ku grafu funkcie \(g(x,y)\) v dotykovom bode \(D\). \par \columnbreak
\textbf{(d)} Pomocou lineárnej aproximácie odhadnite funkčnú hodnotu funkcie \(g(x,y)\) v bode, ktorý je posunutý o vektor \((\Delta x, \Delta y)\) od bodu \(D_0\). \\[1mm]
\textbf{(e)} Vypočítajte \(h(x,y)\) v jej lokálnom extréme. \\[1mm]
\textbf{(f)} Nájdite bod viazeného lokálneho extrému funkcie \(h(x,y)\) na priamke. \\[1mm]
\end{multicols}
}
\end{center}
\end{minipage}
\\[4mm]
\begin{minipage}{0.79\linewidth}
\begin{center}
\begin{varwidth}{\linewidth}
\begin{enumerate}
\normalsizerr
\item $f(x,y) = \sqrt{5x^2y + 5xy^2}\enspace , \enspace \pdv{f}{y}$\quad \dotfill\; ???\;\dotfill \quad $\frac{5x^2+10xy}{2\,\sqrt{5x^2y+5xy^2}}$
\item $A = [3, 2]$\quad \dotfill\; ???\;\dotfill \quad $\left( \frac{160}{2\,\sqrt{150}} \;,\; \frac{105}{2\,\sqrt{150}} \right)$
\item $g(x,y) = \sin\frac{x}{y} \enspace , \enspace D = [0, -3, ?]$\quad \dotfill\; ???\;\dotfill \quad $x+3z=0$
\item $\Delta x = {-0.03} \enspace,\enspace \Delta y = {0.1}$\quad \dotfill\; ???\;\dotfill \quad $0.09$
\item $h(x,y) = 3x^2+2y^2+xy-3x$\quad \dotfill\; ???\;\dotfill \quad $\nicefrac{-18}{23}$
\item $y = -2x-3$\quad \dotfill\; ???\;\dotfill \quad $\left[-1\;,\;\frac{-9}{13}\right]$
\end{enumerate}
\end{varwidth}
\end{center}
\end{minipage}
\begin{minipage}{0.20\linewidth}
\begin{center}
{\Huge\bfseries 2.} \\[2mm]
\includegraphics[height=37mm]{images/braille.png}
{\small Braillove písmeno}
\end{center}
\end{minipage}
\end{center}
\end{minipage}
\\ \hdashline
\begin{minipage}[c][104.5mm][t]{0.5\linewidth}
\begin{center}
\vspace{7mm}
{\huge \LARGE Funkcia dvoch premenných, skupina \textit{Eta $\eta$} -\romannumeral3}\\[5mm]
\begin{minipage}{0.95\linewidth}
\begin{center}
{\small \begin{multicols}{2} \raggedright
\textbf{(a)} Vypočítajte zadanú parciálnu deriváciu funkcie \(f(x,y)\). \\[1mm]
\textbf{(b)} Určte gradient \(f(x,y)\) v bode \(A\). \\[1mm]
\textbf{(c)} Napíšte rovnicu dotykovej roviny ku grafu funkcie \(g(x,y)\) v dotykovom bode \(D\). \par \columnbreak
\textbf{(d)} Pomocou lineárnej aproximácie odhadnite funkčnú hodnotu funkcie \(g(x,y)\) v bode, ktorý je posunutý o vektor \((\Delta x, \Delta y)\) od bodu \(D_0\). \\[1mm]
\textbf{(e)} Vypočítajte \(h(x,y)\) v jej lokálnom extréme. \\[1mm]
\textbf{(f)} Nájdite bod viazeného lokálneho extrému funkcie \(h(x,y)\) na priamke. \\[1mm]
\end{multicols}
}
\end{center}
\end{minipage}
\\[4mm]
\begin{minipage}{0.79\linewidth}
\begin{center}
\begin{varwidth}{\linewidth}
\begin{enumerate}
\normalsizerr
\item $f(x,y) = \sqrt[3]{7x^2y + 3xy^2}\enspace , \enspace \pdv{f}{x}$\quad \dotfill\; ???\;\dotfill \quad $\frac{14xy+3y^2}{3\,\sqrt[3]{(7x^2y+3xy^2)^{2}}}$
\item $A = [4, 4]$\quad \dotfill\; ???\;\dotfill \quad $\left( \frac{272}{3\,\sqrt[3]{640^2}} \;,\; \frac{208}{3\,\sqrt[3]{640^2}} \right)$
\item $g(x,y) = \sin\frac{x}{y} \enspace , \enspace D = [0, -3, ?]$\quad \dotfill\; ???\;\dotfill \quad $x-3z=0$
\item $\Delta x = {-0.2} \enspace,\enspace \Delta y = {-0.01}$\quad \dotfill\; ???\;\dotfill \quad $0.003$
\item $h(x,y) = -6x^2-5y^2-3xy+x$\quad \dotfill\; ???\;\dotfill \quad $\nicefrac{5}{111}$
\item $y = 2x+1$\quad \dotfill\; ???\;\dotfill \quad $\left[\frac{9}{32}\;,\;\frac{5}{16}\right]$
\end{enumerate}
\end{varwidth}
\end{center}
\end{minipage}
\begin{minipage}{0.20\linewidth}
\begin{center}
{\Huge\bfseries 3.} \\[2mm]
\includegraphics[height=37mm]{images/braille.png}
{\small Braillove písmeno}
\end{center}
\end{minipage}
\end{center}
\end{minipage}
&
\begin{minipage}[c][104.5mm][t]{0.5\linewidth}
\begin{center}
\vspace{7mm}
{\huge \LARGE Funkcia dvoch premenných, skupina \textit{Eta $\eta$} -\romannumeral4}\\[5mm]
\begin{minipage}{0.95\linewidth}
\begin{center}
{\small \begin{multicols}{2} \raggedright
\textbf{(a)} Vypočítajte zadanú parciálnu deriváciu funkcie \(f(x,y)\). \\[1mm]
\textbf{(b)} Určte gradient \(f(x,y)\) v bode \(A\). \\[1mm]
\textbf{(c)} Napíšte rovnicu dotykovej roviny ku grafu funkcie \(g(x,y)\) v dotykovom bode \(D\). \par \columnbreak
\textbf{(d)} Pomocou lineárnej aproximácie odhadnite funkčnú hodnotu funkcie \(g(x,y)\) v bode, ktorý je posunutý o vektor \((\Delta x, \Delta y)\) od bodu \(D_0\). \\[1mm]
\textbf{(e)} Vypočítajte \(h(x,y)\) v jej lokálnom extréme. \\[1mm]
\textbf{(f)} Nájdite bod viazeného lokálneho extrému funkcie \(h(x,y)\) na priamke. \\[1mm]
\end{multicols}
}
\end{center}
\end{minipage}
\\[4mm]
\begin{minipage}{0.79\linewidth}
\begin{center}
\begin{varwidth}{\linewidth}
\begin{enumerate}
\normalsizerr
\item $f(x,y) = \sqrt[3]{6x^2y + xy^2}\enspace , \enspace \pdv{f}{x}$\quad \dotfill\; ???\;\dotfill \quad $\frac{12xy+y^2}{3\,\sqrt[3]{(6x^2y+xy^2)^{2}}}$
\item $A = [1, 1]$\quad \dotfill\; ???\;\dotfill \quad $\left( \frac{13}{3\,\sqrt[3]{7^2}} \;,\; \frac{16}{3\,\sqrt[3]{7^2}} \right)$
\item $g(x,y) = \sin\frac{x}{y} \enspace , \enspace D = [0, 1, ?]$\quad \dotfill\; ???\;\dotfill \quad $x-z=1$
\item $\Delta x = {-0.03} \enspace,\enspace \Delta y = {0.02}$\quad \dotfill\; ???\;\dotfill \quad $0.02$
\item $h(x,y) = -6x^2-2y^2+2xy+2x$\quad \dotfill\; ???\;\dotfill \quad $\nicefrac{9}{44}$
\item $y = -2x+7$\quad \dotfill\; ???\;\dotfill \quad $\left[2\;,\;\frac{27}{5}\right]$
\end{enumerate}
\end{varwidth}
\end{center}
\end{minipage}
\begin{minipage}{0.20\linewidth}
\begin{center}
{\Huge\bfseries 4.} \\[2mm]
\includegraphics[height=37mm]{images/braille.png}
{\small Braillove písmeno}
\end{center}
\end{minipage}
\end{center}
\end{minipage}
%
\end{tabular}
\newpage
\thispagestyle{empty}
\begin{tabular}{c:c}
\begin{minipage}[c][104.5mm][t]{0.5\linewidth}
\begin{center}
\vspace{7mm}
{\huge \LARGE Funkcia dvoch premenných, skupina \textit{Theta $\theta$} -\romannumeral1}\\[5mm]
\begin{minipage}{0.95\linewidth}
\begin{center}
{\small \begin{multicols}{2} \raggedright
\textbf{(a)} Vypočítajte zadanú parciálnu deriváciu funkcie \(f(x,y)\). \\[1mm]
\textbf{(b)} Určte gradient \(f(x,y)\) v bode \(A\). \\[1mm]
\textbf{(c)} Napíšte rovnicu dotykovej roviny ku grafu funkcie \(g(x,y)\) v dotykovom bode \(D\). \par \columnbreak
\textbf{(d)} Pomocou lineárnej aproximácie odhadnite funkčnú hodnotu funkcie \(g(x,y)\) v bode, ktorý je posunutý o vektor \((\Delta x, \Delta y)\) od bodu \(D_0\). \\[1mm]
\textbf{(e)} Vypočítajte \(h(x,y)\) v jej lokálnom extréme. \\[1mm]
\textbf{(f)} Nájdite bod viazeného lokálneho extrému funkcie \(h(x,y)\) na priamke. \\[1mm]
\end{multicols}
}
\end{center}
\end{minipage}
\\[4mm]
\begin{minipage}{0.79\linewidth}
\begin{center}
\begin{varwidth}{\linewidth}
\begin{enumerate}
\normalsizerr
\item $f(x,y) = \sqrt[4]{2x^2y + 6xy^2}\enspace , \enspace \pdv{f}{x}$\quad \dotfill\; ???\;\dotfill \quad $\frac{4xy+6y^2}{4\,\sqrt[4]{(2x^2y+6xy^2)^{3}}}$
\item $A = [4, 2]$\quad \dotfill\; ???\;\dotfill \quad $\left( \frac{56}{4\,\sqrt[4]{160^3}} \;,\; \frac{128}{4\,\sqrt[4]{160^3}} \right)$
\item $g(x,y) = \ln\frac{x}{y} \enspace , \enspace D=[e,1,?]$\quad \dotfill\; ???\;\dotfill \quad $\frac{x}{e} - y - z + 1 = 0$
\item $\Delta x = {-0.1} \enspace,\enspace \Delta y = {0.2}$\quad \dotfill\; ???\;\dotfill \quad $0.763$
\item $h(x,y) = -3x^2-6y^2+3xy-x$\quad \dotfill\; ???\;\dotfill \quad $\nicefrac{2}{27}$
\item $y = x+4$\quad \dotfill\; ???\;\dotfill \quad $\left[\frac{59}{12}\;,\;\frac{11}{12}\right]$
\end{enumerate}
\end{varwidth}
\end{center}
\end{minipage}
\begin{minipage}{0.20\linewidth}
\begin{center}
{\Huge\bfseries 1.} \\[2mm]
\includegraphics[height=37mm]{images/braille.png}
{\small Braillove písmeno}
\end{center}
\end{minipage}
\end{center}
\end{minipage}
&
\begin{minipage}[c][104.5mm][t]{0.5\linewidth}
\begin{center}
\vspace{7mm}
{\huge \LARGE Funkcia dvoch premenných, skupina \textit{Theta $\theta$} -\romannumeral2}\\[5mm]
\begin{minipage}{0.95\linewidth}
\begin{center}
{\small \begin{multicols}{2} \raggedright
\textbf{(a)} Vypočítajte zadanú parciálnu deriváciu funkcie \(f(x,y)\). \\[1mm]
\textbf{(b)} Určte gradient \(f(x,y)\) v bode \(A\). \\[1mm]
\textbf{(c)} Napíšte rovnicu dotykovej roviny ku grafu funkcie \(g(x,y)\) v dotykovom bode \(D\). \par \columnbreak
\textbf{(d)} Pomocou lineárnej aproximácie odhadnite funkčnú hodnotu funkcie \(g(x,y)\) v bode, ktorý je posunutý o vektor \((\Delta x, \Delta y)\) od bodu \(D_0\). \\[1mm]
\textbf{(e)} Vypočítajte \(h(x,y)\) v jej lokálnom extréme. \\[1mm]
\textbf{(f)} Nájdite bod viazeného lokálneho extrému funkcie \(h(x,y)\) na priamke. \\[1mm]
\end{multicols}
}
\end{center}
\end{minipage}
\\[4mm]
\begin{minipage}{0.79\linewidth}
\begin{center}
\begin{varwidth}{\linewidth}
\begin{enumerate}
\normalsizerr
\item $f(x,y) = \sqrt[4]{4x^2y + 6xy^2}\enspace , \enspace \pdv{f}{y}$\quad \dotfill\; ???\;\dotfill \quad $\frac{4x^2+12xy}{4\,\sqrt[4]{(4x^2y+6xy^2)^{3}}}$
\item $A = [1, 2]$\quad \dotfill\; ???\;\dotfill \quad $\left( \frac{40}{4\,\sqrt[4]{32^3}} \;,\; \frac{28}{4\,\sqrt[4]{32^3}} \right)$
\item $g(x,y) = \ln\frac{x}{y} \enspace , \enspace D=[e^{3},1,?]$\quad \dotfill\; ???\;\dotfill \quad $\frac{x}{e^3} - y - z + 3 = 0$
\item $\Delta x = {-0.1} \enspace,\enspace \Delta y = {0.03}$\quad \dotfill\; ???\;\dotfill \quad $3.095$
\item $h(x,y) = -5x^2-y^2+3xy+2x$\quad \dotfill\; ???\;\dotfill \quad $\nicefrac{5}{11}$
\item $y = 2x-9$\quad \dotfill\; ???\;\dotfill \quad $\left[\frac{11}{6}\;,\;\frac{-16}{15}\right]$
\end{enumerate}
\end{varwidth}
\end{center}
\end{minipage}
\begin{minipage}{0.20\linewidth}
\begin{center}
{\Huge\bfseries 2.} \\[2mm]
\includegraphics[height=37mm]{images/braille.png}
{\small Braillove písmeno}
\end{center}
\end{minipage}
\end{center}
\end{minipage}
\\ \hdashline
\begin{minipage}[c][104.5mm][t]{0.5\linewidth}
\begin{center}
\vspace{7mm}
{\huge \LARGE Funkcia dvoch premenných, skupina \textit{Theta $\theta$} -\romannumeral3}\\[5mm]
\begin{minipage}{0.95\linewidth}
\begin{center}
{\small \begin{multicols}{2} \raggedright
\textbf{(a)} Vypočítajte zadanú parciálnu deriváciu funkcie \(f(x,y)\). \\[1mm]
\textbf{(b)} Určte gradient \(f(x,y)\) v bode \(A\). \\[1mm]
\textbf{(c)} Napíšte rovnicu dotykovej roviny ku grafu funkcie \(g(x,y)\) v dotykovom bode \(D\). \par \columnbreak
\textbf{(d)} Pomocou lineárnej aproximácie odhadnite funkčnú hodnotu funkcie \(g(x,y)\) v bode, ktorý je posunutý o vektor \((\Delta x, \Delta y)\) od bodu \(D_0\). \\[1mm]
\textbf{(e)} Vypočítajte \(h(x,y)\) v jej lokálnom extréme. \\[1mm]
\textbf{(f)} Nájdite bod viazeného lokálneho extrému funkcie \(h(x,y)\) na priamke. \\[1mm]
\end{multicols}
}
\end{center}
\end{minipage}
\\[4mm]
\begin{minipage}{0.79\linewidth}
\begin{center}
\begin{varwidth}{\linewidth}
\begin{enumerate}
\normalsizerr
\item $f(x,y) = \sqrt{3x^2y + 5xy^2}\enspace , \enspace \pdv{f}{y}$\quad \dotfill\; ???\;\dotfill \quad $\frac{3x^2+10xy}{2\,\sqrt{3x^2y+5xy^2}}$
\item $A = [3, 1]$\quad \dotfill\; ???\;\dotfill \quad $\left( \frac{46}{2\,\sqrt{42}} \;,\; \frac{57}{2\,\sqrt{42}} \right)$
\item $g(x,y) = \sin\frac{x}{y} \enspace , \enspace D = [0, -2, ?]$\quad \dotfill\; ???\;\dotfill \quad $x+2z=1$
\item $\Delta x = {0.02} \enspace,\enspace \Delta y = {0.03}$\quad \dotfill\; ???\;\dotfill \quad $-0.04$
\item $h(x,y) = x^2+5y^2+2xy+2x$\quad \dotfill\; ???\;\dotfill \quad $\nicefrac{-5}{4}$
\item $y = -x+8$\quad \dotfill\; ???\;\dotfill \quad $\left[\frac{31}{4}\;,\;\frac{1}{8}\right]$
\end{enumerate}
\end{varwidth}
\end{center}
\end{minipage}
\begin{minipage}{0.20\linewidth}
\begin{center}
{\Huge\bfseries 3.} \\[2mm]
\includegraphics[height=37mm]{images/braille.png}
{\small Braillove písmeno}
\end{center}
\end{minipage}
\end{center}
\end{minipage}
&
\begin{minipage}[c][104.5mm][t]{0.5\linewidth}
\begin{center}
\vspace{7mm}
{\huge \LARGE Funkcia dvoch premenných, skupina \textit{Theta $\theta$} -\romannumeral4}\\[5mm]
\begin{minipage}{0.95\linewidth}
\begin{center}
{\small \begin{multicols}{2} \raggedright
\textbf{(a)} Vypočítajte zadanú parciálnu deriváciu funkcie \(f(x,y)\). \\[1mm]
\textbf{(b)} Určte gradient \(f(x,y)\) v bode \(A\). \\[1mm]
\textbf{(c)} Napíšte rovnicu dotykovej roviny ku grafu funkcie \(g(x,y)\) v dotykovom bode \(D\). \par \columnbreak
\textbf{(d)} Pomocou lineárnej aproximácie odhadnite funkčnú hodnotu funkcie \(g(x,y)\) v bode, ktorý je posunutý o vektor \((\Delta x, \Delta y)\) od bodu \(D_0\). \\[1mm]
\textbf{(e)} Vypočítajte \(h(x,y)\) v jej lokálnom extréme. \\[1mm]
\textbf{(f)} Nájdite bod viazeného lokálneho extrému funkcie \(h(x,y)\) na priamke. \\[1mm]
\end{multicols}
}
\end{center}
\end{minipage}
\\[4mm]
\begin{minipage}{0.79\linewidth}
\begin{center}
\begin{varwidth}{\linewidth}
\begin{enumerate}
\normalsizerr
\item $f(x,y) = \sqrt[4]{x^2y + 3xy^2}\enspace , \enspace \pdv{f}{x}$\quad \dotfill\; ???\;\dotfill \quad $\frac{2xy+3y^2}{4\,\sqrt[4]{(x^2y+3xy^2)^{3}}}$
\item $A = [4, 4]$\quad \dotfill\; ???\;\dotfill \quad $\left( \frac{80}{4\,\sqrt[4]{256^3}} \;,\; \frac{224}{4\,\sqrt[4]{256^3}} \right)$
\item $g(x,y) = \sin\frac{x}{y} \enspace , \enspace D = [0, -1, ?]$\quad \dotfill\; ???\;\dotfill \quad $x+z=0$
\item $\Delta x = {0.03} \enspace,\enspace \Delta y = {0.3}$\quad \dotfill\; ???\;\dotfill \quad $-0.3$
\item $h(x,y) = 5x^2+3y^2-3xy+2x$\quad \dotfill\; ???\;\dotfill \quad $\nicefrac{-11}{51}$
\item $y = -2x+3$\quad \dotfill\; ???\;\dotfill \quad $\left[\frac{43}{46}\;,\;\frac{26}{11}\right]$
\end{enumerate}
\end{varwidth}
\end{center}
\end{minipage}
\begin{minipage}{0.20\linewidth}
\begin{center}
{\Huge\bfseries 4.} \\[2mm]
\includegraphics[height=37mm]{images/braille.png}
{\small Braillove písmeno}
\end{center}
\end{minipage}
\end{center}
\end{minipage}
%
\end{tabular}
\newpage
\thispagestyle{empty}
\begin{tabular}{c:c}
\begin{minipage}[c][104.5mm][t]{0.5\linewidth}
\begin{center}
\vspace{7mm}
{\huge \LARGE Funkcia dvoch premenných, skupina \textit{Iota $\iota$} -\romannumeral1}\\[5mm]
\begin{minipage}{0.95\linewidth}
\begin{center}
{\small \begin{multicols}{2} \raggedright
\textbf{(a)} Vypočítajte zadanú parciálnu deriváciu funkcie \(f(x,y)\). \\[1mm]
\textbf{(b)} Určte gradient \(f(x,y)\) v bode \(A\). \\[1mm]
\textbf{(c)} Napíšte rovnicu dotykovej roviny ku grafu funkcie \(g(x,y)\) v dotykovom bode \(D\). \par \columnbreak
\textbf{(d)} Pomocou lineárnej aproximácie odhadnite funkčnú hodnotu funkcie \(g(x,y)\) v bode, ktorý je posunutý o vektor \((\Delta x, \Delta y)\) od bodu \(D_0\). \\[1mm]
\textbf{(e)} Vypočítajte \(h(x,y)\) v jej lokálnom extréme. \\[1mm]
\textbf{(f)} Nájdite bod viazeného lokálneho extrému funkcie \(h(x,y)\) na priamke. \\[1mm]
\end{multicols}
}
\end{center}
\end{minipage}
\\[4mm]
\begin{minipage}{0.79\linewidth}
\begin{center}
\begin{varwidth}{\linewidth}
\begin{enumerate}
\normalsizerr
\item $f(x,y) = \sqrt{3x^2y + 2xy^2}\enspace , \enspace \pdv{f}{y}$\quad \dotfill\; ???\;\dotfill \quad $\frac{3x^2+4xy}{2\,\sqrt{3x^2y+2xy^2}}$
\item $A = [1, 2]$\quad \dotfill\; ???\;\dotfill \quad $\left( \frac{40}{2\,\sqrt{14}} \;,\; \frac{11}{2\,\sqrt{14}} \right)$
\item $g(x,y) = \ln\frac{x}{y} \enspace , \enspace D=[1,1,?]$\quad \dotfill\; ???\;\dotfill \quad $x - y - z + 1 = 0$
\item $\Delta x = {-0.2} \enspace,\enspace \Delta y = {0.1}$\quad \dotfill\; ???\;\dotfill \quad $-0.1$
\item $h(x,y) = -2x^2-3y^2-2xy+4x$\quad \dotfill\; ???\;\dotfill \quad $\nicefrac{12}{5}$
\item $y = 2x+4$\quad \dotfill\; ???\;\dotfill \quad $\left[\frac{11}{9}\;,\;\frac{10}{9}\right]$
\end{enumerate}
\end{varwidth}
\end{center}
\end{minipage}
\begin{minipage}{0.20\linewidth}
\begin{center}
{\Huge\bfseries 1.} \\[2mm]
\includegraphics[height=37mm]{images/braille.png}
{\small Braillove písmeno}
\end{center}
\end{minipage}
\end{center}
\end{minipage}
&
\begin{minipage}[c][104.5mm][t]{0.5\linewidth}
\begin{center}
\vspace{7mm}
{\huge \LARGE Funkcia dvoch premenných, skupina \textit{Iota $\iota$} -\romannumeral2}\\[5mm]
\begin{minipage}{0.95\linewidth}
\begin{center}
{\small \begin{multicols}{2} \raggedright
\textbf{(a)} Vypočítajte zadanú parciálnu deriváciu funkcie \(f(x,y)\). \\[1mm]
\textbf{(b)} Určte gradient \(f(x,y)\) v bode \(A\). \\[1mm]
\textbf{(c)} Napíšte rovnicu dotykovej roviny ku grafu funkcie \(g(x,y)\) v dotykovom bode \(D\). \par \columnbreak
\textbf{(d)} Pomocou lineárnej aproximácie odhadnite funkčnú hodnotu funkcie \(g(x,y)\) v bode, ktorý je posunutý o vektor \((\Delta x, \Delta y)\) od bodu \(D_0\). \\[1mm]
\textbf{(e)} Vypočítajte \(h(x,y)\) v jej lokálnom extréme. \\[1mm]
\textbf{(f)} Nájdite bod viazeného lokálneho extrému funkcie \(h(x,y)\) na priamke. \\[1mm]
\end{multicols}
}
\end{center}
\end{minipage}
\\[4mm]
\begin{minipage}{0.79\linewidth}
\begin{center}
\begin{varwidth}{\linewidth}
\begin{enumerate}
\normalsizerr
\item $f(x,y) = \sqrt{6x^2y + 5xy^2}\enspace , \enspace \pdv{f}{x}$\quad \dotfill\; ???\;\dotfill \quad $\frac{12xy+5y^2}{2\,\sqrt{6x^2y+5xy^2}}$
\item $A = [1, 3]$\quad \dotfill\; ???\;\dotfill \quad $\left( \frac{81}{2\,\sqrt{63}} \;,\; \frac{36}{2\,\sqrt{63}} \right)$
\item $g(x,y) = \ln\frac{x}{y} \enspace , \enspace D=[e^{4},1,?]$\quad \dotfill\; ???\;\dotfill \quad $\frac{x}{e^4} - y - z + 4 = 0$
\item $\Delta x = {-0.2} \enspace,\enspace \Delta y = {0.2}$\quad \dotfill\; ???\;\dotfill \quad $4.196$
\item $h(x,y) = -4x^2-4y^2-2xy+3x$\quad \dotfill\; ???\;\dotfill \quad $\nicefrac{3}{5}$
\item $y = x+6$\quad \dotfill\; ???\;\dotfill \quad $\left[\frac{-57}{20}\;,\;\frac{21}{4}\right]$
\end{enumerate}
\end{varwidth}
\end{center}
\end{minipage}
\begin{minipage}{0.20\linewidth}
\begin{center}
{\Huge\bfseries 2.} \\[2mm]
\includegraphics[height=37mm]{images/braille.png}
{\small Braillove písmeno}
\end{center}
\end{minipage}
\end{center}
\end{minipage}
\\ \hdashline
\begin{minipage}[c][104.5mm][t]{0.5\linewidth}
\begin{center}
\vspace{7mm}
{\huge \LARGE Funkcia dvoch premenných, skupina \textit{Iota $\iota$} -\romannumeral3}\\[5mm]
\begin{minipage}{0.95\linewidth}
\begin{center}
{\small \begin{multicols}{2} \raggedright
\textbf{(a)} Vypočítajte zadanú parciálnu deriváciu funkcie \(f(x,y)\). \\[1mm]
\textbf{(b)} Určte gradient \(f(x,y)\) v bode \(A\). \\[1mm]
\textbf{(c)} Napíšte rovnicu dotykovej roviny ku grafu funkcie \(g(x,y)\) v dotykovom bode \(D\). \par \columnbreak
\textbf{(d)} Pomocou lineárnej aproximácie odhadnite funkčnú hodnotu funkcie \(g(x,y)\) v bode, ktorý je posunutý o vektor \((\Delta x, \Delta y)\) od bodu \(D_0\). \\[1mm]
\textbf{(e)} Vypočítajte \(h(x,y)\) v jej lokálnom extréme. \\[1mm]
\textbf{(f)} Nájdite bod viazeného lokálneho extrému funkcie \(h(x,y)\) na priamke. \\[1mm]
\end{multicols}
}
\end{center}
\end{minipage}
\\[4mm]
\begin{minipage}{0.79\linewidth}
\begin{center}
\begin{varwidth}{\linewidth}
\begin{enumerate}
\normalsizerr
\item $f(x,y) = \sqrt[3]{2x^2y + 5xy^2}\enspace , \enspace \pdv{f}{x}$\quad \dotfill\; ???\;\dotfill \quad $\frac{4xy-5y^2}{3\,\sqrt[3]{(2x^2y+5xy^2)^{2}}}$
\item $A = [1, 4]$\quad \dotfill\; ???\;\dotfill \quad $\left( \frac{96}{3\,\sqrt[3]{88^2}} \;,\; \frac{42}{3\,\sqrt[3]{88^2}} \right)$
\item $g(x,y) = \sin\frac{x}{y} \enspace , \enspace D = [0, -1, ?]$\quad \dotfill\; ???\;\dotfill \quad $x-z=0$
\item $\Delta x = {-0.03} \enspace,\enspace \Delta y = {0.03}$\quad \dotfill\; ???\;\dotfill \quad $-0.03$
\item $h(x,y) = 2x^2+7y^2-3xy-2x$\quad \dotfill\; ???\;\dotfill \quad $\nicefrac{-27}{47}$
\item $y = -x+5$\quad \dotfill\; ???\;\dotfill \quad $\left[\frac{29}{8}\;,\;\frac{11}{8}\right]$
\end{enumerate}
\end{varwidth}
\end{center}
\end{minipage}
\begin{minipage}{0.20\linewidth}
\begin{center}
{\Huge\bfseries 3.} \\[2mm]
\includegraphics[height=37mm]{images/braille.png}
{\small Braillove písmeno}
\end{center}
\end{minipage}
\end{center}
\end{minipage}
&
\begin{minipage}[c][104.5mm][t]{0.5\linewidth}
\begin{center}
\vspace{7mm}
{\huge \LARGE Funkcia dvoch premenných, skupina \textit{Iota $\iota$} -\romannumeral4}\\[5mm]
\begin{minipage}{0.95\linewidth}
\begin{center}
{\small \begin{multicols}{2} \raggedright
\textbf{(a)} Vypočítajte zadanú parciálnu deriváciu funkcie \(f(x,y)\). \\[1mm]
\textbf{(b)} Určte gradient \(f(x,y)\) v bode \(A\). \\[1mm]
\textbf{(c)} Napíšte rovnicu dotykovej roviny ku grafu funkcie \(g(x,y)\) v dotykovom bode \(D\). \par \columnbreak
\textbf{(d)} Pomocou lineárnej aproximácie odhadnite funkčnú hodnotu funkcie \(g(x,y)\) v bode, ktorý je posunutý o vektor \((\Delta x, \Delta y)\) od bodu \(D_0\). \\[1mm]
\textbf{(e)} Vypočítajte \(h(x,y)\) v jej lokálnom extréme. \\[1mm]
\textbf{(f)} Nájdite bod viazeného lokálneho extrému funkcie \(h(x,y)\) na priamke. \\[1mm]
\end{multicols}
}
\end{center}
\end{minipage}
\\[4mm]
\begin{minipage}{0.79\linewidth}
\begin{center}
\begin{varwidth}{\linewidth}
\begin{enumerate}
\normalsizerr
\item $f(x,y) = \sqrt[3]{x^2y + 7xy^2}\enspace , \enspace \pdv{f}{x}$\quad \dotfill\; ???\;\dotfill \quad $\frac{2xy+7y^2}{3\,\sqrt[3]{(x^2y+7xy^2)^{2}}}$
\item $A = [3, 2]$\quad \dotfill\; ???\;\dotfill \quad $\left( \frac{80}{3\,\sqrt[3]{102^2}} \;,\; \frac{93}{3\,\sqrt[3]{102^2}} \right)$
\item $g(x,y) = \sin\frac{x}{y} \enspace , \enspace D = [0, -2, ?]$\quad \dotfill\; ???\;\dotfill \quad $x+2z=0$
\item $\Delta x = {0.2} \enspace,\enspace \Delta y = {-0.03}$\quad \dotfill\; ???\;\dotfill \quad $0.015$
\item $h(x,y) = -2x^2-3y^2-3xy+3x$\quad \dotfill\; ???\;\dotfill \quad $\nicefrac{9}{11}$
\item $y = 5x-1$\quad \dotfill\; ???\;\dotfill \quad $\left[\frac{9}{46}\;,\;\frac{-1}{31}\right]$
\end{enumerate}
\end{varwidth}
\end{center}
\end{minipage}
\begin{minipage}{0.20\linewidth}
\begin{center}
{\Huge\bfseries 4.} \\[2mm]
\includegraphics[height=37mm]{images/braille.png}
{\small Braillove písmeno}
\end{center}
\end{minipage}
\end{center}
\end{minipage}
%
\end{tabular}
\newpage
\thispagestyle{empty}
\begin{tabular}{c:c}
\begin{minipage}[c][104.5mm][t]{0.5\linewidth}
\begin{center}
\vspace{7mm}
{\huge \LARGE Funkcia dvoch premenných, skupina \textit{Kappa $\kappa$} -\romannumeral1}\\[5mm]
\begin{minipage}{0.95\linewidth}
\begin{center}
{\small \begin{multicols}{2} \raggedright
\textbf{(a)} Vypočítajte zadanú parciálnu deriváciu funkcie \(f(x,y)\). \\[1mm]
\textbf{(b)} Určte gradient \(f(x,y)\) v bode \(A\). \\[1mm]
\textbf{(c)} Napíšte rovnicu dotykovej roviny ku grafu funkcie \(g(x,y)\) v dotykovom bode \(D\). \par \columnbreak
\textbf{(d)} Pomocou lineárnej aproximácie odhadnite funkčnú hodnotu funkcie \(g(x,y)\) v bode, ktorý je posunutý o vektor \((\Delta x, \Delta y)\) od bodu \(D_0\). \\[1mm]
\textbf{(e)} Vypočítajte \(h(x,y)\) v jej lokálnom extréme. \\[1mm]
\textbf{(f)} Nájdite bod viazeného lokálneho extrému funkcie \(h(x,y)\) na priamke. \\[1mm]
\end{multicols}
}
\end{center}
\end{minipage}
\\[4mm]
\begin{minipage}{0.79\linewidth}
\begin{center}
\begin{varwidth}{\linewidth}
\begin{enumerate}
\normalsizerr
\item $f(x,y) = \sqrt{x^2y + 5xy^2}\enspace , \enspace \pdv{f}{x}$\quad \dotfill\; ???\;\dotfill \quad $\frac{2xy+5y^2}{2\,\sqrt{x^2y+5xy^2}}$
\item $A = [4, 2]$\quad \dotfill\; ???\;\dotfill \quad $\left( \frac{36}{2\,\sqrt{112}} \;,\; \frac{192}{2\,\sqrt{112}} \right)$
\item $g(x,y) = \sin\frac{x}{y} \enspace , \enspace D = [0, -2, ?]$\quad \dotfill\; ???\;\dotfill \quad $x+2z=0$
\item $\Delta x = {-0.1} \enspace,\enspace \Delta y = {0.3}$\quad \dotfill\; ???\;\dotfill \quad $-0.15$
\item $h(x,y) = -2x^2-4y^2+3xy+2x$\quad \dotfill\; ???\;\dotfill \quad $\nicefrac{17}{23}$
\item $y = 3x-1$\quad \dotfill\; ???\;\dotfill \quad $\left[\frac{-25}{58}\;,\;\frac{11}{58}\right]$
\end{enumerate}
\end{varwidth}
\end{center}
\end{minipage}
\begin{minipage}{0.20\linewidth}
\begin{center}
{\Huge\bfseries 1.} \\[2mm]
\includegraphics[height=37mm]{images/braille.png}
{\small Braillove písmeno}
\end{center}
\end{minipage}
\end{center}
\end{minipage}
&
\begin{minipage}[c][104.5mm][t]{0.5\linewidth}
\begin{center}
\vspace{7mm}
{\huge \LARGE Funkcia dvoch premenných, skupina \textit{Kappa $\kappa$} -\romannumeral2}\\[5mm]
\begin{minipage}{0.95\linewidth}
\begin{center}
{\small \begin{multicols}{2} \raggedright
\textbf{(a)} Vypočítajte zadanú parciálnu deriváciu funkcie \(f(x,y)\). \\[1mm]
\textbf{(b)} Určte gradient \(f(x,y)\) v bode \(A\). \\[1mm]
\textbf{(c)} Napíšte rovnicu dotykovej roviny ku grafu funkcie \(g(x,y)\) v dotykovom bode \(D\). \par \columnbreak
\textbf{(d)} Pomocou lineárnej aproximácie odhadnite funkčnú hodnotu funkcie \(g(x,y)\) v bode, ktorý je posunutý o vektor \((\Delta x, \Delta y)\) od bodu \(D_0\). \\[1mm]
\textbf{(e)} Vypočítajte \(h(x,y)\) v jej lokálnom extréme. \\[1mm]
\textbf{(f)} Nájdite bod viazeného lokálneho extrému funkcie \(h(x,y)\) na priamke. \\[1mm]
\end{multicols}
}
\end{center}
\end{minipage}
\\[4mm]
\begin{minipage}{0.79\linewidth}
\begin{center}
\begin{varwidth}{\linewidth}
\begin{enumerate}
\normalsizerr
\item $f(x,y) = \sqrt{6x^2y + 3xy^2}\enspace , \enspace \pdv{f}{y}$\quad \dotfill\; ???\;\dotfill \quad $\frac{6x^2+6xy}{2\,\sqrt{6x^2y+3xy^2}}$
\item $A = [4, 4]$\quad \dotfill\; ???\;\dotfill \quad $\left( \frac{480}{2\,\sqrt{576}} \;,\; \frac{192}{2\,\sqrt{576}} \right)$
\item $g(x,y) = \sin\frac{x}{y} \enspace , \enspace D = [0, -1, ?]$\quad \dotfill\; ???\;\dotfill \quad $x+z=0$
\item $\Delta x = {-0.3} \enspace,\enspace \Delta y = {0.3}$\quad \dotfill\; ???\;\dotfill \quad $-0.3$
\item $h(x,y) = 3x^2+y^2+xy+3x$\quad \dotfill\; ???\;\dotfill \quad $\nicefrac{-9}{11}$
\item $y = -4x+4$\quad \dotfill\; ???\;\dotfill \quad $\left[\frac{5}{6}\;,\;\frac{10}{23}\right]$
\end{enumerate}
\end{varwidth}
\end{center}
\end{minipage}
\begin{minipage}{0.20\linewidth}
\begin{center}
{\Huge\bfseries 2.} \\[2mm]
\includegraphics[height=37mm]{images/braille.png}
{\small Braillove písmeno}
\end{center}
\end{minipage}
\end{center}
\end{minipage}
\\ \hdashline
\begin{minipage}[c][104.5mm][t]{0.5\linewidth}
\begin{center}
\vspace{7mm}
{\huge \LARGE Funkcia dvoch premenných, skupina \textit{Kappa $\kappa$} -\romannumeral3}\\[5mm]
\begin{minipage}{0.95\linewidth}
\begin{center}
{\small \begin{multicols}{2} \raggedright
\textbf{(a)} Vypočítajte zadanú parciálnu deriváciu funkcie \(f(x,y)\). \\[1mm]
\textbf{(b)} Určte gradient \(f(x,y)\) v bode \(A\). \\[1mm]
\textbf{(c)} Napíšte rovnicu dotykovej roviny ku grafu funkcie \(g(x,y)\) v dotykovom bode \(D\). \par \columnbreak
\textbf{(d)} Pomocou lineárnej aproximácie odhadnite funkčnú hodnotu funkcie \(g(x,y)\) v bode, ktorý je posunutý o vektor \((\Delta x, \Delta y)\) od bodu \(D_0\). \\[1mm]
\textbf{(e)} Vypočítajte \(h(x,y)\) v jej lokálnom extréme. \\[1mm]
\textbf{(f)} Nájdite bod viazeného lokálneho extrému funkcie \(h(x,y)\) na priamke. \\[1mm]
\end{multicols}
}
\end{center}
\end{minipage}
\\[4mm]
\begin{minipage}{0.79\linewidth}
\begin{center}
\begin{varwidth}{\linewidth}
\begin{enumerate}
\normalsizerr
\item $f(x,y) = \sqrt{4x^2y + 3xy^2}\enspace , \enspace \pdv{f}{y}$\quad \dotfill\; ???\;\dotfill \quad $\frac{4x^2+6xy}{\sqrt{4x^2y+6xy^2}}$
\item $A = [3, 1]$\quad \dotfill\; ???\;\dotfill \quad $\left( \frac{27}{2\,\sqrt{45}} \;,\; \frac{54}{2\,\sqrt{45}} \right)$
\item $g(x,y) = \sin\frac{x}{y} \enspace , \enspace D = [0, 3, ?]$\quad \dotfill\; ???\;\dotfill \quad $x-3z=0$
\item $\Delta x = {0.02} \enspace,\enspace \Delta y = {0.02}$\quad \dotfill\; ???\;\dotfill \quad $0.06$
\item $h(x,y) = -6x^2-5y^2+3xy-2x$\quad \dotfill\; ???\;\dotfill \quad $\nicefrac{20}{129}$
\item $y = -2x+4$\quad \dotfill\; ???\;\dotfill \quad $\left[\frac{45}{32}\;,\;\frac{19}{16}\right]$
\end{enumerate}
\end{varwidth}
\end{center}
\end{minipage}
\begin{minipage}{0.20\linewidth}
\begin{center}
{\Huge\bfseries 3.} \\[2mm]
\includegraphics[height=37mm]{images/braille.png}
{\small Braillove písmeno}
\end{center}
\end{minipage}
\end{center}
\end{minipage}
&
\begin{minipage}[c][104.5mm][t]{0.5\linewidth}
\begin{center}
\vspace{7mm}
{\huge \LARGE Funkcia dvoch premenných, skupina \textit{Kappa $\kappa$} -\romannumeral4}\\[5mm]
\begin{minipage}{0.95\linewidth}
\begin{center}
{\small \begin{multicols}{2} \raggedright
\textbf{(a)} Vypočítajte zadanú parciálnu deriváciu funkcie \(f(x,y)\). \\[1mm]
\textbf{(b)} Určte gradient \(f(x,y)\) v bode \(A\). \\[1mm]
\textbf{(c)} Napíšte rovnicu dotykovej roviny ku grafu funkcie \(g(x,y)\) v dotykovom bode \(D\). \par \columnbreak
\textbf{(d)} Pomocou lineárnej aproximácie odhadnite funkčnú hodnotu funkcie \(g(x,y)\) v bode, ktorý je posunutý o vektor \((\Delta x, \Delta y)\) od bodu \(D_0\). \\[1mm]
\textbf{(e)} Vypočítajte \(h(x,y)\) v jej lokálnom extréme. \\[1mm]
\textbf{(f)} Nájdite bod viazeného lokálneho extrému funkcie \(h(x,y)\) na priamke. \\[1mm]
\end{multicols}
}
\end{center}
\end{minipage}
\\[4mm]
\begin{minipage}{0.79\linewidth}
\begin{center}
\begin{varwidth}{\linewidth}
\begin{enumerate}
\normalsizerr
\item $f(x,y) = \sqrt[4]{3x^2y + 3xy^2}\enspace , \enspace \pdv{f}{y}$\quad \dotfill\; ???\;\dotfill \quad $\frac{3x^2+6xy}{4\,\sqrt[4]{(3x^2y+3xy^2)^{3}}}$
\item $A = [3, 2]$\quad \dotfill\; ???\;\dotfill \quad $\left( \frac{48}{4\,\sqrt[4]{90^3}} \;,\; \frac{63}{4\,\sqrt[4]{90^3}} \right)$
\item $g(x,y) = \sin\frac{x}{y} \enspace , \enspace D = [0, -3, ?]$\quad \dotfill\; ???\;\dotfill \quad $x-3z=0$
\item $\Delta x = {0.1} \enspace,\enspace \Delta y = {-0.2}$\quad \dotfill\; ???\;\dotfill \quad $-0.033$
\item $h(x,y) = 7x^2+2y^2+5xy-4x$\quad \dotfill\; ???\;\dotfill \quad $\nicefrac{-32}{31}$
\item $y = x+2$\quad \dotfill\; ???\;\dotfill \quad $\left[\frac{1}{14}\;,\;\frac{3}{2}\right]$
\end{enumerate}
\end{varwidth}
\end{center}
\end{minipage}
\begin{minipage}{0.20\linewidth}
\begin{center}
{\Huge\bfseries 4.} \\[2mm]
\includegraphics[height=37mm]{images/braille.png}
{\small Braillove písmeno}
\end{center}
\end{minipage}
\end{center}
\end{minipage}
%
\end{tabular}
\newpage
\thispagestyle{empty}
\begin{tabular}{c:c}
\begin{minipage}[c][104.5mm][t]{0.5\linewidth}
\begin{center}
\vspace{7mm}
{\huge \LARGE Funkcia dvoch premenných, skupina \textit{Lambda $\lambda$} -\romannumeral1}\\[5mm]
\begin{minipage}{0.95\linewidth}
\begin{center}
{\small \begin{multicols}{2} \raggedright
\textbf{(a)} Vypočítajte zadanú parciálnu deriváciu funkcie \(f(x,y)\). \\[1mm]
\textbf{(b)} Určte gradient \(f(x,y)\) v bode \(A\). \\[1mm]
\textbf{(c)} Napíšte rovnicu dotykovej roviny ku grafu funkcie \(g(x,y)\) v dotykovom bode \(D\). \par \columnbreak
\textbf{(d)} Pomocou lineárnej aproximácie odhadnite funkčnú hodnotu funkcie \(g(x,y)\) v bode, ktorý je posunutý o vektor \((\Delta x, \Delta y)\) od bodu \(D_0\). \\[1mm]
\textbf{(e)} Vypočítajte \(h(x,y)\) v jej lokálnom extréme. \\[1mm]
\textbf{(f)} Nájdite bod viazeného lokálneho extrému funkcie \(h(x,y)\) na priamke. \\[1mm]
\end{multicols}
}
\end{center}
\end{minipage}
\\[4mm]
\begin{minipage}{0.79\linewidth}
\begin{center}
\begin{varwidth}{\linewidth}
\begin{enumerate}
\normalsizerr
\item $f(x,y) = \sqrt[3]{x^2y + 2xy^2}\enspace , \enspace \pdv{f}{y}$\quad \dotfill\; ???\;\dotfill \quad $\frac{x^2-4xy}{3\,\sqrt[3]{(x^2y+2xy^2)^{2}}}$
\item $A = [2, 2]$\quad \dotfill\; ???\;\dotfill \quad $\left( \frac{16}{3\,\sqrt[3]{24^2}} \;,\; \frac{20}{3\,\sqrt[3]{24^2}} \right)$
\item $g(x,y) = \ln\frac{x}{y} \enspace , \enspace D=[1,1,?]$\quad \dotfill\; ???\;\dotfill \quad $x + y + z + 0 = 0$
\item $\Delta x = {-0.1} \enspace,\enspace \Delta y = {0.2}$\quad \dotfill\; ???\;\dotfill \quad $0.0$
\item $h(x,y) = -2x^2-2y^2+3xy+3x$\quad \dotfill\; ???\;\dotfill \quad $\nicefrac{19}{7}$
\item $y = -2x+3$\quad \dotfill\; ???\;\dotfill \quad $\left[\frac{9}{8}\;,\;\frac{3}{4}\right]$
\end{enumerate}
\end{varwidth}
\end{center}
\end{minipage}
\begin{minipage}{0.20\linewidth}
\begin{center}
{\Huge\bfseries 1.} \\[2mm]
\includegraphics[height=37mm]{images/braille.png}
{\small Braillove písmeno}
\end{center}
\end{minipage}
\end{center}
\end{minipage}
&
\begin{minipage}[c][104.5mm][t]{0.5\linewidth}
\begin{center}
\vspace{7mm}
{\huge \LARGE Funkcia dvoch premenných, skupina \textit{Lambda $\lambda$} -\romannumeral2}\\[5mm]
\begin{minipage}{0.95\linewidth}
\begin{center}
{\small \begin{multicols}{2} \raggedright
\textbf{(a)} Vypočítajte zadanú parciálnu deriváciu funkcie \(f(x,y)\). \\[1mm]
\textbf{(b)} Určte gradient \(f(x,y)\) v bode \(A\). \\[1mm]
\textbf{(c)} Napíšte rovnicu dotykovej roviny ku grafu funkcie \(g(x,y)\) v dotykovom bode \(D\). \par \columnbreak
\textbf{(d)} Pomocou lineárnej aproximácie odhadnite funkčnú hodnotu funkcie \(g(x,y)\) v bode, ktorý je posunutý o vektor \((\Delta x, \Delta y)\) od bodu \(D_0\). \\[1mm]
\textbf{(e)} Vypočítajte \(h(x,y)\) v jej lokálnom extréme. \\[1mm]
\textbf{(f)} Nájdite bod viazeného lokálneho extrému funkcie \(h(x,y)\) na priamke. \\[1mm]
\end{multicols}
}
\end{center}
\end{minipage}
\\[4mm]
\begin{minipage}{0.79\linewidth}
\begin{center}
\begin{varwidth}{\linewidth}
\begin{enumerate}
\normalsizerr
\item $f(x,y) = \sqrt[3]{2x^2y + 7xy^2}\enspace , \enspace \pdv{f}{y}$\quad \dotfill\; ???\;\dotfill \quad $\frac{2x^2+14xy}{3\,\sqrt[3]{(2x^2y+7xy^2)^{2}}}$
\item $A = [4, 3]$\quad \dotfill\; ???\;\dotfill \quad $\left( \frac{111}{3\,\sqrt[3]{348^2}} \;,\; \frac{200}{3\,\sqrt[3]{348^2}} \right)$
\item $g(x,y) = \sin\frac{x}{y} \enspace , \enspace D = [0, 3, ?]$\quad \dotfill\; ???\;\dotfill \quad $x-3z=0$
\item $\Delta x = {-0.2} \enspace,\enspace \Delta y = {-0.3}$\quad \dotfill\; ???\;\dotfill \quad $-0.067$
\item $h(x,y) = -4x^2-y^2-3xy+2x$\quad \dotfill\; ???\;\dotfill \quad $\nicefrac{4}{25}$
\item $y = 2x-1$\quad \dotfill\; ???\;\dotfill \quad $\left[\frac{1}{28}\;,\;\frac{-5}{14}\right]$
\end{enumerate}
\end{varwidth}
\end{center}
\end{minipage}
\begin{minipage}{0.20\linewidth}
\begin{center}
{\Huge\bfseries 2.} \\[2mm]
\includegraphics[height=37mm]{images/braille.png}
{\small Braillove písmeno}
\end{center}
\end{minipage}
\end{center}
\end{minipage}
\\ \hdashline
\begin{minipage}[c][104.5mm][t]{0.5\linewidth}
\begin{center}
\vspace{7mm}
{\huge \LARGE Funkcia dvoch premenných, skupina \textit{Lambda $\lambda$} -\romannumeral3}\\[5mm]
\begin{minipage}{0.95\linewidth}
\begin{center}
{\small \begin{multicols}{2} \raggedright
\textbf{(a)} Vypočítajte zadanú parciálnu deriváciu funkcie \(f(x,y)\). \\[1mm]
\textbf{(b)} Určte gradient \(f(x,y)\) v bode \(A\). \\[1mm]
\textbf{(c)} Napíšte rovnicu dotykovej roviny ku grafu funkcie \(g(x,y)\) v dotykovom bode \(D\). \par \columnbreak
\textbf{(d)} Pomocou lineárnej aproximácie odhadnite funkčnú hodnotu funkcie \(g(x,y)\) v bode, ktorý je posunutý o vektor \((\Delta x, \Delta y)\) od bodu \(D_0\). \\[1mm]
\textbf{(e)} Vypočítajte \(h(x,y)\) v jej lokálnom extréme. \\[1mm]
\textbf{(f)} Nájdite bod viazeného lokálneho extrému funkcie \(h(x,y)\) na priamke. \\[1mm]
\end{multicols}
}
\end{center}
\end{minipage}
\\[4mm]
\begin{minipage}{0.79\linewidth}
\begin{center}
\begin{varwidth}{\linewidth}
\begin{enumerate}
\normalsizerr
\item $f(x,y) = \sqrt[3]{7x^2y + 5xy^2}\enspace , \enspace \pdv{f}{x}$\quad \dotfill\; ???\;\dotfill \quad $\frac{14xy+5y^2}{3\,\sqrt[3]{(7x^2y+5xy^2)^{2}}}$
\item $A = [2, 1]$\quad \dotfill\; ???\;\dotfill \quad $\left( \frac{33}{3\,\sqrt[3]{38^2}} \;,\; \frac{96}{3\,\sqrt[3]{38^2}} \right)$
\item $g(x,y) = \sin\frac{x}{y} \enspace , \enspace D = [0, 1, ?]$\quad \dotfill\; ???\;\dotfill \quad $x-z=1$
\item $\Delta x = {0.2} \enspace,\enspace \Delta y = {0.1}$\quad \dotfill\; ???\;\dotfill \quad $0.1$
\item $h(x,y) = -5x^2-2y^2+4xy-4x$\quad \dotfill\; ???\;\dotfill \quad $\nicefrac{11}{8}$
\item $y = -x+2$\quad \dotfill\; ???\;\dotfill \quad $\left[\frac{6}{11}\;,\;\frac{16}{3}\right]$
\end{enumerate}
\end{varwidth}
\end{center}
\end{minipage}
\begin{minipage}{0.20\linewidth}
\begin{center}
{\Huge\bfseries 3.} \\[2mm]
\includegraphics[height=37mm]{images/braille.png}
{\small Braillove písmeno}
\end{center}
\end{minipage}
\end{center}
\end{minipage}
&
\begin{minipage}[c][104.5mm][t]{0.5\linewidth}
\begin{center}
\vspace{7mm}
{\huge \LARGE Funkcia dvoch premenných, skupina \textit{Lambda $\lambda$} -\romannumeral4}\\[5mm]
\begin{minipage}{0.95\linewidth}
\begin{center}
{\small \begin{multicols}{2} \raggedright
\textbf{(a)} Vypočítajte zadanú parciálnu deriváciu funkcie \(f(x,y)\). \\[1mm]
\textbf{(b)} Určte gradient \(f(x,y)\) v bode \(A\). \\[1mm]
\textbf{(c)} Napíšte rovnicu dotykovej roviny ku grafu funkcie \(g(x,y)\) v dotykovom bode \(D\). \par \columnbreak
\textbf{(d)} Pomocou lineárnej aproximácie odhadnite funkčnú hodnotu funkcie \(g(x,y)\) v bode, ktorý je posunutý o vektor \((\Delta x, \Delta y)\) od bodu \(D_0\). \\[1mm]
\textbf{(e)} Vypočítajte \(h(x,y)\) v jej lokálnom extréme. \\[1mm]
\textbf{(f)} Nájdite bod viazeného lokálneho extrému funkcie \(h(x,y)\) na priamke. \\[1mm]
\end{multicols}
}
\end{center}
\end{minipage}
\\[4mm]
\begin{minipage}{0.79\linewidth}
\begin{center}
\begin{varwidth}{\linewidth}
\begin{enumerate}
\normalsizerr
\item $f(x,y) = \sqrt{2x^2y + 6xy^2}\enspace , \enspace \pdv{f}{y}$\quad \dotfill\; ???\;\dotfill \quad $\frac{2x^2+12xy}{2\,\sqrt{2x^2y+6xy^2}}$
\item $A = [3, 1]$\quad \dotfill\; ???\;\dotfill \quad $\left( \frac{18}{2\,\sqrt{36}} \;,\; \frac{108}{2\,\sqrt{36}} \right)$
\item $g(x,y) = \ln\frac{x}{y} \enspace , \enspace D=[1,1,?]$\quad \dotfill\; ???\;\dotfill \quad $x - y - z + 0 = 0$
\item $\Delta x = {0.3} \enspace,\enspace \Delta y = {-0.01}$\quad \dotfill\; ???\;\dotfill \quad $0.29$
\item $h(x,y) = 7x^2+4y^2+4xy+x$\quad \dotfill\; ???\;\dotfill \quad $\nicefrac{-1}{24}$
\item $y = 2x+1$\quad \dotfill\; ???\;\dotfill \quad $\left[\frac{-21}{62}\;,\;\frac{10}{31}\right]$
\end{enumerate}
\end{varwidth}
\end{center}
\end{minipage}
\begin{minipage}{0.20\linewidth}
\begin{center}
{\Huge\bfseries 4.} \\[2mm]
\includegraphics[height=37mm]{images/braille.png}
{\small Braillove písmeno}
\end{center}
\end{minipage}
\end{center}
\end{minipage}
%
\end{tabular}
\newpage
\thispagestyle{empty}
\begin{tabular}{c:c}
\begin{minipage}[c][104.5mm][t]{0.5\linewidth}
\begin{center}
\vspace{7mm}
{\huge \LARGE Funkcia dvoch premenných, skupina \textit{Mu $\mu$} -\romannumeral1}\\[5mm]
\begin{minipage}{0.95\linewidth}
\begin{center}
{\small \begin{multicols}{2} \raggedright
\textbf{(a)} Vypočítajte zadanú parciálnu deriváciu funkcie \(f(x,y)\). \\[1mm]
\textbf{(b)} Určte gradient \(f(x,y)\) v bode \(A\). \\[1mm]
\textbf{(c)} Napíšte rovnicu dotykovej roviny ku grafu funkcie \(g(x,y)\) v dotykovom bode \(D\). \par \columnbreak
\textbf{(d)} Pomocou lineárnej aproximácie odhadnite funkčnú hodnotu funkcie \(g(x,y)\) v bode, ktorý je posunutý o vektor \((\Delta x, \Delta y)\) od bodu \(D_0\). \\[1mm]
\textbf{(e)} Vypočítajte \(h(x,y)\) v jej lokálnom extréme. \\[1mm]
\textbf{(f)} Nájdite bod viazeného lokálneho extrému funkcie \(h(x,y)\) na priamke. \\[1mm]
\end{multicols}
}
\end{center}
\end{minipage}
\\[4mm]
\begin{minipage}{0.79\linewidth}
\begin{center}
\begin{varwidth}{\linewidth}
\begin{enumerate}
\normalsizerr
\item $f(x,y) = \sqrt{7x^2y + 2xy^2}\enspace , \enspace \pdv{f}{x}$\quad \dotfill\; ???\;\dotfill \quad $\frac{14xy+2y^2}{2\,\sqrt{7x^2y+2xy^2}}$
\item $A = [4, 2]$\quad \dotfill\; ???\;\dotfill \quad $\left( \frac{120}{2\,\sqrt{256}} \;,\; \frac{144}{2\,\sqrt{256}} \right)$
\item $g(x,y) = \sin\frac{x}{y} \enspace , \enspace D = [0, -1, ?]$\quad \dotfill\; ???\;\dotfill \quad $x+z=0$
\item $\Delta x = {0.2} \enspace,\enspace \Delta y = {0.02}$\quad \dotfill\; ???\;\dotfill \quad $-0.02$
\item $h(x,y) = -3x^2-2y^2-2xy+3x$\quad \dotfill\; ???\;\dotfill \quad $\nicefrac{19}{20}$
\item $y = 3x+3$\quad \dotfill\; ???\;\dotfill \quad $\left[\frac{-13}{18}\;,\;\frac{5}{6}\right]$
\end{enumerate}
\end{varwidth}
\end{center}
\end{minipage}
\begin{minipage}{0.20\linewidth}
\begin{center}
{\Huge\bfseries 1.} \\[2mm]
\includegraphics[height=37mm]{images/braille.png}
{\small Braillove písmeno}
\end{center}
\end{minipage}
\end{center}
\end{minipage}
&
\begin{minipage}[c][104.5mm][t]{0.5\linewidth}
\begin{center}
\vspace{7mm}
{\huge \LARGE Funkcia dvoch premenných, skupina \textit{Mu $\mu$} -\romannumeral2}\\[5mm]
\begin{minipage}{0.95\linewidth}
\begin{center}
{\small \begin{multicols}{2} \raggedright
\textbf{(a)} Vypočítajte zadanú parciálnu deriváciu funkcie \(f(x,y)\). \\[1mm]
\textbf{(b)} Určte gradient \(f(x,y)\) v bode \(A\). \\[1mm]
\textbf{(c)} Napíšte rovnicu dotykovej roviny ku grafu funkcie \(g(x,y)\) v dotykovom bode \(D\). \par \columnbreak
\textbf{(d)} Pomocou lineárnej aproximácie odhadnite funkčnú hodnotu funkcie \(g(x,y)\) v bode, ktorý je posunutý o vektor \((\Delta x, \Delta y)\) od bodu \(D_0\). \\[1mm]
\textbf{(e)} Vypočítajte \(h(x,y)\) v jej lokálnom extréme. \\[1mm]
\textbf{(f)} Nájdite bod viazeného lokálneho extrému funkcie \(h(x,y)\) na priamke. \\[1mm]
\end{multicols}
}
\end{center}
\end{minipage}
\\[4mm]
\begin{minipage}{0.79\linewidth}
\begin{center}
\begin{varwidth}{\linewidth}
\begin{enumerate}
\normalsizerr
\item $f(x,y) = \sqrt[4]{6x^2y + xy^2}\enspace , \enspace \pdv{f}{x}$\quad \dotfill\; ???\;\dotfill \quad $\frac{12xy-y^2}{4\,\sqrt[4]{(6x^2y+xy^2)^{3}}}$
\item $A = [2, 1]$\quad \dotfill\; ???\;\dotfill \quad $\left( \frac{25}{4\,\sqrt[4]{26^3}} \;,\; \frac{28}{4\,\sqrt[4]{26^3}} \right)$
\item $g(x,y) = \ln\frac{x}{y} \enspace , \enspace D=[e^{5},1,?]$\quad \dotfill\; ???\;\dotfill \quad $\frac{x}{e^6} - y - z + 6 = 0$
\item $\Delta x = {-0.1} \enspace,\enspace \Delta y = {0.03}$\quad \dotfill\; ???\;\dotfill \quad $5.099$
\item $h(x,y) = -5x^2-4y^2+2xy+3x$\quad \dotfill\; ???\;\dotfill \quad $\nicefrac{3}{7}$
\item $y = x-1$\quad \dotfill\; ???\;\dotfill \quad $\left[\frac{9}{14}\;,\;\frac{-5}{14}\right]$
\end{enumerate}
\end{varwidth}
\end{center}
\end{minipage}
\begin{minipage}{0.20\linewidth}
\begin{center}
{\Huge\bfseries 2.} \\[2mm]
\includegraphics[height=37mm]{images/braille.png}
{\small Braillove písmeno}
\end{center}
\end{minipage}
\end{center}
\end{minipage}
\\ \hdashline
\begin{minipage}[c][104.5mm][t]{0.5\linewidth}
\begin{center}
\vspace{7mm}
{\huge \LARGE Funkcia dvoch premenných, skupina \textit{Mu $\mu$} -\romannumeral3}\\[5mm]
\begin{minipage}{0.95\linewidth}
\begin{center}
{\small \begin{multicols}{2} \raggedright
\textbf{(a)} Vypočítajte zadanú parciálnu deriváciu funkcie \(f(x,y)\). \\[1mm]
\textbf{(b)} Určte gradient \(f(x,y)\) v bode \(A\). \\[1mm]
\textbf{(c)} Napíšte rovnicu dotykovej roviny ku grafu funkcie \(g(x,y)\) v dotykovom bode \(D\). \par \columnbreak
\textbf{(d)} Pomocou lineárnej aproximácie odhadnite funkčnú hodnotu funkcie \(g(x,y)\) v bode, ktorý je posunutý o vektor \((\Delta x, \Delta y)\) od bodu \(D_0\). \\[1mm]
\textbf{(e)} Vypočítajte \(h(x,y)\) v jej lokálnom extréme. \\[1mm]
\textbf{(f)} Nájdite bod viazeného lokálneho extrému funkcie \(h(x,y)\) na priamke. \\[1mm]
\end{multicols}
}
\end{center}
\end{minipage}
\\[4mm]
\begin{minipage}{0.79\linewidth}
\begin{center}
\begin{varwidth}{\linewidth}
\begin{enumerate}
\normalsizerr
\item $f(x,y) = \sqrt[4]{x^2y + 4xy^2}\enspace , \enspace \pdv{f}{x}$\quad \dotfill\; ???\;\dotfill \quad $\frac{2xy+4y^2}{4\,\sqrt[4]{(x^2y+4xy^2)^{3}}}$
\item $A = [2, 2]$\quad \dotfill\; ???\;\dotfill \quad $\left( \frac{24}{4\,\sqrt[4]{40^3}} \;,\; \frac{36}{4\,\sqrt[4]{40^3}} \right)$
\item $g(x,y) = \sin\frac{x}{y} \enspace , \enspace D = [0, 1, ?]$\quad \dotfill\; ???\;\dotfill \quad $x-z=0$
\item $\Delta x = {0.01} \enspace,\enspace \Delta y = {0.03}$\quad \dotfill\; ???\;\dotfill \quad $0.01$
\item $h(x,y) = -x^2-3y^2-2xy-5x$\quad \dotfill\; ???\;\dotfill \quad $\nicefrac{19}{2}$
\item $y = -2x-1$\quad \dotfill\; ???\;\dotfill \quad $\left[\frac{-5}{6}\;,\;\frac{6}{17}\right]$
\end{enumerate}
\end{varwidth}
\end{center}
\end{minipage}
\begin{minipage}{0.20\linewidth}
\begin{center}
{\Huge\bfseries 3.} \\[2mm]
\includegraphics[height=37mm]{images/braille.png}
{\small Braillove písmeno}
\end{center}
\end{minipage}
\end{center}
\end{minipage}
&
\begin{minipage}[c][104.5mm][t]{0.5\linewidth}
\begin{center}
\vspace{7mm}
{\huge \LARGE Funkcia dvoch premenných, skupina \textit{Mu $\mu$} -\romannumeral4}\\[5mm]
\begin{minipage}{0.95\linewidth}
\begin{center}
{\small \begin{multicols}{2} \raggedright
\textbf{(a)} Vypočítajte zadanú parciálnu deriváciu funkcie \(f(x,y)\). \\[1mm]
\textbf{(b)} Určte gradient \(f(x,y)\) v bode \(A\). \\[1mm]
\textbf{(c)} Napíšte rovnicu dotykovej roviny ku grafu funkcie \(g(x,y)\) v dotykovom bode \(D\). \par \columnbreak
\textbf{(d)} Pomocou lineárnej aproximácie odhadnite funkčnú hodnotu funkcie \(g(x,y)\) v bode, ktorý je posunutý o vektor \((\Delta x, \Delta y)\) od bodu \(D_0\). \\[1mm]
\textbf{(e)} Vypočítajte \(h(x,y)\) v jej lokálnom extréme. \\[1mm]
\textbf{(f)} Nájdite bod viazeného lokálneho extrému funkcie \(h(x,y)\) na priamke. \\[1mm]
\end{multicols}
}
\end{center}
\end{minipage}
\\[4mm]
\begin{minipage}{0.79\linewidth}
\begin{center}
\begin{varwidth}{\linewidth}
\begin{enumerate}
\normalsizerr
\item $f(x,y) = \sqrt{5x^2y + 4xy^2}\enspace , \enspace \pdv{f}{y}$\quad \dotfill\; ???\;\dotfill \quad $\frac{5x^2+8xy}{2\,\sqrt{5x^2y+4xy^2}}$
\item $A = [4, 2]$\quad \dotfill\; ???\;\dotfill \quad $\left( \frac{96}{2\,\sqrt{224}} \;,\; \frac{288}{2\,\sqrt{224}} \right)$
\item $g(x,y) = \sin\frac{x}{y} \enspace , \enspace D = [0, -2, ?]$\quad \dotfill\; ???\;\dotfill \quad $x+2z=0$
\item $\Delta x = {0.3} \enspace,\enspace \Delta y = {-0.02}$\quad \dotfill\; ???\;\dotfill \quad $0.01$
\item $h(x,y) = 2x^2+5y^2-2xy-3x$\quad \dotfill\; ???\;\dotfill \quad $\nicefrac{-5}{4}$
\item $y = -3x+2$\quad \dotfill\; ???\;\dotfill \quad $\left[\frac{-1}{2}\;,\;\frac{11}{106}\right]$
\end{enumerate}
\end{varwidth}
\end{center}
\end{minipage}
\begin{minipage}{0.20\linewidth}
\begin{center}
{\Huge\bfseries 4.} \\[2mm]
\includegraphics[height=37mm]{images/braille.png}
{\small Braillove písmeno}
\end{center}
\end{minipage}
\end{center}
\end{minipage}
%
\end{tabular}
\newpage
\thispagestyle{empty}
\begin{tabular}{c:c}
\begin{minipage}[c][104.5mm][t]{0.5\linewidth}
\begin{center}
\vspace{7mm}
{\huge \LARGE Funkcia dvoch premenných, skupina \textit{Nu $\nu$} -\romannumeral1}\\[5mm]
\begin{minipage}{0.95\linewidth}
\begin{center}
{\small \begin{multicols}{2} \raggedright
\textbf{(a)} Vypočítajte zadanú parciálnu deriváciu funkcie \(f(x,y)\). \\[1mm]
\textbf{(b)} Určte gradient \(f(x,y)\) v bode \(A\). \\[1mm]
\textbf{(c)} Napíšte rovnicu dotykovej roviny ku grafu funkcie \(g(x,y)\) v dotykovom bode \(D\). \par \columnbreak
\textbf{(d)} Pomocou lineárnej aproximácie odhadnite funkčnú hodnotu funkcie \(g(x,y)\) v bode, ktorý je posunutý o vektor \((\Delta x, \Delta y)\) od bodu \(D_0\). \\[1mm]
\textbf{(e)} Vypočítajte \(h(x,y)\) v jej lokálnom extréme. \\[1mm]
\textbf{(f)} Nájdite bod viazeného lokálneho extrému funkcie \(h(x,y)\) na priamke. \\[1mm]
\end{multicols}
}
\end{center}
\end{minipage}
\\[4mm]
\begin{minipage}{0.79\linewidth}
\begin{center}
\begin{varwidth}{\linewidth}
\begin{enumerate}
\normalsizerr
\item $f(x,y) = \sqrt{5x^2y + 2xy^2}\enspace , \enspace \pdv{f}{x}$\quad \dotfill\; ???\;\dotfill \quad $\frac{10xy+2y^2}{2\,\sqrt{5x^2y+2xy^2}}$
\item $A = [4, 2]$\quad \dotfill\; ???\;\dotfill \quad $\left( \frac{88}{2\,\sqrt{192}} \;,\; \frac{112}{2\,\sqrt{192}} \right)$
\item $g(x,y) = \sin\frac{x}{y} \enspace , \enspace D = [0, 3, ?]$\quad \dotfill\; ???\;\dotfill \quad $x-3z=0$
\item $\Delta x = {0.03} \enspace,\enspace \Delta y = {-0.2}$\quad \dotfill\; ???\;\dotfill \quad $0.09$
\item $h(x,y) = -2x^2-5y^2+4xy+3x$\quad \dotfill\; ???\;\dotfill \quad $\nicefrac{23}{12}$
\item $y = -2x-4$\quad \dotfill\; ???\;\dotfill \quad $\left[\frac{-31}{20}\;,\;\frac{-9}{10}\right]$
\end{enumerate}
\end{varwidth}
\end{center}
\end{minipage}
\begin{minipage}{0.20\linewidth}
\begin{center}
{\Huge\bfseries 1.} \\[2mm]
\includegraphics[height=37mm]{images/braille.png}
{\small Braillove písmeno}
\end{center}
\end{minipage}
\end{center}
\end{minipage}
&
\begin{minipage}[c][104.5mm][t]{0.5\linewidth}
\begin{center}
\vspace{7mm}
{\huge \LARGE Funkcia dvoch premenných, skupina \textit{Nu $\nu$} -\romannumeral2}\\[5mm]
\begin{minipage}{0.95\linewidth}
\begin{center}
{\small \begin{multicols}{2} \raggedright
\textbf{(a)} Vypočítajte zadanú parciálnu deriváciu funkcie \(f(x,y)\). \\[1mm]
\textbf{(b)} Určte gradient \(f(x,y)\) v bode \(A\). \\[1mm]
\textbf{(c)} Napíšte rovnicu dotykovej roviny ku grafu funkcie \(g(x,y)\) v dotykovom bode \(D\). \par \columnbreak
\textbf{(d)} Pomocou lineárnej aproximácie odhadnite funkčnú hodnotu funkcie \(g(x,y)\) v bode, ktorý je posunutý o vektor \((\Delta x, \Delta y)\) od bodu \(D_0\). \\[1mm]
\textbf{(e)} Vypočítajte \(h(x,y)\) v jej lokálnom extréme. \\[1mm]
\textbf{(f)} Nájdite bod viazeného lokálneho extrému funkcie \(h(x,y)\) na priamke. \\[1mm]
\end{multicols}
}
\end{center}
\end{minipage}
\\[4mm]
\begin{minipage}{0.79\linewidth}
\begin{center}
\begin{varwidth}{\linewidth}
\begin{enumerate}
\normalsizerr
\item $f(x,y) = \sqrt[4]{6x^2y + 4xy^2}\enspace , \enspace \pdv{f}{x}$\quad \dotfill\; ???\;\dotfill \quad $\frac{12xy+4y^2}{4\,\sqrt[4]{(6x^2y+4xy^2)^{3}}}$
\item $A = [4, 3]$\quad \dotfill\; ???\;\dotfill \quad $\left( \frac{180}{4\,\sqrt[4]{432^3}} \;,\; \frac{192}{4\,\sqrt[4]{432^3}} \right)$
\item $g(x,y) = \sin\frac{x}{y} \enspace , \enspace D = [0, -2, ?]$\quad \dotfill\; ???\;\dotfill \quad $x+2z=0$
\item $\Delta x = {-0.01} \enspace,\enspace \Delta y = {0.02}$\quad \dotfill\; ???\;\dotfill \quad $0.02$
\item $h(x,y) = 4x^2+4y^2+xy-2x$\quad \dotfill\; ???\;\dotfill \quad $\nicefrac{-16}{63}$
\item $y = x-5$\quad \dotfill\; ???\;\dotfill \quad $\left[\frac{-11}{6}\;,\;\frac{-43}{18}\right]$
\end{enumerate}
\end{varwidth}
\end{center}
\end{minipage}
\begin{minipage}{0.20\linewidth}
\begin{center}
{\Huge\bfseries 2.} \\[2mm]
\includegraphics[height=37mm]{images/braille.png}
{\small Braillove písmeno}
\end{center}
\end{minipage}
\end{center}
\end{minipage}
\\ \hdashline
\begin{minipage}[c][104.5mm][t]{0.5\linewidth}
\begin{center}
\vspace{7mm}
{\huge \LARGE Funkcia dvoch premenných, skupina \textit{Nu $\nu$} -\romannumeral3}\\[5mm]
\begin{minipage}{0.95\linewidth}
\begin{center}
{\small \begin{multicols}{2} \raggedright
\textbf{(a)} Vypočítajte zadanú parciálnu deriváciu funkcie \(f(x,y)\). \\[1mm]
\textbf{(b)} Určte gradient \(f(x,y)\) v bode \(A\). \\[1mm]
\textbf{(c)} Napíšte rovnicu dotykovej roviny ku grafu funkcie \(g(x,y)\) v dotykovom bode \(D\). \par \columnbreak
\textbf{(d)} Pomocou lineárnej aproximácie odhadnite funkčnú hodnotu funkcie \(g(x,y)\) v bode, ktorý je posunutý o vektor \((\Delta x, \Delta y)\) od bodu \(D_0\). \\[1mm]
\textbf{(e)} Vypočítajte \(h(x,y)\) v jej lokálnom extréme. \\[1mm]
\textbf{(f)} Nájdite bod viazeného lokálneho extrému funkcie \(h(x,y)\) na priamke. \\[1mm]
\end{multicols}
}
\end{center}
\end{minipage}
\\[4mm]
\begin{minipage}{0.79\linewidth}
\begin{center}
\begin{varwidth}{\linewidth}
\begin{enumerate}
\normalsizerr
\item $f(x,y) = \sqrt{2x^2y + 2xy^2}\enspace , \enspace \pdv{f}{x}$\quad \dotfill\; ???\;\dotfill \quad $\frac{2xy+2y^2}{2\,\sqrt{2x^2y+2xy^2}}$
\item $A = [2, 1]$\quad \dotfill\; ???\;\dotfill \quad $\left( \frac{10}{2\,\sqrt{12}} \;,\; \frac{16}{2\,\sqrt{12}} \right)$
\item $g(x,y) = \ln\frac{x}{y} \enspace , \enspace D=[e^{2},1,?]$\quad \dotfill\; ???\;\dotfill \quad $\frac{x}{e^2} - y - z + 2 = 0$
\item $\Delta x = {-0.3} \enspace,\enspace \Delta y = {-0.3}$\quad \dotfill\; ???\;\dotfill \quad $1.659$
\item $h(x,y) = -4x^2-y^2+xy+3x$\quad \dotfill\; ???\;\dotfill \quad $\nicefrac{9}{17}$
\item $y = -4x+1$\quad \dotfill\; ???\;\dotfill \quad $\left[\frac{1}{4}\;,\;0\right]$
\end{enumerate}
\end{varwidth}
\end{center}
\end{minipage}
\begin{minipage}{0.20\linewidth}
\begin{center}
{\Huge\bfseries 3.} \\[2mm]
\includegraphics[height=37mm]{images/braille.png}
{\small Braillove písmeno}
\end{center}
\end{minipage}
\end{center}
\end{minipage}
&
\begin{minipage}[c][104.5mm][t]{0.5\linewidth}
\begin{center}
\vspace{7mm}
{\huge \LARGE Funkcia dvoch premenných, skupina \textit{Nu $\nu$} -\romannumeral4}\\[5mm]
\begin{minipage}{0.95\linewidth}
\begin{center}
{\small \begin{multicols}{2} \raggedright
\textbf{(a)} Vypočítajte zadanú parciálnu deriváciu funkcie \(f(x,y)\). \\[1mm]
\textbf{(b)} Určte gradient \(f(x,y)\) v bode \(A\). \\[1mm]
\textbf{(c)} Napíšte rovnicu dotykovej roviny ku grafu funkcie \(g(x,y)\) v dotykovom bode \(D\). \par \columnbreak
\textbf{(d)} Pomocou lineárnej aproximácie odhadnite funkčnú hodnotu funkcie \(g(x,y)\) v bode, ktorý je posunutý o vektor \((\Delta x, \Delta y)\) od bodu \(D_0\). \\[1mm]
\textbf{(e)} Vypočítajte \(h(x,y)\) v jej lokálnom extréme. \\[1mm]
\textbf{(f)} Nájdite bod viazeného lokálneho extrému funkcie \(h(x,y)\) na priamke. \\[1mm]
\end{multicols}
}
\end{center}
\end{minipage}
\\[4mm]
\begin{minipage}{0.79\linewidth}
\begin{center}
\begin{varwidth}{\linewidth}
\begin{enumerate}
\normalsizerr
\item $f(x,y) = \sqrt[4]{2x^2y + 6xy^2}\enspace , \enspace \pdv{f}{x}$\quad \dotfill\; ???\;\dotfill \quad $\frac{4xy-6y^2}{4\,\sqrt[4]{(2x^2y+6xy^2)^{3}}}$
\item $A = [3, 4]$\quad \dotfill\; ???\;\dotfill \quad $\left( \frac{144}{4\,\sqrt[4]{360^3}} \;,\; \frac{162}{4\,\sqrt[4]{360^3}} \right)$
\item $g(x,y) = \ln\frac{x}{y} \enspace , \enspace D=[e^{3},1,?]$\quad \dotfill\; ???\;\dotfill \quad $\frac{x}{e^3} - y - z + 3 = 0$
\item $\Delta x = {-0.02} \enspace,\enspace \Delta y = {0.03}$\quad \dotfill\; ???\;\dotfill \quad $3.019$
\item $h(x,y) = -5x^2-5y^2-3xy+2x$\quad \dotfill\; ???\;\dotfill \quad $\nicefrac{20}{91}$
\item $y = -x-6$\quad \dotfill\; ???\;\dotfill \quad $\left[\frac{-20}{7}\;,\;\frac{-22}{7}\right]$
\end{enumerate}
\end{varwidth}
\end{center}
\end{minipage}
\begin{minipage}{0.20\linewidth}
\begin{center}
{\Huge\bfseries 4.} \\[2mm]
\includegraphics[height=37mm]{images/braille.png}
{\small Braillove písmeno}
\end{center}
\end{minipage}
\end{center}
\end{minipage}
%
\end{tabular}
\newpage
\thispagestyle{empty}
\begin{tabular}{c:c}
\begin{minipage}[c][104.5mm][t]{0.5\linewidth}
\begin{center}
\vspace{7mm}
{\huge \LARGE Funkcia dvoch premenných, skupina \textit{Xi $\xi$} -\romannumeral1}\\[5mm]
\begin{minipage}{0.95\linewidth}
\begin{center}
{\small \begin{multicols}{2} \raggedright
\textbf{(a)} Vypočítajte zadanú parciálnu deriváciu funkcie \(f(x,y)\). \\[1mm]
\textbf{(b)} Určte gradient \(f(x,y)\) v bode \(A\). \\[1mm]
\textbf{(c)} Napíšte rovnicu dotykovej roviny ku grafu funkcie \(g(x,y)\) v dotykovom bode \(D\). \par \columnbreak
\textbf{(d)} Pomocou lineárnej aproximácie odhadnite funkčnú hodnotu funkcie \(g(x,y)\) v bode, ktorý je posunutý o vektor \((\Delta x, \Delta y)\) od bodu \(D_0\). \\[1mm]
\textbf{(e)} Vypočítajte \(h(x,y)\) v jej lokálnom extréme. \\[1mm]
\textbf{(f)} Nájdite bod viazeného lokálneho extrému funkcie \(h(x,y)\) na priamke. \\[1mm]
\end{multicols}
}
\end{center}
\end{minipage}
\\[4mm]
\begin{minipage}{0.79\linewidth}
\begin{center}
\begin{varwidth}{\linewidth}
\begin{enumerate}
\normalsizerr
\item $f(x,y) = \sqrt[3]{6x^2y + 8xy^2}\enspace , \enspace \pdv{f}{x}$\quad \dotfill\; ???\;\dotfill \quad $\frac{12xy+8y^2}{3\,\sqrt[3]{(6x^2y+8xy^2)^{2}}}$
\item $A = [1, 3]$\quad \dotfill\; ???\;\dotfill \quad $\left( \frac{108}{3\,\sqrt[3]{90^2}} \;,\; \frac{54}{3\,\sqrt[3]{90^2}} \right)$
\item $g(x,y) = \ln\frac{x}{y} \enspace , \enspace D=[e^{2},1,?]$\quad \dotfill\; ???\;\dotfill \quad $\frac{x}{e^2} - y - z + 2 = 0$
\item $\Delta x = {-0.3} \enspace,\enspace \Delta y = {-0.2}$\quad \dotfill\; ???\;\dotfill \quad $1.759$
\item $h(x,y) = x^2+2y^2-xy-4x$\quad \dotfill\; ???\;\dotfill \quad $\nicefrac{-32}{7}$
\item $y = -7x+1$\quad \dotfill\; ???\;\dotfill \quad $\left[\frac{33}{212}\;,\;\frac{-19}{184}\right]$
\end{enumerate}
\end{varwidth}
\end{center}
\end{minipage}
\begin{minipage}{0.20\linewidth}
\begin{center}
{\Huge\bfseries 1.} \\[2mm]
\includegraphics[height=37mm]{images/braille.png}
{\small Braillove písmeno}
\end{center}
\end{minipage}
\end{center}
\end{minipage}
&
\begin{minipage}[c][104.5mm][t]{0.5\linewidth}
\begin{center}
\vspace{7mm}
{\huge \LARGE Funkcia dvoch premenných, skupina \textit{Xi $\xi$} -\romannumeral2}\\[5mm]
\begin{minipage}{0.95\linewidth}
\begin{center}
{\small \begin{multicols}{2} \raggedright
\textbf{(a)} Vypočítajte zadanú parciálnu deriváciu funkcie \(f(x,y)\). \\[1mm]
\textbf{(b)} Určte gradient \(f(x,y)\) v bode \(A\). \\[1mm]
\textbf{(c)} Napíšte rovnicu dotykovej roviny ku grafu funkcie \(g(x,y)\) v dotykovom bode \(D\). \par \columnbreak
\textbf{(d)} Pomocou lineárnej aproximácie odhadnite funkčnú hodnotu funkcie \(g(x,y)\) v bode, ktorý je posunutý o vektor \((\Delta x, \Delta y)\) od bodu \(D_0\). \\[1mm]
\textbf{(e)} Vypočítajte \(h(x,y)\) v jej lokálnom extréme. \\[1mm]
\textbf{(f)} Nájdite bod viazeného lokálneho extrému funkcie \(h(x,y)\) na priamke. \\[1mm]
\end{multicols}
}
\end{center}
\end{minipage}
\\[4mm]
\begin{minipage}{0.79\linewidth}
\begin{center}
\begin{varwidth}{\linewidth}
\begin{enumerate}
\normalsizerr
\item $f(x,y) = \sqrt[3]{2x^2y + 6xy^2}\enspace , \enspace \pdv{f}{x}$\quad \dotfill\; ???\;\dotfill \quad $\frac{4xy+6y^2}{3\,\sqrt[3]{(2x^2y+6xy^2)^{2}}}$
\item $A = [2, 4]$\quad \dotfill\; ???\;\dotfill \quad $\left( \frac{256}{3\,\sqrt[3]{224^2}} \;,\; \frac{104}{3\,\sqrt[3]{224^2}} \right)$
\item $g(x,y) = \sin\frac{x}{y} \enspace , \enspace D = [0, -3, ?]$\quad \dotfill\; ???\;\dotfill \quad $x+3z=0$
\item $\Delta x = {0.01} \enspace,\enspace \Delta y = {0.3}$\quad \dotfill\; ???\;\dotfill \quad $-0.03$
\item $h(x,y) = -x^2-6y^2-xy-3x$\quad \dotfill\; ???\;\dotfill \quad $\nicefrac{54}{23}$
\item $y = -x-4$\quad \dotfill\; ???\;\dotfill \quad $\left[\frac{-47}{12}\;,\;\frac{-1}{16}\right]$
\end{enumerate}
\end{varwidth}
\end{center}
\end{minipage}
\begin{minipage}{0.20\linewidth}
\begin{center}
{\Huge\bfseries 2.} \\[2mm]
\includegraphics[height=37mm]{images/braille.png}
{\small Braillove písmeno}
\end{center}
\end{minipage}
\end{center}
\end{minipage}
\\ \hdashline
\begin{minipage}[c][104.5mm][t]{0.5\linewidth}
\begin{center}
\vspace{7mm}
{\huge \LARGE Funkcia dvoch premenných, skupina \textit{Xi $\xi$} -\romannumeral3}\\[5mm]
\begin{minipage}{0.95\linewidth}
\begin{center}
{\small \begin{multicols}{2} \raggedright
\textbf{(a)} Vypočítajte zadanú parciálnu deriváciu funkcie \(f(x,y)\). \\[1mm]
\textbf{(b)} Určte gradient \(f(x,y)\) v bode \(A\). \\[1mm]
\textbf{(c)} Napíšte rovnicu dotykovej roviny ku grafu funkcie \(g(x,y)\) v dotykovom bode \(D\). \par \columnbreak
\textbf{(d)} Pomocou lineárnej aproximácie odhadnite funkčnú hodnotu funkcie \(g(x,y)\) v bode, ktorý je posunutý o vektor \((\Delta x, \Delta y)\) od bodu \(D_0\). \\[1mm]
\textbf{(e)} Vypočítajte \(h(x,y)\) v jej lokálnom extréme. \\[1mm]
\textbf{(f)} Nájdite bod viazeného lokálneho extrému funkcie \(h(x,y)\) na priamke. \\[1mm]
\end{multicols}
}
\end{center}
\end{minipage}
\\[4mm]
\begin{minipage}{0.79\linewidth}
\begin{center}
\begin{varwidth}{\linewidth}
\begin{enumerate}
\normalsizerr
\item $f(x,y) = \sqrt[3]{2x^2y + 2xy^2}\enspace , \enspace \pdv{f}{y}$\quad \dotfill\; ???\;\dotfill \quad $\frac{2x^2+4xy}{3\,\sqrt[3]{(2x^2y+2xy^2)^{2}}}$
\item $A = [4, 1]$\quad \dotfill\; ???\;\dotfill \quad $\left( \frac{18}{3\,\sqrt[3]{40^2}} \;,\; \frac{48}{3\,\sqrt[3]{40^2}} \right)$
\item $g(x,y) = \ln\frac{x}{y} \enspace , \enspace D=[e^{5},1,?]$\quad \dotfill\; ???\;\dotfill \quad $\frac{x}{e^5} - y - z + 5 = 0$
\item $\Delta x = {-0.2} \enspace,\enspace \Delta y = {-0.01}$\quad \dotfill\; ???\;\dotfill \quad $5.199$
\item $h(x,y) = -3x^2-y^2-2xy-2x$\quad \dotfill\; ???\;\dotfill \quad $\nicefrac{5}{8}$
\item $y = -2x-5$\quad \dotfill\; ???\;\dotfill \quad $\left[-2\;,\;-1\right]$
\end{enumerate}
\end{varwidth}
\end{center}
\end{minipage}
\begin{minipage}{0.20\linewidth}
\begin{center}
{\Huge\bfseries 3.} \\[2mm]
\includegraphics[height=37mm]{images/braille.png}
{\small Braillove písmeno}
\end{center}
\end{minipage}
\end{center}
\end{minipage}
&
\begin{minipage}[c][104.5mm][t]{0.5\linewidth}
\begin{center}
\vspace{7mm}
{\huge \LARGE Funkcia dvoch premenných, skupina \textit{Xi $\xi$} -\romannumeral4}\\[5mm]
\begin{minipage}{0.95\linewidth}
\begin{center}
{\small \begin{multicols}{2} \raggedright
\textbf{(a)} Vypočítajte zadanú parciálnu deriváciu funkcie \(f(x,y)\). \\[1mm]
\textbf{(b)} Určte gradient \(f(x,y)\) v bode \(A\). \\[1mm]
\textbf{(c)} Napíšte rovnicu dotykovej roviny ku grafu funkcie \(g(x,y)\) v dotykovom bode \(D\). \par \columnbreak
\textbf{(d)} Pomocou lineárnej aproximácie odhadnite funkčnú hodnotu funkcie \(g(x,y)\) v bode, ktorý je posunutý o vektor \((\Delta x, \Delta y)\) od bodu \(D_0\). \\[1mm]
\textbf{(e)} Vypočítajte \(h(x,y)\) v jej lokálnom extréme. \\[1mm]
\textbf{(f)} Nájdite bod viazeného lokálneho extrému funkcie \(h(x,y)\) na priamke. \\[1mm]
\end{multicols}
}
\end{center}
\end{minipage}
\\[4mm]
\begin{minipage}{0.79\linewidth}
\begin{center}
\begin{varwidth}{\linewidth}
\begin{enumerate}
\normalsizerr
\item $f(x,y) = \sqrt{5x^2y + xy^2}\enspace , \enspace \pdv{f}{y}$\quad \dotfill\; ???\;\dotfill \quad $\frac{5x^2+2xy}{2\,\sqrt{5x^2y+xy^2}}$
\item $A = [4, 4]$\quad \dotfill\; ???\;\dotfill \quad $\left( \frac{352}{2\,\sqrt{384}} \;,\; \frac{112}{2\,\sqrt{384}} \right)$
\item $g(x,y) = \sin\frac{x}{y} \enspace , \enspace D = [0, -2, ?]$\quad \dotfill\; ???\;\dotfill \quad $x-2z=0$
\item $\Delta x = {-0.2} \enspace,\enspace \Delta y = {-0.2}$\quad \dotfill\; ???\;\dotfill \quad $0.4$
\item $h(x,y) = 4x^2+2y^2+5xy+5x$\quad \dotfill\; ???\;\dotfill \quad $\nicefrac{-50}{57}$
\item $y = 3x+1$\quad \dotfill\; ???\;\dotfill \quad $\left[\frac{-11}{37}\;,\;\frac{4}{7}\right]$
\end{enumerate}
\end{varwidth}
\end{center}
\end{minipage}
\begin{minipage}{0.20\linewidth}
\begin{center}
{\Huge\bfseries 4.} \\[2mm]
\includegraphics[height=37mm]{images/braille.png}
{\small Braillove písmeno}
\end{center}
\end{minipage}
\end{center}
\end{minipage}
%
\end{tabular}
\newpage
\thispagestyle{empty}
\begin{tabular}{c:c}
\begin{minipage}[c][104.5mm][t]{0.5\linewidth}
\begin{center}
\vspace{7mm}
{\huge \LARGE Funkcia dvoch premenných, skupina \textit{Omicron $\omicron$} -\romannumeral1}\\[5mm]
\begin{minipage}{0.95\linewidth}
\begin{center}
{\small \begin{multicols}{2} \raggedright
\textbf{(a)} Vypočítajte zadanú parciálnu deriváciu funkcie \(f(x,y)\). \\[1mm]
\textbf{(b)} Určte gradient \(f(x,y)\) v bode \(A\). \\[1mm]
\textbf{(c)} Napíšte rovnicu dotykovej roviny ku grafu funkcie \(g(x,y)\) v dotykovom bode \(D\). \par \columnbreak
\textbf{(d)} Pomocou lineárnej aproximácie odhadnite funkčnú hodnotu funkcie \(g(x,y)\) v bode, ktorý je posunutý o vektor \((\Delta x, \Delta y)\) od bodu \(D_0\). \\[1mm]
\textbf{(e)} Vypočítajte \(h(x,y)\) v jej lokálnom extréme. \\[1mm]
\textbf{(f)} Nájdite bod viazeného lokálneho extrému funkcie \(h(x,y)\) na priamke. \\[1mm]
\end{multicols}
}
\end{center}
\end{minipage}
\\[4mm]
\begin{minipage}{0.79\linewidth}
\begin{center}
\begin{varwidth}{\linewidth}
\begin{enumerate}
\normalsizerr
\item $f(x,y) = \sqrt[4]{4x^2y + 2xy^2}\enspace , \enspace \pdv{f}{y}$\quad \dotfill\; ???\;\dotfill \quad $\frac{4x^2+2xy}{4\,\sqrt[4]{(4x^2y+2xy^2)^{3}}}$
\item $A = [3, 2]$\quad \dotfill\; ???\;\dotfill \quad $\left( \frac{56}{4\,\sqrt[4]{96^3}} \;,\; \frac{60}{4\,\sqrt[4]{96^3}} \right)$
\item $g(x,y) = \sin\frac{x}{y} \enspace , \enspace D = [0, 1, ?]$\quad \dotfill\; ???\;\dotfill \quad $x-z=0$
\item $\Delta x = {0.03} \enspace,\enspace \Delta y = {-0.02}$\quad \dotfill\; ???\;\dotfill \quad $0.03$
\item $h(x,y) = x^2+6y^2+2xy-3x$\quad \dotfill\; ???\;\dotfill \quad $\nicefrac{-27}{14}$
\item $y = -2x-2$\quad \dotfill\; ???\;\dotfill \quad $\left[\frac{-41}{42}\;,\;\frac{-1}{21}\right]$
\end{enumerate}
\end{varwidth}
\end{center}
\end{minipage}
\begin{minipage}{0.20\linewidth}
\begin{center}
{\Huge\bfseries 1.} \\[2mm]
\includegraphics[height=37mm]{images/braille.png}
{\small Braillove písmeno}
\end{center}
\end{minipage}
\end{center}
\end{minipage}
&
\begin{minipage}[c][104.5mm][t]{0.5\linewidth}
\begin{center}
\vspace{7mm}
{\huge \LARGE Funkcia dvoch premenných, skupina \textit{Omicron $\omicron$} -\romannumeral2}\\[5mm]
\begin{minipage}{0.95\linewidth}
\begin{center}
{\small \begin{multicols}{2} \raggedright
\textbf{(a)} Vypočítajte zadanú parciálnu deriváciu funkcie \(f(x,y)\). \\[1mm]
\textbf{(b)} Určte gradient \(f(x,y)\) v bode \(A\). \\[1mm]
\textbf{(c)} Napíšte rovnicu dotykovej roviny ku grafu funkcie \(g(x,y)\) v dotykovom bode \(D\). \par \columnbreak
\textbf{(d)} Pomocou lineárnej aproximácie odhadnite funkčnú hodnotu funkcie \(g(x,y)\) v bode, ktorý je posunutý o vektor \((\Delta x, \Delta y)\) od bodu \(D_0\). \\[1mm]
\textbf{(e)} Vypočítajte \(h(x,y)\) v jej lokálnom extréme. \\[1mm]
\textbf{(f)} Nájdite bod viazeného lokálneho extrému funkcie \(h(x,y)\) na priamke. \\[1mm]
\end{multicols}
}
\end{center}
\end{minipage}
\\[4mm]
\begin{minipage}{0.79\linewidth}
\begin{center}
\begin{varwidth}{\linewidth}
\begin{enumerate}
\normalsizerr
\item $f(x,y) = \sqrt{6x^2y + 7xy^2}\enspace , \enspace \pdv{f}{x}$\quad \dotfill\; ???\;\dotfill \quad $\frac{12xy+7y^2}{2\,\sqrt{6x^2y+7xy^2}}$
\item $A = [2, 2]$\quad \dotfill\; ???\;\dotfill \quad $\left( \frac{152}{2\,\sqrt{104}} \;,\; \frac{80}{2\,\sqrt{104}} \right)$
\item $g(x,y) = \sin\frac{x}{y} \enspace , \enspace D = [0, -1, ?]$\quad \dotfill\; ???\;\dotfill \quad $x-z=0$
\item $\Delta x = {-0.01} \enspace,\enspace \Delta y = {0.01}$\quad \dotfill\; ???\;\dotfill \quad $-0.01$
\item $h(x,y) = -6x^2-6y^2-xy-x$\quad \dotfill\; ???\;\dotfill \quad $\nicefrac{6}{145}$
\item $y = x+3$\quad \dotfill\; ???\;\dotfill \quad $\left[\frac{-20}{13}\;,\;\frac{19}{11}\right]$
\end{enumerate}
\end{varwidth}
\end{center}
\end{minipage}
\begin{minipage}{0.20\linewidth}
\begin{center}
{\Huge\bfseries 2.} \\[2mm]
\includegraphics[height=37mm]{images/braille.png}
{\small Braillove písmeno}
\end{center}
\end{minipage}
\end{center}
\end{minipage}
\\ \hdashline
\begin{minipage}[c][104.5mm][t]{0.5\linewidth}
\begin{center}
\vspace{7mm}
{\huge \LARGE Funkcia dvoch premenných, skupina \textit{Omicron $\omicron$} -\romannumeral3}\\[5mm]
\begin{minipage}{0.95\linewidth}
\begin{center}
{\small \begin{multicols}{2} \raggedright
\textbf{(a)} Vypočítajte zadanú parciálnu deriváciu funkcie \(f(x,y)\). \\[1mm]
\textbf{(b)} Určte gradient \(f(x,y)\) v bode \(A\). \\[1mm]
\textbf{(c)} Napíšte rovnicu dotykovej roviny ku grafu funkcie \(g(x,y)\) v dotykovom bode \(D\). \par \columnbreak
\textbf{(d)} Pomocou lineárnej aproximácie odhadnite funkčnú hodnotu funkcie \(g(x,y)\) v bode, ktorý je posunutý o vektor \((\Delta x, \Delta y)\) od bodu \(D_0\). \\[1mm]
\textbf{(e)} Vypočítajte \(h(x,y)\) v jej lokálnom extréme. \\[1mm]
\textbf{(f)} Nájdite bod viazeného lokálneho extrému funkcie \(h(x,y)\) na priamke. \\[1mm]
\end{multicols}
}
\end{center}
\end{minipage}
\\[4mm]
\begin{minipage}{0.79\linewidth}
\begin{center}
\begin{varwidth}{\linewidth}
\begin{enumerate}
\normalsizerr
\item $f(x,y) = \sqrt[3]{2x^2y + 4xy^2}\enspace , \enspace \pdv{f}{x}$\quad \dotfill\; ???\;\dotfill \quad $\frac{4xy+4y^2}{3\,\sqrt[3]{(2x^2y+4xy^2)^{2}}}$
\item $A = [1, 2]$\quad \dotfill\; ???\;\dotfill \quad $\left( \frac{24}{3\,\sqrt[3]{20^2}} \;,\; \frac{18}{3\,\sqrt[3]{20^2}} \right)$
\item $g(x,y) = \sin\frac{x}{y} \enspace , \enspace D = [0, -3, ?]$\quad \dotfill\; ???\;\dotfill \quad $x-3z=0$
\item $\Delta x = {-0.1} \enspace,\enspace \Delta y = {-0.02}$\quad \dotfill\; ???\;\dotfill \quad $0.007$
\item $h(x,y) = -4x^2-4y^2-xy-5x$\quad \dotfill\; ???\;\dotfill \quad $\nicefrac{101}{63}$
\item $y = -x+1$\quad \dotfill\; ???\;\dotfill \quad $\left[\frac{1}{7}\;,\;\frac{2}{3}\right]$
\end{enumerate}
\end{varwidth}
\end{center}
\end{minipage}
\begin{minipage}{0.20\linewidth}
\begin{center}
{\Huge\bfseries 3.} \\[2mm]
\includegraphics[height=37mm]{images/braille.png}
{\small Braillove písmeno}
\end{center}
\end{minipage}
\end{center}
\end{minipage}
&
\begin{minipage}[c][104.5mm][t]{0.5\linewidth}
\begin{center}
\vspace{7mm}
{\huge \LARGE Funkcia dvoch premenných, skupina \textit{Omicron $\omicron$} -\romannumeral4}\\[5mm]
\begin{minipage}{0.95\linewidth}
\begin{center}
{\small \begin{multicols}{2} \raggedright
\textbf{(a)} Vypočítajte zadanú parciálnu deriváciu funkcie \(f(x,y)\). \\[1mm]
\textbf{(b)} Určte gradient \(f(x,y)\) v bode \(A\). \\[1mm]
\textbf{(c)} Napíšte rovnicu dotykovej roviny ku grafu funkcie \(g(x,y)\) v dotykovom bode \(D\). \par \columnbreak
\textbf{(d)} Pomocou lineárnej aproximácie odhadnite funkčnú hodnotu funkcie \(g(x,y)\) v bode, ktorý je posunutý o vektor \((\Delta x, \Delta y)\) od bodu \(D_0\). \\[1mm]
\textbf{(e)} Vypočítajte \(h(x,y)\) v jej lokálnom extréme. \\[1mm]
\textbf{(f)} Nájdite bod viazeného lokálneho extrému funkcie \(h(x,y)\) na priamke. \\[1mm]
\end{multicols}
}
\end{center}
\end{minipage}
\\[4mm]
\begin{minipage}{0.79\linewidth}
\begin{center}
\begin{varwidth}{\linewidth}
\begin{enumerate}
\normalsizerr
\item $f(x,y) = \sqrt[3]{7x^2y + 7xy^2}\enspace , \enspace \pdv{f}{y}$\quad \dotfill\; ???\;\dotfill \quad $\frac{7x^2+14xy}{3\,\sqrt[3]{(7x^2y+7xy^2)^{2}}}$
\item $A = [3, 1]$\quad \dotfill\; ???\;\dotfill \quad $\left( \frac{49}{3\,\sqrt[3]{84^2}} \;,\; \frac{210}{3\,\sqrt[3]{84^2}} \right)$
\item $g(x,y) = \ln\frac{x}{y} \enspace , \enspace D=[e^{4},1,?]$\quad \dotfill\; ???\;\dotfill \quad $\frac{x}{e^4} + y + z + 4 = 0$
\item $\Delta x = {-0.02} \enspace,\enspace \Delta y = {0.01}$\quad \dotfill\; ???\;\dotfill \quad $4.01$
\item $h(x,y) = -2x^2-y^2-2xy-x$\quad \dotfill\; ???\;\dotfill \quad $\nicefrac{1}{12}$
\item $y = 5x+3$\quad \dotfill\; ???\;\dotfill \quad $\left[\frac{23}{74}\;,\;\frac{1}{2}\right]$
\end{enumerate}
\end{varwidth}
\end{center}
\end{minipage}
\begin{minipage}{0.20\linewidth}
\begin{center}
{\Huge\bfseries 4.} \\[2mm]
\includegraphics[height=37mm]{images/braille.png}
{\small Braillove písmeno}
\end{center}
\end{minipage}
\end{center}
\end{minipage}
%
\end{tabular}
\newpage
\thispagestyle{empty}
\begin{tabular}{c:c}
\begin{minipage}[c][104.5mm][t]{0.5\linewidth}
\begin{center}
\vspace{7mm}
{\huge \LARGE Funkcia dvoch premenných, skupina \textit{Pi $\pi$} -\romannumeral1}\\[5mm]
\begin{minipage}{0.95\linewidth}
\begin{center}
{\small \begin{multicols}{2} \raggedright
\textbf{(a)} Vypočítajte zadanú parciálnu deriváciu funkcie \(f(x,y)\). \\[1mm]
\textbf{(b)} Určte gradient \(f(x,y)\) v bode \(A\). \\[1mm]
\textbf{(c)} Napíšte rovnicu dotykovej roviny ku grafu funkcie \(g(x,y)\) v dotykovom bode \(D\). \par \columnbreak
\textbf{(d)} Pomocou lineárnej aproximácie odhadnite funkčnú hodnotu funkcie \(g(x,y)\) v bode, ktorý je posunutý o vektor \((\Delta x, \Delta y)\) od bodu \(D_0\). \\[1mm]
\textbf{(e)} Vypočítajte \(h(x,y)\) v jej lokálnom extréme. \\[1mm]
\textbf{(f)} Nájdite bod viazeného lokálneho extrému funkcie \(h(x,y)\) na priamke. \\[1mm]
\end{multicols}
}
\end{center}
\end{minipage}
\\[4mm]
\begin{minipage}{0.79\linewidth}
\begin{center}
\begin{varwidth}{\linewidth}
\begin{enumerate}
\normalsizerr
\item $f(x,y) = \sqrt{9x^2y + 8xy^2}\enspace , \enspace \pdv{f}{x}$\quad \dotfill\; ???\;\dotfill \quad $\frac{18xy+8y^2}{2\,\sqrt{9x^2y+8xy^2}}$
\item $A = [2, 1]$\quad \dotfill\; ???\;\dotfill \quad $\left( \frac{44}{2\,\sqrt{52}} \;,\; \frac{68}{2\,\sqrt{52}} \right)$
\item $g(x,y) = \ln\frac{x}{y} \enspace , \enspace D=[1,1,?]$\quad \dotfill\; ???\;\dotfill \quad $x - y - z + 0 = 0$
\item $\Delta x = {-0.1} \enspace,\enspace \Delta y = {0.3}$\quad \dotfill\; ???\;\dotfill \quad $-0.4$
\item $h(x,y) = -6x^2-3y^2-6xy+3x$\quad \dotfill\; ???\;\dotfill \quad $\nicefrac{3}{4}$
\item $y = -x-3$\quad \dotfill\; ???\;\dotfill \quad $\left[\frac{13}{2}\;,\;\frac{-7}{2}\right]$
\end{enumerate}
\end{varwidth}
\end{center}
\end{minipage}
\begin{minipage}{0.20\linewidth}
\begin{center}
{\Huge\bfseries 1.} \\[2mm]
\includegraphics[height=37mm]{images/braille.png}
{\small Braillove písmeno}
\end{center}
\end{minipage}
\end{center}
\end{minipage}
&
\begin{minipage}[c][104.5mm][t]{0.5\linewidth}
\begin{center}
\vspace{7mm}
{\huge \LARGE Funkcia dvoch premenných, skupina \textit{Pi $\pi$} -\romannumeral2}\\[5mm]
\begin{minipage}{0.95\linewidth}
\begin{center}
{\small \begin{multicols}{2} \raggedright
\textbf{(a)} Vypočítajte zadanú parciálnu deriváciu funkcie \(f(x,y)\). \\[1mm]
\textbf{(b)} Určte gradient \(f(x,y)\) v bode \(A\). \\[1mm]
\textbf{(c)} Napíšte rovnicu dotykovej roviny ku grafu funkcie \(g(x,y)\) v dotykovom bode \(D\). \par \columnbreak
\textbf{(d)} Pomocou lineárnej aproximácie odhadnite funkčnú hodnotu funkcie \(g(x,y)\) v bode, ktorý je posunutý o vektor \((\Delta x, \Delta y)\) od bodu \(D_0\). \\[1mm]
\textbf{(e)} Vypočítajte \(h(x,y)\) v jej lokálnom extréme. \\[1mm]
\textbf{(f)} Nájdite bod viazeného lokálneho extrému funkcie \(h(x,y)\) na priamke. \\[1mm]
\end{multicols}
}
\end{center}
\end{minipage}
\\[4mm]
\begin{minipage}{0.79\linewidth}
\begin{center}
\begin{varwidth}{\linewidth}
\begin{enumerate}
\normalsizerr
\item $f(x,y) = \sqrt{2x^2y + xy^2}\enspace , \enspace \pdv{f}{x}$\quad \dotfill\; ???\;\dotfill \quad $\frac{4xy+y^2}{\sqrt{2x^2y+2xy^2}}$
\item $A = [4, 4]$\quad \dotfill\; ???\;\dotfill \quad $\left( \frac{80}{2\,\sqrt{192}} \;,\; \frac{64}{2\,\sqrt{192}} \right)$
\item $g(x,y) = \sin\frac{x}{y} \enspace , \enspace D = [0, -3, ?]$\quad \dotfill\; ???\;\dotfill \quad $x-3z=0$
\item $\Delta x = {0.02} \enspace,\enspace \Delta y = {0.2}$\quad \dotfill\; ???\;\dotfill \quad $-0.067$
\item $h(x,y) = 3x^2+3y^2-3xy-3x$\quad \dotfill\; ???\;\dotfill \quad $\nicefrac{-26}{27}$
\item $y = 3x-1$\quad \dotfill\; ???\;\dotfill \quad $\left[\frac{3}{7}\;,\;\frac{2}{7}\right]$
\end{enumerate}
\end{varwidth}
\end{center}
\end{minipage}
\begin{minipage}{0.20\linewidth}
\begin{center}
{\Huge\bfseries 2.} \\[2mm]
\includegraphics[height=37mm]{images/braille.png}
{\small Braillove písmeno}
\end{center}
\end{minipage}
\end{center}
\end{minipage}
\\ \hdashline
\begin{minipage}[c][104.5mm][t]{0.5\linewidth}
\begin{center}
\vspace{7mm}
{\huge \LARGE Funkcia dvoch premenných, skupina \textit{Pi $\pi$} -\romannumeral3}\\[5mm]
\begin{minipage}{0.95\linewidth}
\begin{center}
{\small \begin{multicols}{2} \raggedright
\textbf{(a)} Vypočítajte zadanú parciálnu deriváciu funkcie \(f(x,y)\). \\[1mm]
\textbf{(b)} Určte gradient \(f(x,y)\) v bode \(A\). \\[1mm]
\textbf{(c)} Napíšte rovnicu dotykovej roviny ku grafu funkcie \(g(x,y)\) v dotykovom bode \(D\). \par \columnbreak
\textbf{(d)} Pomocou lineárnej aproximácie odhadnite funkčnú hodnotu funkcie \(g(x,y)\) v bode, ktorý je posunutý o vektor \((\Delta x, \Delta y)\) od bodu \(D_0\). \\[1mm]
\textbf{(e)} Vypočítajte \(h(x,y)\) v jej lokálnom extréme. \\[1mm]
\textbf{(f)} Nájdite bod viazeného lokálneho extrému funkcie \(h(x,y)\) na priamke. \\[1mm]
\end{multicols}
}
\end{center}
\end{minipage}
\\[4mm]
\begin{minipage}{0.79\linewidth}
\begin{center}
\begin{varwidth}{\linewidth}
\begin{enumerate}
\normalsizerr
\item $f(x,y) = \sqrt[4]{x^2y + 2xy^2}\enspace , \enspace \pdv{f}{x}$\quad \dotfill\; ???\;\dotfill \quad $\frac{2xy+2y^2}{4\,\sqrt[4]{(x^2y+2xy^2)^{3}}}$
\item $A = [1, 1]$\quad \dotfill\; ???\;\dotfill \quad $\left( \frac{4}{4\,\sqrt[4]{3^3}} \;,\; \frac{5}{4\,\sqrt[4]{3^3}} \right)$
\item $g(x,y) = \sin\frac{x}{y} \enspace , \enspace D = [0, -3, ?]$\quad \dotfill\; ???\;\dotfill \quad $x+3z=1$
\item $\Delta x = {-0.1} \enspace,\enspace \Delta y = {-0.2}$\quad \dotfill\; ???\;\dotfill \quad $0.067$
\item $h(x,y) = 3x^2+y^2-3xy-x$\quad \dotfill\; ???\;\dotfill \quad $\nicefrac{-1}{21}$
\item $y = -3x+7$\quad \dotfill\; ???\;\dotfill \quad $\left[\frac{32}{21}\;,\;\frac{17}{7}\right]$
\end{enumerate}
\end{varwidth}
\end{center}
\end{minipage}
\begin{minipage}{0.20\linewidth}
\begin{center}
{\Huge\bfseries 3.} \\[2mm]
\includegraphics[height=37mm]{images/braille.png}
{\small Braillove písmeno}
\end{center}
\end{minipage}
\end{center}
\end{minipage}
&
\begin{minipage}[c][104.5mm][t]{0.5\linewidth}
\begin{center}
\vspace{7mm}
{\huge \LARGE Funkcia dvoch premenných, skupina \textit{Pi $\pi$} -\romannumeral4}\\[5mm]
\begin{minipage}{0.95\linewidth}
\begin{center}
{\small \begin{multicols}{2} \raggedright
\textbf{(a)} Vypočítajte zadanú parciálnu deriváciu funkcie \(f(x,y)\). \\[1mm]
\textbf{(b)} Určte gradient \(f(x,y)\) v bode \(A\). \\[1mm]
\textbf{(c)} Napíšte rovnicu dotykovej roviny ku grafu funkcie \(g(x,y)\) v dotykovom bode \(D\). \par \columnbreak
\textbf{(d)} Pomocou lineárnej aproximácie odhadnite funkčnú hodnotu funkcie \(g(x,y)\) v bode, ktorý je posunutý o vektor \((\Delta x, \Delta y)\) od bodu \(D_0\). \\[1mm]
\textbf{(e)} Vypočítajte \(h(x,y)\) v jej lokálnom extréme. \\[1mm]
\textbf{(f)} Nájdite bod viazeného lokálneho extrému funkcie \(h(x,y)\) na priamke. \\[1mm]
\end{multicols}
}
\end{center}
\end{minipage}
\\[4mm]
\begin{minipage}{0.79\linewidth}
\begin{center}
\begin{varwidth}{\linewidth}
\begin{enumerate}
\normalsizerr
\item $f(x,y) = \sqrt[4]{8x^2y + xy^2}\enspace , \enspace \pdv{f}{x}$\quad \dotfill\; ???\;\dotfill \quad $\frac{16xy+2y^2}{\sqrt[4]{(8x^2y+xy^2)^{3}}}$
\item $A = [3, 1]$\quad \dotfill\; ???\;\dotfill \quad $\left( \frac{49}{4\,\sqrt[4]{75^3}} \;,\; \frac{78}{4\,\sqrt[4]{75^3}} \right)$
\item $g(x,y) = \ln\frac{x}{y} \enspace , \enspace D=[1,1,?]$\quad \dotfill\; ???\;\dotfill \quad $x + y + z + 0 = 0$
\item $\Delta x = {0.2} \enspace,\enspace \Delta y = {-0.01}$\quad \dotfill\; ???\;\dotfill \quad $0.0$
\item $h(x,y) = -x^2-5y^2-xy+2x$\quad \dotfill\; ???\;\dotfill \quad $\nicefrac{20}{21}$
\item $y = x-2$\quad \dotfill\; ???\;\dotfill \quad $\left[\frac{12}{7}\;,\;\frac{-2}{7}\right]$
\end{enumerate}
\end{varwidth}
\end{center}
\end{minipage}
\begin{minipage}{0.20\linewidth}
\begin{center}
{\Huge\bfseries 4.} \\[2mm]
\includegraphics[height=37mm]{images/braille.png}
{\small Braillove písmeno}
\end{center}
\end{minipage}
\end{center}
\end{minipage}
%
\end{tabular}
\newpage
\thispagestyle{empty}
\begin{tabular}{c:c}
\begin{minipage}[c][104.5mm][t]{0.5\linewidth}
\begin{center}
\vspace{7mm}
{\huge \LARGE Funkcia dvoch premenných, skupina \textit{Rho $\rho$} -\romannumeral1}\\[5mm]
\begin{minipage}{0.95\linewidth}
\begin{center}
{\small \begin{multicols}{2} \raggedright
\textbf{(a)} Vypočítajte zadanú parciálnu deriváciu funkcie \(f(x,y)\). \\[1mm]
\textbf{(b)} Určte gradient \(f(x,y)\) v bode \(A\). \\[1mm]
\textbf{(c)} Napíšte rovnicu dotykovej roviny ku grafu funkcie \(g(x,y)\) v dotykovom bode \(D\). \par \columnbreak
\textbf{(d)} Pomocou lineárnej aproximácie odhadnite funkčnú hodnotu funkcie \(g(x,y)\) v bode, ktorý je posunutý o vektor \((\Delta x, \Delta y)\) od bodu \(D_0\). \\[1mm]
\textbf{(e)} Vypočítajte \(h(x,y)\) v jej lokálnom extréme. \\[1mm]
\textbf{(f)} Nájdite bod viazeného lokálneho extrému funkcie \(h(x,y)\) na priamke. \\[1mm]
\end{multicols}
}
\end{center}
\end{minipage}
\\[4mm]
\begin{minipage}{0.79\linewidth}
\begin{center}
\begin{varwidth}{\linewidth}
\begin{enumerate}
\normalsizerr
\item $f(x,y) = \sqrt[3]{8x^2y + xy^2}\enspace , \enspace \pdv{f}{x}$\quad \dotfill\; ???\;\dotfill \quad $\frac{16xy+y^2}{3\,\sqrt[3]{(8x^2y+xy^2)^{2}}}$
\item $A = [1, 2]$\quad \dotfill\; ???\;\dotfill \quad $\left( \frac{36}{3\,\sqrt[3]{20^2}} \;,\; \frac{12}{3\,\sqrt[3]{20^2}} \right)$
\item $g(x,y) = \ln\frac{x}{y} \enspace , \enspace D=[e^{4},1,?]$\quad \dotfill\; ???\;\dotfill \quad $\frac{x}{e^5} - y - z + 5 = 0$
\item $\Delta x = {-0.1} \enspace,\enspace \Delta y = {-0.1}$\quad \dotfill\; ???\;\dotfill \quad $3.898$
\item $h(x,y) = -x^2-3y^2-xy+4x$\quad \dotfill\; ???\;\dotfill \quad $\nicefrac{48}{11}$
\item $y = x+3$\quad \dotfill\; ???\;\dotfill \quad $\left[\frac{19}{10}\;,\;\frac{13}{10}\right]$
\end{enumerate}
\end{varwidth}
\end{center}
\end{minipage}
\begin{minipage}{0.20\linewidth}
\begin{center}
{\Huge\bfseries 1.} \\[2mm]
\includegraphics[height=37mm]{images/braille.png}
{\small Braillove písmeno}
\end{center}
\end{minipage}
\end{center}
\end{minipage}
&
\begin{minipage}[c][104.5mm][t]{0.5\linewidth}
\begin{center}
\vspace{7mm}
{\huge \LARGE Funkcia dvoch premenných, skupina \textit{Rho $\rho$} -\romannumeral2}\\[5mm]
\begin{minipage}{0.95\linewidth}
\begin{center}
{\small \begin{multicols}{2} \raggedright
\textbf{(a)} Vypočítajte zadanú parciálnu deriváciu funkcie \(f(x,y)\). \\[1mm]
\textbf{(b)} Určte gradient \(f(x,y)\) v bode \(A\). \\[1mm]
\textbf{(c)} Napíšte rovnicu dotykovej roviny ku grafu funkcie \(g(x,y)\) v dotykovom bode \(D\). \par \columnbreak
\textbf{(d)} Pomocou lineárnej aproximácie odhadnite funkčnú hodnotu funkcie \(g(x,y)\) v bode, ktorý je posunutý o vektor \((\Delta x, \Delta y)\) od bodu \(D_0\). \\[1mm]
\textbf{(e)} Vypočítajte \(h(x,y)\) v jej lokálnom extréme. \\[1mm]
\textbf{(f)} Nájdite bod viazeného lokálneho extrému funkcie \(h(x,y)\) na priamke. \\[1mm]
\end{multicols}
}
\end{center}
\end{minipage}
\\[4mm]
\begin{minipage}{0.79\linewidth}
\begin{center}
\begin{varwidth}{\linewidth}
\begin{enumerate}
\normalsizerr
\item $f(x,y) = \sqrt[3]{2x^2y + 3xy^2}\enspace , \enspace \pdv{f}{y}$\quad \dotfill\; ???\;\dotfill \quad $\frac{2x^2+6xy}{3\,\sqrt[3]{(2x^2y+3xy^2)^{2}}}$
\item $A = [2, 3]$\quad \dotfill\; ???\;\dotfill \quad $\left( \frac{51}{3\,\sqrt[3]{78^2}} \;,\; \frac{44}{3\,\sqrt[3]{78^2}} \right)$
\item $g(x,y) = \sin\frac{x}{y} \enspace , \enspace D = [0, -2, ?]$\quad \dotfill\; ???\;\dotfill \quad $x+2z=0$
\item $\Delta x = {-0.01} \enspace,\enspace \Delta y = {0.01}$\quad \dotfill\; ???\;\dotfill \quad $-0.005$
\item $h(x,y) = -6x^2-y^2-xy+x$\quad \dotfill\; ???\;\dotfill \quad $\nicefrac{1}{23}$
\item $y = x+7$\quad \dotfill\; ???\;\dotfill \quad $\left[\frac{-5}{4}\;,\;\frac{23}{3}\right]$
\end{enumerate}
\end{varwidth}
\end{center}
\end{minipage}
\begin{minipage}{0.20\linewidth}
\begin{center}
{\Huge\bfseries 2.} \\[2mm]
\includegraphics[height=37mm]{images/braille.png}
{\small Braillove písmeno}
\end{center}
\end{minipage}
\end{center}
\end{minipage}
\\ \hdashline
\begin{minipage}[c][104.5mm][t]{0.5\linewidth}
\begin{center}
\vspace{7mm}
{\huge \LARGE Funkcia dvoch premenných, skupina \textit{Rho $\rho$} -\romannumeral3}\\[5mm]
\begin{minipage}{0.95\linewidth}
\begin{center}
{\small \begin{multicols}{2} \raggedright
\textbf{(a)} Vypočítajte zadanú parciálnu deriváciu funkcie \(f(x,y)\). \\[1mm]
\textbf{(b)} Určte gradient \(f(x,y)\) v bode \(A\). \\[1mm]
\textbf{(c)} Napíšte rovnicu dotykovej roviny ku grafu funkcie \(g(x,y)\) v dotykovom bode \(D\). \par \columnbreak
\textbf{(d)} Pomocou lineárnej aproximácie odhadnite funkčnú hodnotu funkcie \(g(x,y)\) v bode, ktorý je posunutý o vektor \((\Delta x, \Delta y)\) od bodu \(D_0\). \\[1mm]
\textbf{(e)} Vypočítajte \(h(x,y)\) v jej lokálnom extréme. \\[1mm]
\textbf{(f)} Nájdite bod viazeného lokálneho extrému funkcie \(h(x,y)\) na priamke. \\[1mm]
\end{multicols}
}
\end{center}
\end{minipage}
\\[4mm]
\begin{minipage}{0.79\linewidth}
\begin{center}
\begin{varwidth}{\linewidth}
\begin{enumerate}
\normalsizerr
\item $f(x,y) = \sqrt[4]{x^2y + 2xy^2}\enspace , \enspace \pdv{f}{y}$\quad \dotfill\; ???\;\dotfill \quad $\frac{x^2+4xy}{4\,\sqrt[4]{(x^2y+2xy^2)^{3}}}$
\item $A = [1, 3]$\quad \dotfill\; ???\;\dotfill \quad $\left( \frac{48}{4\,\sqrt[4]{21^3}} \;,\; \frac{13}{4\,\sqrt[4]{21^3}} \right)$
\item $g(x,y) = \ln\frac{x}{y} \enspace , \enspace D=[e,1,?]$\quad \dotfill\; ???\;\dotfill \quad $\frac{x}{e} - y - z + 2 = 0$
\item $\Delta x = {0.2} \enspace,\enspace \Delta y = {-0.01}$\quad \dotfill\; ???\;\dotfill \quad $1.064$
\item $h(x,y) = -6x^2-2y^2-4xy+5x$\quad \dotfill\; ???\;\dotfill \quad $\nicefrac{25}{32}$
\item $y = -x+2$\quad \dotfill\; ???\;\dotfill \quad $\left[\frac{5}{8}\;,\;\frac{11}{24}\right]$
\end{enumerate}
\end{varwidth}
\end{center}
\end{minipage}
\begin{minipage}{0.20\linewidth}
\begin{center}
{\Huge\bfseries 3.} \\[2mm]
\includegraphics[height=37mm]{images/braille.png}
{\small Braillove písmeno}
\end{center}
\end{minipage}
\end{center}
\end{minipage}
&
\begin{minipage}[c][104.5mm][t]{0.5\linewidth}
\begin{center}
\vspace{7mm}
{\huge \LARGE Funkcia dvoch premenných, skupina \textit{Rho $\rho$} -\romannumeral4}\\[5mm]
\begin{minipage}{0.95\linewidth}
\begin{center}
{\small \begin{multicols}{2} \raggedright
\textbf{(a)} Vypočítajte zadanú parciálnu deriváciu funkcie \(f(x,y)\). \\[1mm]
\textbf{(b)} Určte gradient \(f(x,y)\) v bode \(A\). \\[1mm]
\textbf{(c)} Napíšte rovnicu dotykovej roviny ku grafu funkcie \(g(x,y)\) v dotykovom bode \(D\). \par \columnbreak
\textbf{(d)} Pomocou lineárnej aproximácie odhadnite funkčnú hodnotu funkcie \(g(x,y)\) v bode, ktorý je posunutý o vektor \((\Delta x, \Delta y)\) od bodu \(D_0\). \\[1mm]
\textbf{(e)} Vypočítajte \(h(x,y)\) v jej lokálnom extréme. \\[1mm]
\textbf{(f)} Nájdite bod viazeného lokálneho extrému funkcie \(h(x,y)\) na priamke. \\[1mm]
\end{multicols}
}
\end{center}
\end{minipage}
\\[4mm]
\begin{minipage}{0.79\linewidth}
\begin{center}
\begin{varwidth}{\linewidth}
\begin{enumerate}
\normalsizerr
\item $f(x,y) = \sqrt[3]{3x^2y + 2xy^2}\enspace , \enspace \pdv{f}{x}$\quad \dotfill\; ???\;\dotfill \quad $\frac{6xy+2y^2}{3\,\sqrt[3]{(3x^2y+2xy^2)^{2}}}$
\item $A = [2, 2]$\quad \dotfill\; ???\;\dotfill \quad $\left( \frac{32}{3\,\sqrt[3]{40^2}} \;,\; \frac{56}{3\,\sqrt[3]{40^2}} \right)$
\item $g(x,y) = \ln\frac{x}{y} \enspace , \enspace D=[1,1,?]$\quad \dotfill\; ???\;\dotfill \quad $x + y + z + 0 = 0$
\item $\Delta x = {0.02} \enspace,\enspace \Delta y = {0.3}$\quad \dotfill\; ???\;\dotfill \quad $0.0$
\item $h(x,y) = 2x^2+4y^2+xy-x$\quad \dotfill\; ???\;\dotfill \quad $\nicefrac{-4}{31}$
\item $y = -2x+1$\quad \dotfill\; ???\;\dotfill \quad $\left[\frac{1}{2}\;,\;0\right]$
\end{enumerate}
\end{varwidth}
\end{center}
\end{minipage}
\begin{minipage}{0.20\linewidth}
\begin{center}
{\Huge\bfseries 4.} \\[2mm]
\includegraphics[height=37mm]{images/braille.png}
{\small Braillove písmeno}
\end{center}
\end{minipage}
\end{center}
\end{minipage}
%
\end{tabular}
\newpage
\thispagestyle{empty}
\begin{tabular}{c:c}
\begin{minipage}[c][104.5mm][t]{0.5\linewidth}
\begin{center}
\vspace{7mm}
{\huge \LARGE Funkcia dvoch premenných, skupina \textit{Sigma $\sigma$} -\romannumeral1}\\[5mm]
\begin{minipage}{0.95\linewidth}
\begin{center}
{\small \begin{multicols}{2} \raggedright
\textbf{(a)} Vypočítajte zadanú parciálnu deriváciu funkcie \(f(x,y)\). \\[1mm]
\textbf{(b)} Určte gradient \(f(x,y)\) v bode \(A\). \\[1mm]
\textbf{(c)} Napíšte rovnicu dotykovej roviny ku grafu funkcie \(g(x,y)\) v dotykovom bode \(D\). \par \columnbreak
\textbf{(d)} Pomocou lineárnej aproximácie odhadnite funkčnú hodnotu funkcie \(g(x,y)\) v bode, ktorý je posunutý o vektor \((\Delta x, \Delta y)\) od bodu \(D_0\). \\[1mm]
\textbf{(e)} Vypočítajte \(h(x,y)\) v jej lokálnom extréme. \\[1mm]
\textbf{(f)} Nájdite bod viazeného lokálneho extrému funkcie \(h(x,y)\) na priamke. \\[1mm]
\end{multicols}
}
\end{center}
\end{minipage}
\\[4mm]
\begin{minipage}{0.79\linewidth}
\begin{center}
\begin{varwidth}{\linewidth}
\begin{enumerate}
\normalsizerr
\item $f(x,y) = \sqrt{7x^2y + 2xy^2}\enspace , \enspace \pdv{f}{x}$\quad \dotfill\; ???\;\dotfill \quad $\frac{14xy+2y^2}{2\,\sqrt{7x^2y+2xy^2}}$
\item $A = [2, 4]$\quad \dotfill\; ???\;\dotfill \quad $\left( \frac{144}{2\,\sqrt{176}} \;,\; \frac{60}{2\,\sqrt{176}} \right)$
\item $g(x,y) = \ln\frac{x}{y} \enspace , \enspace D=[e,1,?]$\quad \dotfill\; ???\;\dotfill \quad $\frac{x}{e} - y - z + 1 = 0$
\item $\Delta x = {-0.01} \enspace,\enspace \Delta y = {0.03}$\quad \dotfill\; ???\;\dotfill \quad $1.006$
\item $h(x,y) = -7x^2-y^2+2xy-4x$\quad \dotfill\; ???\;\dotfill \quad $\nicefrac{17}{24}$
\item $y = 4x-3$\quad \dotfill\; ???\;\dotfill \quad $\left[\frac{-17}{15}\;,\;\frac{-17}{15}\right]$
\end{enumerate}
\end{varwidth}
\end{center}
\end{minipage}
\begin{minipage}{0.20\linewidth}
\begin{center}
{\Huge\bfseries 1.} \\[2mm]
\includegraphics[height=37mm]{images/braille.png}
{\small Braillove písmeno}
\end{center}
\end{minipage}
\end{center}
\end{minipage}
&
\begin{minipage}[c][104.5mm][t]{0.5\linewidth}
\begin{center}
\vspace{7mm}
{\huge \LARGE Funkcia dvoch premenných, skupina \textit{Sigma $\sigma$} -\romannumeral2}\\[5mm]
\begin{minipage}{0.95\linewidth}
\begin{center}
{\small \begin{multicols}{2} \raggedright
\textbf{(a)} Vypočítajte zadanú parciálnu deriváciu funkcie \(f(x,y)\). \\[1mm]
\textbf{(b)} Určte gradient \(f(x,y)\) v bode \(A\). \\[1mm]
\textbf{(c)} Napíšte rovnicu dotykovej roviny ku grafu funkcie \(g(x,y)\) v dotykovom bode \(D\). \par \columnbreak
\textbf{(d)} Pomocou lineárnej aproximácie odhadnite funkčnú hodnotu funkcie \(g(x,y)\) v bode, ktorý je posunutý o vektor \((\Delta x, \Delta y)\) od bodu \(D_0\). \\[1mm]
\textbf{(e)} Vypočítajte \(h(x,y)\) v jej lokálnom extréme. \\[1mm]
\textbf{(f)} Nájdite bod viazeného lokálneho extrému funkcie \(h(x,y)\) na priamke. \\[1mm]
\end{multicols}
}
\end{center}
\end{minipage}
\\[4mm]
\begin{minipage}{0.79\linewidth}
\begin{center}
\begin{varwidth}{\linewidth}
\begin{enumerate}
\normalsizerr
\item $f(x,y) = \sqrt{x^2y + 8xy^2}\enspace , \enspace \pdv{f}{x}$\quad \dotfill\; ???\;\dotfill \quad $\frac{2xy+8y^2}{2\,\sqrt{x^2y+8xy^2}}$
\item $A = [2, 3]$\quad \dotfill\; ???\;\dotfill \quad $\left( \frac{168}{2\,\sqrt{156}} \;,\; \frac{100}{2\,\sqrt{156}} \right)$
\item $g(x,y) = \ln\frac{x}{y} \enspace , \enspace D=[e^{2},1,?]$\quad \dotfill\; ???\;\dotfill \quad $\frac{x}{e^3} - y - z + 3 = 0$
\item $\Delta x = {-0.02} \enspace,\enspace \Delta y = {-0.1}$\quad \dotfill\; ???\;\dotfill \quad $2.017$
\item $h(x,y) = -3x^2-2y^2-2xy-6x$\quad \dotfill\; ???\;\dotfill \quad $\nicefrac{73}{20}$
\item $y = -x+4$\quad \dotfill\; ???\;\dotfill \quad $\left[\frac{1}{3}\;,\;\frac{11}{7}\right]$
\end{enumerate}
\end{varwidth}
\end{center}
\end{minipage}
\begin{minipage}{0.20\linewidth}
\begin{center}
{\Huge\bfseries 2.} \\[2mm]
\includegraphics[height=37mm]{images/braille.png}
{\small Braillove písmeno}
\end{center}
\end{minipage}
\end{center}
\end{minipage}
\\ \hdashline
\begin{minipage}[c][104.5mm][t]{0.5\linewidth}
\begin{center}
\vspace{7mm}
{\huge \LARGE Funkcia dvoch premenných, skupina \textit{Sigma $\sigma$} -\romannumeral3}\\[5mm]
\begin{minipage}{0.95\linewidth}
\begin{center}
{\small \begin{multicols}{2} \raggedright
\textbf{(a)} Vypočítajte zadanú parciálnu deriváciu funkcie \(f(x,y)\). \\[1mm]
\textbf{(b)} Určte gradient \(f(x,y)\) v bode \(A\). \\[1mm]
\textbf{(c)} Napíšte rovnicu dotykovej roviny ku grafu funkcie \(g(x,y)\) v dotykovom bode \(D\). \par \columnbreak
\textbf{(d)} Pomocou lineárnej aproximácie odhadnite funkčnú hodnotu funkcie \(g(x,y)\) v bode, ktorý je posunutý o vektor \((\Delta x, \Delta y)\) od bodu \(D_0\). \\[1mm]
\textbf{(e)} Vypočítajte \(h(x,y)\) v jej lokálnom extréme. \\[1mm]
\textbf{(f)} Nájdite bod viazeného lokálneho extrému funkcie \(h(x,y)\) na priamke. \\[1mm]
\end{multicols}
}
\end{center}
\end{minipage}
\\[4mm]
\begin{minipage}{0.79\linewidth}
\begin{center}
\begin{varwidth}{\linewidth}
\begin{enumerate}
\normalsizerr
\item $f(x,y) = \sqrt{3x^2y + 5xy^2}\enspace , \enspace \pdv{f}{x}$\quad \dotfill\; ???\;\dotfill \quad $\frac{6xy+5y^2}{2\,\sqrt{3x^2y+5xy^2}}$
\item $A = [1, 1]$\quad \dotfill\; ???\;\dotfill \quad $\left( \frac{22}{2\,\sqrt{8}} \;,\; \frac{13}{2\,\sqrt{8}} \right)$
\item $g(x,y) = \sin\frac{x}{y} \enspace , \enspace D = [0, -2, ?]$\quad \dotfill\; ???\;\dotfill \quad $x+2z=0$
\item $\Delta x = {-0.2} \enspace,\enspace \Delta y = {-0.3}$\quad \dotfill\; ???\;\dotfill \quad $0.15$
\item $h(x,y) = 7x^2+y^2+5xy+2x$\quad \dotfill\; ???\;\dotfill \quad $\nicefrac{-4}{53}$
\item $y = 4x+1$\quad \dotfill\; ???\;\dotfill \quad $\left[\frac{-15}{86}\;,\;\frac{13}{43}\right]$
\end{enumerate}
\end{varwidth}
\end{center}
\end{minipage}
\begin{minipage}{0.20\linewidth}
\begin{center}
{\Huge\bfseries 3.} \\[2mm]
\includegraphics[height=37mm]{images/braille.png}
{\small Braillove písmeno}
\end{center}
\end{minipage}
\end{center}
\end{minipage}
&
\begin{minipage}[c][104.5mm][t]{0.5\linewidth}
\begin{center}
\vspace{7mm}
{\huge \LARGE Funkcia dvoch premenných, skupina \textit{Sigma $\sigma$} -\romannumeral4}\\[5mm]
\begin{minipage}{0.95\linewidth}
\begin{center}
{\small \begin{multicols}{2} \raggedright
\textbf{(a)} Vypočítajte zadanú parciálnu deriváciu funkcie \(f(x,y)\). \\[1mm]
\textbf{(b)} Určte gradient \(f(x,y)\) v bode \(A\). \\[1mm]
\textbf{(c)} Napíšte rovnicu dotykovej roviny ku grafu funkcie \(g(x,y)\) v dotykovom bode \(D\). \par \columnbreak
\textbf{(d)} Pomocou lineárnej aproximácie odhadnite funkčnú hodnotu funkcie \(g(x,y)\) v bode, ktorý je posunutý o vektor \((\Delta x, \Delta y)\) od bodu \(D_0\). \\[1mm]
\textbf{(e)} Vypočítajte \(h(x,y)\) v jej lokálnom extréme. \\[1mm]
\textbf{(f)} Nájdite bod viazeného lokálneho extrému funkcie \(h(x,y)\) na priamke. \\[1mm]
\end{multicols}
}
\end{center}
\end{minipage}
\\[4mm]
\begin{minipage}{0.79\linewidth}
\begin{center}
\begin{varwidth}{\linewidth}
\begin{enumerate}
\normalsizerr
\item $f(x,y) = \sqrt{8x^2y + 2xy^2}\enspace , \enspace \pdv{f}{y}$\quad \dotfill\; ???\;\dotfill \quad $\frac{8x^2+4xy}{2\,\sqrt{8x^2y+2xy^2}}$
\item $A = [2, 3]$\quad \dotfill\; ???\;\dotfill \quad $\left( \frac{114}{2\,\sqrt{132}} \;,\; \frac{112}{2\,\sqrt{132}} \right)$
\item $g(x,y) = \sin\frac{x}{y} \enspace , \enspace D = [0, -1, ?]$\quad \dotfill\; ???\;\dotfill \quad $x+z=1$
\item $\Delta x = {-0.3} \enspace,\enspace \Delta y = {-0.01}$\quad \dotfill\; ???\;\dotfill \quad $0.01$
\item $h(x,y) = -2x^2-2y^2-3xy+2x$\quad \dotfill\; ???\;\dotfill \quad $\nicefrac{9}{7}$
\item $y = 4x-2$\quad \dotfill\; ???\;\dotfill \quad $\left[\frac{10}{23}\;,\;\frac{-6}{11}\right]$
\end{enumerate}
\end{varwidth}
\end{center}
\end{minipage}
\begin{minipage}{0.20\linewidth}
\begin{center}
{\Huge\bfseries 4.} \\[2mm]
\includegraphics[height=37mm]{images/braille.png}
{\small Braillove písmeno}
\end{center}
\end{minipage}
\end{center}
\end{minipage}
%
\end{tabular}
\newpage
\thispagestyle{empty}
\begin{tabular}{c:c}
\begin{minipage}[c][104.5mm][t]{0.5\linewidth}
\begin{center}
\vspace{7mm}
{\huge \LARGE Funkcia dvoch premenných, skupina \textit{Tau $\tau$} -\romannumeral1}\\[5mm]
\begin{minipage}{0.95\linewidth}
\begin{center}
{\small \begin{multicols}{2} \raggedright
\textbf{(a)} Vypočítajte zadanú parciálnu deriváciu funkcie \(f(x,y)\). \\[1mm]
\textbf{(b)} Určte gradient \(f(x,y)\) v bode \(A\). \\[1mm]
\textbf{(c)} Napíšte rovnicu dotykovej roviny ku grafu funkcie \(g(x,y)\) v dotykovom bode \(D\). \par \columnbreak
\textbf{(d)} Pomocou lineárnej aproximácie odhadnite funkčnú hodnotu funkcie \(g(x,y)\) v bode, ktorý je posunutý o vektor \((\Delta x, \Delta y)\) od bodu \(D_0\). \\[1mm]
\textbf{(e)} Vypočítajte \(h(x,y)\) v jej lokálnom extréme. \\[1mm]
\textbf{(f)} Nájdite bod viazeného lokálneho extrému funkcie \(h(x,y)\) na priamke. \\[1mm]
\end{multicols}
}
\end{center}
\end{minipage}
\\[4mm]
\begin{minipage}{0.79\linewidth}
\begin{center}
\begin{varwidth}{\linewidth}
\begin{enumerate}
\normalsizerr
\item $f(x,y) = \sqrt[3]{x^2y + 2xy^2}\enspace , \enspace \pdv{f}{y}$\quad \dotfill\; ???\;\dotfill \quad $\frac{2x^2+4xy}{\sqrt[3]{(x^2y+2xy^2)^{2}}}$
\item $A = [4, 1]$\quad \dotfill\; ???\;\dotfill \quad $\left( \frac{10}{3\,\sqrt[3]{24^2}} \;,\; \frac{32}{3\,\sqrt[3]{24^2}} \right)$
\item $g(x,y) = \sin\frac{x}{y} \enspace , \enspace D = [0, 2, ?]$\quad \dotfill\; ???\;\dotfill \quad $x-2z=0$
\item $\Delta x = {-0.01} \enspace,\enspace \Delta y = {-0.01}$\quad \dotfill\; ???\;\dotfill \quad $-0.005$
\item $h(x,y) = 3x^2+4y^2+5xy+2x$\quad \dotfill\; ???\;\dotfill \quad $\nicefrac{-16}{73}$
\item $y = -3x+1$\quad \dotfill\; ???\;\dotfill \quad $\left[\frac{17}{48}\;,\;\frac{-1}{16}\right]$
\end{enumerate}
\end{varwidth}
\end{center}
\end{minipage}
\begin{minipage}{0.20\linewidth}
\begin{center}
{\Huge\bfseries 1.} \\[2mm]
\includegraphics[height=37mm]{images/braille.png}
{\small Braillove písmeno}
\end{center}
\end{minipage}
\end{center}
\end{minipage}
&
\begin{minipage}[c][104.5mm][t]{0.5\linewidth}
\begin{center}
\vspace{7mm}
{\huge \LARGE Funkcia dvoch premenných, skupina \textit{Tau $\tau$} -\romannumeral2}\\[5mm]
\begin{minipage}{0.95\linewidth}
\begin{center}
{\small \begin{multicols}{2} \raggedright
\textbf{(a)} Vypočítajte zadanú parciálnu deriváciu funkcie \(f(x,y)\). \\[1mm]
\textbf{(b)} Určte gradient \(f(x,y)\) v bode \(A\). \\[1mm]
\textbf{(c)} Napíšte rovnicu dotykovej roviny ku grafu funkcie \(g(x,y)\) v dotykovom bode \(D\). \par \columnbreak
\textbf{(d)} Pomocou lineárnej aproximácie odhadnite funkčnú hodnotu funkcie \(g(x,y)\) v bode, ktorý je posunutý o vektor \((\Delta x, \Delta y)\) od bodu \(D_0\). \\[1mm]
\textbf{(e)} Vypočítajte \(h(x,y)\) v jej lokálnom extréme. \\[1mm]
\textbf{(f)} Nájdite bod viazeného lokálneho extrému funkcie \(h(x,y)\) na priamke. \\[1mm]
\end{multicols}
}
\end{center}
\end{minipage}
\\[4mm]
\begin{minipage}{0.79\linewidth}
\begin{center}
\begin{varwidth}{\linewidth}
\begin{enumerate}
\normalsizerr
\item $f(x,y) = \sqrt[3]{6x^2y + 5xy^2}\enspace , \enspace \pdv{f}{x}$\quad \dotfill\; ???\;\dotfill \quad $\frac{12xy+5y^2}{3\,\sqrt[3]{(6x^2y+5xy^2)^{2}}}$
\item $A = [1, 2]$\quad \dotfill\; ???\;\dotfill \quad $\left( \frac{88}{3\,\sqrt[3]{32^2}} \;,\; \frac{26}{3\,\sqrt[3]{32^2}} \right)$
\item $g(x,y) = \sin\frac{x}{y} \enspace , \enspace D = [0, 3, ?]$\quad \dotfill\; ???\;\dotfill \quad $x+3z=0$
\item $\Delta x = {-0.02} \enspace,\enspace \Delta y = {0.3}$\quad \dotfill\; ???\;\dotfill \quad $-0.06$
\item $h(x,y) = -5x^2-y^2-2xy+5x$\quad \dotfill\; ???\;\dotfill \quad $\nicefrac{25}{24}$
\item $y = -3x+3$\quad \dotfill\; ???\;\dotfill \quad $\left[\frac{17}{16}\;,\;\frac{-3}{40}\right]$
\end{enumerate}
\end{varwidth}
\end{center}
\end{minipage}
\begin{minipage}{0.20\linewidth}
\begin{center}
{\Huge\bfseries 2.} \\[2mm]
\includegraphics[height=37mm]{images/braille.png}
{\small Braillove písmeno}
\end{center}
\end{minipage}
\end{center}
\end{minipage}
\\ \hdashline
\begin{minipage}[c][104.5mm][t]{0.5\linewidth}
\begin{center}
\vspace{7mm}
{\huge \LARGE Funkcia dvoch premenných, skupina \textit{Tau $\tau$} -\romannumeral3}\\[5mm]
\begin{minipage}{0.95\linewidth}
\begin{center}
{\small \begin{multicols}{2} \raggedright
\textbf{(a)} Vypočítajte zadanú parciálnu deriváciu funkcie \(f(x,y)\). \\[1mm]
\textbf{(b)} Určte gradient \(f(x,y)\) v bode \(A\). \\[1mm]
\textbf{(c)} Napíšte rovnicu dotykovej roviny ku grafu funkcie \(g(x,y)\) v dotykovom bode \(D\). \par \columnbreak
\textbf{(d)} Pomocou lineárnej aproximácie odhadnite funkčnú hodnotu funkcie \(g(x,y)\) v bode, ktorý je posunutý o vektor \((\Delta x, \Delta y)\) od bodu \(D_0\). \\[1mm]
\textbf{(e)} Vypočítajte \(h(x,y)\) v jej lokálnom extréme. \\[1mm]
\textbf{(f)} Nájdite bod viazeného lokálneho extrému funkcie \(h(x,y)\) na priamke. \\[1mm]
\end{multicols}
}
\end{center}
\end{minipage}
\\[4mm]
\begin{minipage}{0.79\linewidth}
\begin{center}
\begin{varwidth}{\linewidth}
\begin{enumerate}
\normalsizerr
\item $f(x,y) = \sqrt[4]{3x^2y + 4xy^2}\enspace , \enspace \pdv{f}{x}$\quad \dotfill\; ???\;\dotfill \quad $\frac{6xy+4y^2}{4\,\sqrt[4]{(3x^2y+4xy^2)^{3}}}$
\item $A = [2, 1]$\quad \dotfill\; ???\;\dotfill \quad $\left( \frac{16}{4\,\sqrt[4]{20^3}} \;,\; \frac{28}{4\,\sqrt[4]{20^3}} \right)$
\item $g(x,y) = \sin\frac{x}{y} \enspace , \enspace D = [0, 2, ?]$\quad \dotfill\; ???\;\dotfill \quad $x+2z=0$
\item $\Delta x = {-0.01} \enspace,\enspace \Delta y = {0.02}$\quad \dotfill\; ???\;\dotfill \quad $-0.02$
\item $h(x,y) = -4x^2-3y^2+3xy-x$\quad \dotfill\; ???\;\dotfill \quad $\nicefrac{4}{39}$
\item $y = -5x+1$\quad \dotfill\; ???\;\dotfill \quad $\left[\frac{8}{47}\;,\;\frac{7}{32}\right]$
\end{enumerate}
\end{varwidth}
\end{center}
\end{minipage}
\begin{minipage}{0.20\linewidth}
\begin{center}
{\Huge\bfseries 3.} \\[2mm]
\includegraphics[height=37mm]{images/braille.png}
{\small Braillove písmeno}
\end{center}
\end{minipage}
\end{center}
\end{minipage}
&
\begin{minipage}[c][104.5mm][t]{0.5\linewidth}
\begin{center}
\vspace{7mm}
{\huge \LARGE Funkcia dvoch premenných, skupina \textit{Tau $\tau$} -\romannumeral4}\\[5mm]
\begin{minipage}{0.95\linewidth}
\begin{center}
{\small \begin{multicols}{2} \raggedright
\textbf{(a)} Vypočítajte zadanú parciálnu deriváciu funkcie \(f(x,y)\). \\[1mm]
\textbf{(b)} Určte gradient \(f(x,y)\) v bode \(A\). \\[1mm]
\textbf{(c)} Napíšte rovnicu dotykovej roviny ku grafu funkcie \(g(x,y)\) v dotykovom bode \(D\). \par \columnbreak
\textbf{(d)} Pomocou lineárnej aproximácie odhadnite funkčnú hodnotu funkcie \(g(x,y)\) v bode, ktorý je posunutý o vektor \((\Delta x, \Delta y)\) od bodu \(D_0\). \\[1mm]
\textbf{(e)} Vypočítajte \(h(x,y)\) v jej lokálnom extréme. \\[1mm]
\textbf{(f)} Nájdite bod viazeného lokálneho extrému funkcie \(h(x,y)\) na priamke. \\[1mm]
\end{multicols}
}
\end{center}
\end{minipage}
\\[4mm]
\begin{minipage}{0.79\linewidth}
\begin{center}
\begin{varwidth}{\linewidth}
\begin{enumerate}
\normalsizerr
\item $f(x,y) = \sqrt[3]{9x^2y + 2xy^2}\enspace , \enspace \pdv{f}{x}$\quad \dotfill\; ???\;\dotfill \quad $\frac{18xy+2y^2}{3\,\sqrt[3]{(9x^2y+2xy^2)^{2}}}$
\item $A = [2, 2]$\quad \dotfill\; ???\;\dotfill \quad $\left( \frac{160}{3\,\sqrt[3]{88^2}} \;,\; \frac{52}{3\,\sqrt[3]{88^2}} \right)$
\item $g(x,y) = \sin\frac{x}{y} \enspace , \enspace D = [0, 3, ?]$\quad \dotfill\; ???\;\dotfill \quad $x-3z=1$
\item $\Delta x = {-0.01} \enspace,\enspace \Delta y = {0.01}$\quad \dotfill\; ???\;\dotfill \quad $-0.03$
\item $h(x,y) = 2x^2+3y^2-2xy+4x$\quad \dotfill\; ???\;\dotfill \quad $\nicefrac{-12}{7}$
\item $y = -3x-2$\quad \dotfill\; ???\;\dotfill \quad $\left[\frac{2}{5}\;,\;\frac{-4}{35}\right]$
\end{enumerate}
\end{varwidth}
\end{center}
\end{minipage}
\begin{minipage}{0.20\linewidth}
\begin{center}
{\Huge\bfseries 4.} \\[2mm]
\includegraphics[height=37mm]{images/braille.png}
{\small Braillove písmeno}
\end{center}
\end{minipage}
\end{center}
\end{minipage}
%
\end{tabular}
\newpage
\thispagestyle{empty}
\begin{tabular}{c:c}
\begin{minipage}[c][104.5mm][t]{0.5\linewidth}
\begin{center}
\vspace{7mm}
{\huge \LARGE Funkcia dvoch premenných, skupina \textit{Upsilon $\upsilon$} -\romannumeral1}\\[5mm]
\begin{minipage}{0.95\linewidth}
\begin{center}
{\small \begin{multicols}{2} \raggedright
\textbf{(a)} Vypočítajte zadanú parciálnu deriváciu funkcie \(f(x,y)\). \\[1mm]
\textbf{(b)} Určte gradient \(f(x,y)\) v bode \(A\). \\[1mm]
\textbf{(c)} Napíšte rovnicu dotykovej roviny ku grafu funkcie \(g(x,y)\) v dotykovom bode \(D\). \par \columnbreak
\textbf{(d)} Pomocou lineárnej aproximácie odhadnite funkčnú hodnotu funkcie \(g(x,y)\) v bode, ktorý je posunutý o vektor \((\Delta x, \Delta y)\) od bodu \(D_0\). \\[1mm]
\textbf{(e)} Vypočítajte \(h(x,y)\) v jej lokálnom extréme. \\[1mm]
\textbf{(f)} Nájdite bod viazeného lokálneho extrému funkcie \(h(x,y)\) na priamke. \\[1mm]
\end{multicols}
}
\end{center}
\end{minipage}
\\[4mm]
\begin{minipage}{0.79\linewidth}
\begin{center}
\begin{varwidth}{\linewidth}
\begin{enumerate}
\normalsizerr
\item $f(x,y) = \sqrt[4]{2x^2y + 6xy^2}\enspace , \enspace \pdv{f}{y}$\quad \dotfill\; ???\;\dotfill \quad $\frac{2x^2+12xy}{4\,\sqrt[4]{(2x^2y+6xy^2)^{3}}}$
\item $A = [3, 1]$\quad \dotfill\; ???\;\dotfill \quad $\left( \frac{18}{4\,\sqrt[4]{36^3}} \;,\; \frac{54}{4\,\sqrt[4]{36^3}} \right)$
\item $g(x,y) = \sin\frac{x}{y} \enspace , \enspace D = [0, 2, ?]$\quad \dotfill\; ???\;\dotfill \quad $x-2z=0$
\item $\Delta x = {-0.02} \enspace,\enspace \Delta y = {0.2}$\quad \dotfill\; ???\;\dotfill \quad $-0.01$
\item $h(x,y) = 6x^2+4y^2-2xy-x$\quad \dotfill\; ???\;\dotfill \quad $\nicefrac{-1}{25}$
\item $y = -2x+1$\quad \dotfill\; ???\;\dotfill \quad $\left[\frac{-1}{4}\;,\;\frac{7}{26}\right]$
\end{enumerate}
\end{varwidth}
\end{center}
\end{minipage}
\begin{minipage}{0.20\linewidth}
\begin{center}
{\Huge\bfseries 1.} \\[2mm]
\includegraphics[height=37mm]{images/braille.png}
{\small Braillove písmeno}
\end{center}
\end{minipage}
\end{center}
\end{minipage}
&
\begin{minipage}[c][104.5mm][t]{0.5\linewidth}
\begin{center}
\vspace{7mm}
{\huge \LARGE Funkcia dvoch premenných, skupina \textit{Upsilon $\upsilon$} -\romannumeral2}\\[5mm]
\begin{minipage}{0.95\linewidth}
\begin{center}
{\small \begin{multicols}{2} \raggedright
\textbf{(a)} Vypočítajte zadanú parciálnu deriváciu funkcie \(f(x,y)\). \\[1mm]
\textbf{(b)} Určte gradient \(f(x,y)\) v bode \(A\). \\[1mm]
\textbf{(c)} Napíšte rovnicu dotykovej roviny ku grafu funkcie \(g(x,y)\) v dotykovom bode \(D\). \par \columnbreak
\textbf{(d)} Pomocou lineárnej aproximácie odhadnite funkčnú hodnotu funkcie \(g(x,y)\) v bode, ktorý je posunutý o vektor \((\Delta x, \Delta y)\) od bodu \(D_0\). \\[1mm]
\textbf{(e)} Vypočítajte \(h(x,y)\) v jej lokálnom extréme. \\[1mm]
\textbf{(f)} Nájdite bod viazeného lokálneho extrému funkcie \(h(x,y)\) na priamke. \\[1mm]
\end{multicols}
}
\end{center}
\end{minipage}
\\[4mm]
\begin{minipage}{0.79\linewidth}
\begin{center}
\begin{varwidth}{\linewidth}
\begin{enumerate}
\normalsizerr
\item $f(x,y) = \sqrt[4]{3x^2y + 2xy^2}\enspace , \enspace \pdv{f}{x}$\quad \dotfill\; ???\;\dotfill \quad $\frac{6xy+2y^2}{4\,\sqrt[4]{(3x^2y+2xy^2)^{3}}}$
\item $A = [4, 3]$\quad \dotfill\; ???\;\dotfill \quad $\left( \frac{180}{4\,\sqrt[4]{216^3}} \;,\; \frac{96}{4\,\sqrt[4]{216^3}} \right)$
\item $g(x,y) = \ln\frac{x}{y} \enspace , \enspace D=[e^{4},1,?]$\quad \dotfill\; ???\;\dotfill \quad $\frac{x}{e^4} - y - z + 4 = 0$
\item $\Delta x = {0.1} \enspace,\enspace \Delta y = {0.1}$\quad \dotfill\; ???\;\dotfill \quad $3.902$
\item $h(x,y) = 4x^2+7y^2+4xy-x$\quad \dotfill\; ???\;\dotfill \quad $\nicefrac{-7}{96}$
\item $y = x-4$\quad \dotfill\; ???\;\dotfill \quad $\left[\frac{73}{30}\;,\;\frac{-47}{14}\right]$
\end{enumerate}
\end{varwidth}
\end{center}
\end{minipage}
\begin{minipage}{0.20\linewidth}
\begin{center}
{\Huge\bfseries 2.} \\[2mm]
\includegraphics[height=37mm]{images/braille.png}
{\small Braillove písmeno}
\end{center}
\end{minipage}
\end{center}
\end{minipage}
\\ \hdashline
\begin{minipage}[c][104.5mm][t]{0.5\linewidth}
\begin{center}
\vspace{7mm}
{\huge \LARGE Funkcia dvoch premenných, skupina \textit{Upsilon $\upsilon$} -\romannumeral3}\\[5mm]
\begin{minipage}{0.95\linewidth}
\begin{center}
{\small \begin{multicols}{2} \raggedright
\textbf{(a)} Vypočítajte zadanú parciálnu deriváciu funkcie \(f(x,y)\). \\[1mm]
\textbf{(b)} Určte gradient \(f(x,y)\) v bode \(A\). \\[1mm]
\textbf{(c)} Napíšte rovnicu dotykovej roviny ku grafu funkcie \(g(x,y)\) v dotykovom bode \(D\). \par \columnbreak
\textbf{(d)} Pomocou lineárnej aproximácie odhadnite funkčnú hodnotu funkcie \(g(x,y)\) v bode, ktorý je posunutý o vektor \((\Delta x, \Delta y)\) od bodu \(D_0\). \\[1mm]
\textbf{(e)} Vypočítajte \(h(x,y)\) v jej lokálnom extréme. \\[1mm]
\textbf{(f)} Nájdite bod viazeného lokálneho extrému funkcie \(h(x,y)\) na priamke. \\[1mm]
\end{multicols}
}
\end{center}
\end{minipage}
\\[4mm]
\begin{minipage}{0.79\linewidth}
\begin{center}
\begin{varwidth}{\linewidth}
\begin{enumerate}
\normalsizerr
\item $f(x,y) = \sqrt{6x^2y + xy^2}\enspace , \enspace \pdv{f}{y}$\quad \dotfill\; ???\;\dotfill \quad $\frac{6x^2+2xy}{2\,\sqrt{6x^2y+xy^2}}$
\item $A = [2, 1]$\quad \dotfill\; ???\;\dotfill \quad $\left( \frac{25}{2\,\sqrt{26}} \;,\; \frac{56}{2\,\sqrt{26}} \right)$
\item $g(x,y) = \sin\frac{x}{y} \enspace , \enspace D = [0, 3, ?]$\quad \dotfill\; ???\;\dotfill \quad $x-3z=0$
\item $\Delta x = {-0.02} \enspace,\enspace \Delta y = {-0.02}$\quad \dotfill\; ???\;\dotfill \quad $-0.06$
\item $h(x,y) = -3x^2-7y^2+xy-2x$\quad \dotfill\; ???\;\dotfill \quad $\nicefrac{28}{83}$
\item $y = x-3$\quad \dotfill\; ???\;\dotfill \quad $\left[\frac{-47}{18}\;,\;\frac{-17}{18}\right]$
\end{enumerate}
\end{varwidth}
\end{center}
\end{minipage}
\begin{minipage}{0.20\linewidth}
\begin{center}
{\Huge\bfseries 3.} \\[2mm]
\includegraphics[height=37mm]{images/braille.png}
{\small Braillove písmeno}
\end{center}
\end{minipage}
\end{center}
\end{minipage}
&
\begin{minipage}[c][104.5mm][t]{0.5\linewidth}
\begin{center}
\vspace{7mm}
{\huge \LARGE Funkcia dvoch premenných, skupina \textit{Upsilon $\upsilon$} -\romannumeral4}\\[5mm]
\begin{minipage}{0.95\linewidth}
\begin{center}
{\small \begin{multicols}{2} \raggedright
\textbf{(a)} Vypočítajte zadanú parciálnu deriváciu funkcie \(f(x,y)\). \\[1mm]
\textbf{(b)} Určte gradient \(f(x,y)\) v bode \(A\). \\[1mm]
\textbf{(c)} Napíšte rovnicu dotykovej roviny ku grafu funkcie \(g(x,y)\) v dotykovom bode \(D\). \par \columnbreak
\textbf{(d)} Pomocou lineárnej aproximácie odhadnite funkčnú hodnotu funkcie \(g(x,y)\) v bode, ktorý je posunutý o vektor \((\Delta x, \Delta y)\) od bodu \(D_0\). \\[1mm]
\textbf{(e)} Vypočítajte \(h(x,y)\) v jej lokálnom extréme. \\[1mm]
\textbf{(f)} Nájdite bod viazeného lokálneho extrému funkcie \(h(x,y)\) na priamke. \\[1mm]
\end{multicols}
}
\end{center}
\end{minipage}
\\[4mm]
\begin{minipage}{0.79\linewidth}
\begin{center}
\begin{varwidth}{\linewidth}
\begin{enumerate}
\normalsizerr
\item $f(x,y) = \sqrt[3]{2x^2y + 3xy^2}\enspace , \enspace \pdv{f}{y}$\quad \dotfill\; ???\;\dotfill \quad $\frac{2x^2+6xy}{3\,\sqrt[3]{(2x^2y+3xy^2)^{2}}}$
\item $A = [3, 1]$\quad \dotfill\; ???\;\dotfill \quad $\left( \frac{15}{3\,\sqrt[3]{27^2}} \;,\; \frac{36}{3\,\sqrt[3]{27^2}} \right)$
\item $g(x,y) = \sin\frac{x}{y} \enspace , \enspace D = [0, -1, ?]$\quad \dotfill\; ???\;\dotfill \quad $x+z=0$
\item $\Delta x = {-0.3} \enspace,\enspace \Delta y = {-0.03}$\quad \dotfill\; ???\;\dotfill \quad $0.03$
\item $h(x,y) = -5x^2-5y^2+3xy-x$\quad \dotfill\; ???\;\dotfill \quad $\nicefrac{5}{91}$
\item $y = -3x-1$\quad \dotfill\; ???\;\dotfill \quad $\left[\frac{-17}{59}\;,\;\frac{-8}{41}\right]$
\end{enumerate}
\end{varwidth}
\end{center}
\end{minipage}
\begin{minipage}{0.20\linewidth}
\begin{center}
{\Huge\bfseries 4.} \\[2mm]
\includegraphics[height=37mm]{images/braille.png}
{\small Braillove písmeno}
\end{center}
\end{minipage}
\end{center}
\end{minipage}
%
\end{tabular}
\newpage
\thispagestyle{empty}
\begin{tabular}{c:c}
\begin{minipage}[c][104.5mm][t]{0.5\linewidth}
\begin{center}
\vspace{7mm}
{\huge \LARGE Funkcia dvoch premenných, skupina \textit{Phi $\phi$} -\romannumeral1}\\[5mm]
\begin{minipage}{0.95\linewidth}
\begin{center}
{\small \begin{multicols}{2} \raggedright
\textbf{(a)} Vypočítajte zadanú parciálnu deriváciu funkcie \(f(x,y)\). \\[1mm]
\textbf{(b)} Určte gradient \(f(x,y)\) v bode \(A\). \\[1mm]
\textbf{(c)} Napíšte rovnicu dotykovej roviny ku grafu funkcie \(g(x,y)\) v dotykovom bode \(D\). \par \columnbreak
\textbf{(d)} Pomocou lineárnej aproximácie odhadnite funkčnú hodnotu funkcie \(g(x,y)\) v bode, ktorý je posunutý o vektor \((\Delta x, \Delta y)\) od bodu \(D_0\). \\[1mm]
\textbf{(e)} Vypočítajte \(h(x,y)\) v jej lokálnom extréme. \\[1mm]
\textbf{(f)} Nájdite bod viazeného lokálneho extrému funkcie \(h(x,y)\) na priamke. \\[1mm]
\end{multicols}
}
\end{center}
\end{minipage}
\\[4mm]
\begin{minipage}{0.79\linewidth}
\begin{center}
\begin{varwidth}{\linewidth}
\begin{enumerate}
\normalsizerr
\item $f(x,y) = \sqrt[4]{x^2y + xy^2}\enspace , \enspace \pdv{f}{x}$\quad \dotfill\; ???\;\dotfill \quad $\frac{2xy+2y^2}{\sqrt[4]{(x^2y+xy^2)^{3}}}$
\item $A = [3, 1]$\quad \dotfill\; ???\;\dotfill \quad $\left( \frac{7}{4\,\sqrt[4]{12^3}} \;,\; \frac{15}{4\,\sqrt[4]{12^3}} \right)$
\item $g(x,y) = \sin\frac{x}{y} \enspace , \enspace D = [0, -3, ?]$\quad \dotfill\; ???\;\dotfill \quad $x+3z=0$
\item $\Delta x = {0.02} \enspace,\enspace \Delta y = {-0.1}$\quad \dotfill\; ???\;\dotfill \quad $-0.06$
\item $h(x,y) = 6x^2+7y^2+4xy+2x$\quad \dotfill\; ???\;\dotfill \quad $\nicefrac{-27}{152}$
\item $y = x+2$\quad \dotfill\; ???\;\dotfill \quad $\left[\frac{-19}{17}\;,\;\frac{15}{17}\right]$
\end{enumerate}
\end{varwidth}
\end{center}
\end{minipage}
\begin{minipage}{0.20\linewidth}
\begin{center}
{\Huge\bfseries 1.} \\[2mm]
\includegraphics[height=37mm]{images/braille.png}
{\small Braillove písmeno}
\end{center}
\end{minipage}
\end{center}
\end{minipage}
&
\begin{minipage}[c][104.5mm][t]{0.5\linewidth}
\begin{center}
\vspace{7mm}
{\huge \LARGE Funkcia dvoch premenných, skupina \textit{Phi $\phi$} -\romannumeral2}\\[5mm]
\begin{minipage}{0.95\linewidth}
\begin{center}
{\small \begin{multicols}{2} \raggedright
\textbf{(a)} Vypočítajte zadanú parciálnu deriváciu funkcie \(f(x,y)\). \\[1mm]
\textbf{(b)} Určte gradient \(f(x,y)\) v bode \(A\). \\[1mm]
\textbf{(c)} Napíšte rovnicu dotykovej roviny ku grafu funkcie \(g(x,y)\) v dotykovom bode \(D\). \par \columnbreak
\textbf{(d)} Pomocou lineárnej aproximácie odhadnite funkčnú hodnotu funkcie \(g(x,y)\) v bode, ktorý je posunutý o vektor \((\Delta x, \Delta y)\) od bodu \(D_0\). \\[1mm]
\textbf{(e)} Vypočítajte \(h(x,y)\) v jej lokálnom extréme. \\[1mm]
\textbf{(f)} Nájdite bod viazeného lokálneho extrému funkcie \(h(x,y)\) na priamke. \\[1mm]
\end{multicols}
}
\end{center}
\end{minipage}
\\[4mm]
\begin{minipage}{0.79\linewidth}
\begin{center}
\begin{varwidth}{\linewidth}
\begin{enumerate}
\normalsizerr
\item $f(x,y) = \sqrt{9x^2y + 4xy^2}\enspace , \enspace \pdv{f}{y}$\quad \dotfill\; ???\;\dotfill \quad $\frac{9x^2+8xy}{2\,\sqrt{9x^2y+4xy^2}}$
\item $A = [2, 2]$\quad \dotfill\; ???\;\dotfill \quad $\left( \frac{176}{2\,\sqrt{104}} \;,\; \frac{68}{2\,\sqrt{104}} \right)$
\item $g(x,y) = \sin\frac{x}{y} \enspace , \enspace D = [0, -2, ?]$\quad \dotfill\; ???\;\dotfill \quad $x+2z=1$
\item $\Delta x = {-0.1} \enspace,\enspace \Delta y = {-0.1}$\quad \dotfill\; ???\;\dotfill \quad $0.2$
\item $h(x,y) = 2x^2+2y^2-3xy-4x$\quad \dotfill\; ???\;\dotfill \quad $\nicefrac{-32}{7}$
\item $y = -2x+1$\quad \dotfill\; ???\;\dotfill \quad $\left[\frac{-1}{32}\;,\;\frac{1}{16}\right]$
\end{enumerate}
\end{varwidth}
\end{center}
\end{minipage}
\begin{minipage}{0.20\linewidth}
\begin{center}
{\Huge\bfseries 2.} \\[2mm]
\includegraphics[height=37mm]{images/braille.png}
{\small Braillove písmeno}
\end{center}
\end{minipage}
\end{center}
\end{minipage}
\\ \hdashline
\begin{minipage}[c][104.5mm][t]{0.5\linewidth}
\begin{center}
\vspace{7mm}
{\huge \LARGE Funkcia dvoch premenných, skupina \textit{Phi $\phi$} -\romannumeral3}\\[5mm]
\begin{minipage}{0.95\linewidth}
\begin{center}
{\small \begin{multicols}{2} \raggedright
\textbf{(a)} Vypočítajte zadanú parciálnu deriváciu funkcie \(f(x,y)\). \\[1mm]
\textbf{(b)} Určte gradient \(f(x,y)\) v bode \(A\). \\[1mm]
\textbf{(c)} Napíšte rovnicu dotykovej roviny ku grafu funkcie \(g(x,y)\) v dotykovom bode \(D\). \par \columnbreak
\textbf{(d)} Pomocou lineárnej aproximácie odhadnite funkčnú hodnotu funkcie \(g(x,y)\) v bode, ktorý je posunutý o vektor \((\Delta x, \Delta y)\) od bodu \(D_0\). \\[1mm]
\textbf{(e)} Vypočítajte \(h(x,y)\) v jej lokálnom extréme. \\[1mm]
\textbf{(f)} Nájdite bod viazeného lokálneho extrému funkcie \(h(x,y)\) na priamke. \\[1mm]
\end{multicols}
}
\end{center}
\end{minipage}
\\[4mm]
\begin{minipage}{0.79\linewidth}
\begin{center}
\begin{varwidth}{\linewidth}
\begin{enumerate}
\normalsizerr
\item $f(x,y) = \sqrt[4]{6x^2y + 2xy^2}\enspace , \enspace \pdv{f}{x}$\quad \dotfill\; ???\;\dotfill \quad $\frac{12xy+2y^2}{4\,\sqrt[4]{(6x^2y+2xy^2)^{3}}}$
\item $A = [1, 1]$\quad \dotfill\; ???\;\dotfill \quad $\left( \frac{28}{4\,\sqrt[4]{8^3}} \;,\; \frac{10}{4\,\sqrt[4]{8^3}} \right)$
\item $g(x,y) = \sin\frac{x}{y} \enspace , \enspace D = [0, -2, ?]$\quad \dotfill\; ???\;\dotfill \quad $x+2z=0$
\item $\Delta x = {0.3} \enspace,\enspace \Delta y = {0.2}$\quad \dotfill\; ???\;\dotfill \quad $-0.1$
\item $h(x,y) = -5x^2-y^2+2xy-9x$\quad \dotfill\; ???\;\dotfill \quad $\nicefrac{81}{16}$
\item $y = -2x-1$\quad \dotfill\; ???\;\dotfill \quad $\left[\frac{-15}{26}\;,\;\frac{2}{13}\right]$
\end{enumerate}
\end{varwidth}
\end{center}
\end{minipage}
\begin{minipage}{0.20\linewidth}
\begin{center}
{\Huge\bfseries 3.} \\[2mm]
\includegraphics[height=37mm]{images/braille.png}
{\small Braillove písmeno}
\end{center}
\end{minipage}
\end{center}
\end{minipage}
&
\begin{minipage}[c][104.5mm][t]{0.5\linewidth}
\begin{center}
\vspace{7mm}
{\huge \LARGE Funkcia dvoch premenných, skupina \textit{Phi $\phi$} -\romannumeral4}\\[5mm]
\begin{minipage}{0.95\linewidth}
\begin{center}
{\small \begin{multicols}{2} \raggedright
\textbf{(a)} Vypočítajte zadanú parciálnu deriváciu funkcie \(f(x,y)\). \\[1mm]
\textbf{(b)} Určte gradient \(f(x,y)\) v bode \(A\). \\[1mm]
\textbf{(c)} Napíšte rovnicu dotykovej roviny ku grafu funkcie \(g(x,y)\) v dotykovom bode \(D\). \par \columnbreak
\textbf{(d)} Pomocou lineárnej aproximácie odhadnite funkčnú hodnotu funkcie \(g(x,y)\) v bode, ktorý je posunutý o vektor \((\Delta x, \Delta y)\) od bodu \(D_0\). \\[1mm]
\textbf{(e)} Vypočítajte \(h(x,y)\) v jej lokálnom extréme. \\[1mm]
\textbf{(f)} Nájdite bod viazeného lokálneho extrému funkcie \(h(x,y)\) na priamke. \\[1mm]
\end{multicols}
}
\end{center}
\end{minipage}
\\[4mm]
\begin{minipage}{0.79\linewidth}
\begin{center}
\begin{varwidth}{\linewidth}
\begin{enumerate}
\normalsizerr
\item $f(x,y) = \sqrt[4]{3x^2y + xy^2}\enspace , \enspace \pdv{f}{y}$\quad \dotfill\; ???\;\dotfill \quad $\frac{3x^2+2xy}{4\,\sqrt[4]{(3x^2y+xy^2)^{3}}}$
\item $A = [3, 2]$\quad \dotfill\; ???\;\dotfill \quad $\left( \frac{80}{4\,\sqrt[4]{66^3}} \;,\; \frac{39}{4\,\sqrt[4]{66^3}} \right)$
\item $g(x,y) = \ln\frac{x}{y} \enspace , \enspace D=[e^{3},1,?]$\quad \dotfill\; ???\;\dotfill \quad $\frac{x}{e^3} - y - z + 3 = 0$
\item $\Delta x = {0.2} \enspace,\enspace \Delta y = {0.03}$\quad \dotfill\; ???\;\dotfill \quad $2.81$
\item $h(x,y) = 2x^2+6y^2+2xy+4x$\quad \dotfill\; ???\;\dotfill \quad $\nicefrac{-24}{11}$
\item $y = x+3$\quad \dotfill\; ???\;\dotfill \quad $\left[\frac{13}{10}\;,\;\frac{7}{10}\right]$
\end{enumerate}
\end{varwidth}
\end{center}
\end{minipage}
\begin{minipage}{0.20\linewidth}
\begin{center}
{\Huge\bfseries 4.} \\[2mm]
\includegraphics[height=37mm]{images/braille.png}
{\small Braillove písmeno}
\end{center}
\end{minipage}
\end{center}
\end{minipage}
%
\end{tabular}
\newpage
\thispagestyle{empty}
\begin{tabular}{c:c}
\begin{minipage}[c][104.5mm][t]{0.5\linewidth}
\begin{center}
\vspace{7mm}
{\huge \LARGE Funkcia dvoch premenných, skupina \textit{Chi $\chi$} -\romannumeral1}\\[5mm]
\begin{minipage}{0.95\linewidth}
\begin{center}
{\small \begin{multicols}{2} \raggedright
\textbf{(a)} Vypočítajte zadanú parciálnu deriváciu funkcie \(f(x,y)\). \\[1mm]
\textbf{(b)} Určte gradient \(f(x,y)\) v bode \(A\). \\[1mm]
\textbf{(c)} Napíšte rovnicu dotykovej roviny ku grafu funkcie \(g(x,y)\) v dotykovom bode \(D\). \par \columnbreak
\textbf{(d)} Pomocou lineárnej aproximácie odhadnite funkčnú hodnotu funkcie \(g(x,y)\) v bode, ktorý je posunutý o vektor \((\Delta x, \Delta y)\) od bodu \(D_0\). \\[1mm]
\textbf{(e)} Vypočítajte \(h(x,y)\) v jej lokálnom extréme. \\[1mm]
\textbf{(f)} Nájdite bod viazeného lokálneho extrému funkcie \(h(x,y)\) na priamke. \\[1mm]
\end{multicols}
}
\end{center}
\end{minipage}
\\[4mm]
\begin{minipage}{0.79\linewidth}
\begin{center}
\begin{varwidth}{\linewidth}
\begin{enumerate}
\normalsizerr
\item $f(x,y) = \sqrt[4]{6x^2y + 7xy^2}\enspace , \enspace \pdv{f}{x}$\quad \dotfill\; ???\;\dotfill \quad $\frac{12xy+7y^2}{4\,\sqrt[4]{(6x^2y+7xy^2)^{3}}}$
\item $A = [3, 3]$\quad \dotfill\; ???\;\dotfill \quad $\left( \frac{171}{4\,\sqrt[4]{351^3}} \;,\; \frac{180}{4\,\sqrt[4]{351^3}} \right)$
\item $g(x,y) = \ln\frac{x}{y} \enspace , \enspace D=[e^{5},1,?]$\quad \dotfill\; ???\;\dotfill \quad $\frac{x}{e^5} - y - z + 5 = 0$
\item $\Delta x = {0.2} \enspace,\enspace \Delta y = {0.3}$\quad \dotfill\; ???\;\dotfill \quad $4.701$
\item $h(x,y) = -5x^2-7y^2+2xy+3x$\quad \dotfill\; ???\;\dotfill \quad $\nicefrac{8}{17}$
\item $y = x-2$\quad \dotfill\; ???\;\dotfill \quad $\left[\frac{-29}{20}\;,\;\frac{-13}{20}\right]$
\end{enumerate}
\end{varwidth}
\end{center}
\end{minipage}
\begin{minipage}{0.20\linewidth}
\begin{center}
{\Huge\bfseries 1.} \\[2mm]
\includegraphics[height=37mm]{images/braille.png}
{\small Braillove písmeno}
\end{center}
\end{minipage}
\end{center}
\end{minipage}
&
\begin{minipage}[c][104.5mm][t]{0.5\linewidth}
\begin{center}
\vspace{7mm}
{\huge \LARGE Funkcia dvoch premenných, skupina \textit{Chi $\chi$} -\romannumeral2}\\[5mm]
\begin{minipage}{0.95\linewidth}
\begin{center}
{\small \begin{multicols}{2} \raggedright
\textbf{(a)} Vypočítajte zadanú parciálnu deriváciu funkcie \(f(x,y)\). \\[1mm]
\textbf{(b)} Určte gradient \(f(x,y)\) v bode \(A\). \\[1mm]
\textbf{(c)} Napíšte rovnicu dotykovej roviny ku grafu funkcie \(g(x,y)\) v dotykovom bode \(D\). \par \columnbreak
\textbf{(d)} Pomocou lineárnej aproximácie odhadnite funkčnú hodnotu funkcie \(g(x,y)\) v bode, ktorý je posunutý o vektor \((\Delta x, \Delta y)\) od bodu \(D_0\). \\[1mm]
\textbf{(e)} Vypočítajte \(h(x,y)\) v jej lokálnom extréme. \\[1mm]
\textbf{(f)} Nájdite bod viazeného lokálneho extrému funkcie \(h(x,y)\) na priamke. \\[1mm]
\end{multicols}
}
\end{center}
\end{minipage}
\\[4mm]
\begin{minipage}{0.79\linewidth}
\begin{center}
\begin{varwidth}{\linewidth}
\begin{enumerate}
\normalsizerr
\item $f(x,y) = \sqrt[4]{2x^2y + 4xy^2}\enspace , \enspace \pdv{f}{x}$\quad \dotfill\; ???\;\dotfill \quad $\frac{4xy+4y^2}{4\,\sqrt[4]{(2x^2y+4xy^2)^{3}}}$
\item $A = [1, 1]$\quad \dotfill\; ???\;\dotfill \quad $\left( \frac{8}{4\,\sqrt[4]{6^3}} \;,\; \frac{10}{4\,\sqrt[4]{6^3}} \right)$
\item $g(x,y) = \sin\frac{x}{y} \enspace , \enspace D = [0, 2, ?]$\quad \dotfill\; ???\;\dotfill \quad $x-2z=0$
\item $\Delta x = {0.01} \enspace,\enspace \Delta y = {0.1}$\quad \dotfill\; ???\;\dotfill \quad $0.02$
\item $h(x,y) = -3x^2-3y^2+5xy+2x$\quad \dotfill\; ???\;\dotfill \quad $\nicefrac{12}{61}$
\item $y = -x+1$\quad \dotfill\; ???\;\dotfill \quad $\left[\frac{13}{22}\;,\;\frac{9}{2}\right]$
\end{enumerate}
\end{varwidth}
\end{center}
\end{minipage}
\begin{minipage}{0.20\linewidth}
\begin{center}
{\Huge\bfseries 2.} \\[2mm]
\includegraphics[height=37mm]{images/braille.png}
{\small Braillove písmeno}
\end{center}
\end{minipage}
\end{center}
\end{minipage}
\\ \hdashline
\begin{minipage}[c][104.5mm][t]{0.5\linewidth}
\begin{center}
\vspace{7mm}
{\huge \LARGE Funkcia dvoch premenných, skupina \textit{Chi $\chi$} -\romannumeral3}\\[5mm]
\begin{minipage}{0.95\linewidth}
\begin{center}
{\small \begin{multicols}{2} \raggedright
\textbf{(a)} Vypočítajte zadanú parciálnu deriváciu funkcie \(f(x,y)\). \\[1mm]
\textbf{(b)} Určte gradient \(f(x,y)\) v bode \(A\). \\[1mm]
\textbf{(c)} Napíšte rovnicu dotykovej roviny ku grafu funkcie \(g(x,y)\) v dotykovom bode \(D\). \par \columnbreak
\textbf{(d)} Pomocou lineárnej aproximácie odhadnite funkčnú hodnotu funkcie \(g(x,y)\) v bode, ktorý je posunutý o vektor \((\Delta x, \Delta y)\) od bodu \(D_0\). \\[1mm]
\textbf{(e)} Vypočítajte \(h(x,y)\) v jej lokálnom extréme. \\[1mm]
\textbf{(f)} Nájdite bod viazeného lokálneho extrému funkcie \(h(x,y)\) na priamke. \\[1mm]
\end{multicols}
}
\end{center}
\end{minipage}
\\[4mm]
\begin{minipage}{0.79\linewidth}
\begin{center}
\begin{varwidth}{\linewidth}
\begin{enumerate}
\normalsizerr
\item $f(x,y) = \sqrt[4]{6x^2y + 5xy^2}\enspace , \enspace \pdv{f}{x}$\quad \dotfill\; ???\;\dotfill \quad $\frac{12xy+5y^2}{4\,\sqrt[4]{(6x^2y+5xy^2)^{3}}}$
\item $A = [1, 4]$\quad \dotfill\; ???\;\dotfill \quad $\left( \frac{256}{4\,\sqrt[4]{104^3}} \;,\; \frac{46}{4\,\sqrt[4]{104^3}} \right)$
\item $g(x,y) = \ln\frac{x}{y} \enspace , \enspace D=[1,1,?]$\quad \dotfill\; ???\;\dotfill \quad $x + y + z + 0 = 0$
\item $\Delta x = {0.2} \enspace,\enspace \Delta y = {-0.03}$\quad \dotfill\; ???\;\dotfill \quad $0.17$
\item $h(x,y) = -3x^2-6y^2-xy-2x$\quad \dotfill\; ???\;\dotfill \quad $\nicefrac{24}{73}$
\item $y = x+6$\quad \dotfill\; ???\;\dotfill \quad $\left[-4\;,\;2\right]$
\end{enumerate}
\end{varwidth}
\end{center}
\end{minipage}
\begin{minipage}{0.20\linewidth}
\begin{center}
{\Huge\bfseries 3.} \\[2mm]
\includegraphics[height=37mm]{images/braille.png}
{\small Braillove písmeno}
\end{center}
\end{minipage}
\end{center}
\end{minipage}
&
\begin{minipage}[c][104.5mm][t]{0.5\linewidth}
\begin{center}
\vspace{7mm}
{\huge \LARGE Funkcia dvoch premenných, skupina \textit{Chi $\chi$} -\romannumeral4}\\[5mm]
\begin{minipage}{0.95\linewidth}
\begin{center}
{\small \begin{multicols}{2} \raggedright
\textbf{(a)} Vypočítajte zadanú parciálnu deriváciu funkcie \(f(x,y)\). \\[1mm]
\textbf{(b)} Určte gradient \(f(x,y)\) v bode \(A\). \\[1mm]
\textbf{(c)} Napíšte rovnicu dotykovej roviny ku grafu funkcie \(g(x,y)\) v dotykovom bode \(D\). \par \columnbreak
\textbf{(d)} Pomocou lineárnej aproximácie odhadnite funkčnú hodnotu funkcie \(g(x,y)\) v bode, ktorý je posunutý o vektor \((\Delta x, \Delta y)\) od bodu \(D_0\). \\[1mm]
\textbf{(e)} Vypočítajte \(h(x,y)\) v jej lokálnom extréme. \\[1mm]
\textbf{(f)} Nájdite bod viazeného lokálneho extrému funkcie \(h(x,y)\) na priamke. \\[1mm]
\end{multicols}
}
\end{center}
\end{minipage}
\\[4mm]
\begin{minipage}{0.79\linewidth}
\begin{center}
\begin{varwidth}{\linewidth}
\begin{enumerate}
\normalsizerr
\item $f(x,y) = \sqrt[4]{2x^2y + 2xy^2}\enspace , \enspace \pdv{f}{x}$\quad \dotfill\; ???\;\dotfill \quad $\frac{4xy-2y^2}{4\,\sqrt[4]{(2x^2y+2xy^2)^{3}}}$
\item $A = [4, 1]$\quad \dotfill\; ???\;\dotfill \quad $\left( \frac{18}{4\,\sqrt[4]{40^3}} \;,\; \frac{48}{4\,\sqrt[4]{40^3}} \right)$
\item $g(x,y) = \ln\frac{x}{y} \enspace , \enspace D=[e,1,?]$\quad \dotfill\; ???\;\dotfill \quad $\frac{x}{e} - y - z + 2 = 0$
\item $\Delta x = {0.3} \enspace,\enspace \Delta y = {0.2}$\quad \dotfill\; ???\;\dotfill \quad $0.81$
\item $h(x,y) = 5x^2+2y^2-xy+3x$\quad \dotfill\; ???\;\dotfill \quad $\nicefrac{-17}{39}$
\item $y = -x+2$\quad \dotfill\; ???\;\dotfill \quad $\left[\frac{7}{16}\;,\;\frac{25}{16}\right]$
\end{enumerate}
\end{varwidth}
\end{center}
\end{minipage}
\begin{minipage}{0.20\linewidth}
\begin{center}
{\Huge\bfseries 4.} \\[2mm]
\includegraphics[height=37mm]{images/braille.png}
{\small Braillove písmeno}
\end{center}
\end{minipage}
\end{center}
\end{minipage}
%
\end{tabular}
\newpage
\thispagestyle{empty}
\begin{tabular}{c:c}
\begin{minipage}[c][104.5mm][t]{0.5\linewidth}
\begin{center}
\vspace{7mm}
{\huge \LARGE Funkcia dvoch premenných, skupina \textit{Psi $\psi$} -\romannumeral1}\\[5mm]
\begin{minipage}{0.95\linewidth}
\begin{center}
{\small \begin{multicols}{2} \raggedright
\textbf{(a)} Vypočítajte zadanú parciálnu deriváciu funkcie \(f(x,y)\). \\[1mm]
\textbf{(b)} Určte gradient \(f(x,y)\) v bode \(A\). \\[1mm]
\textbf{(c)} Napíšte rovnicu dotykovej roviny ku grafu funkcie \(g(x,y)\) v dotykovom bode \(D\). \par \columnbreak
\textbf{(d)} Pomocou lineárnej aproximácie odhadnite funkčnú hodnotu funkcie \(g(x,y)\) v bode, ktorý je posunutý o vektor \((\Delta x, \Delta y)\) od bodu \(D_0\). \\[1mm]
\textbf{(e)} Vypočítajte \(h(x,y)\) v jej lokálnom extréme. \\[1mm]
\textbf{(f)} Nájdite bod viazeného lokálneho extrému funkcie \(h(x,y)\) na priamke. \\[1mm]
\end{multicols}
}
\end{center}
\end{minipage}
\\[4mm]
\begin{minipage}{0.79\linewidth}
\begin{center}
\begin{varwidth}{\linewidth}
\begin{enumerate}
\normalsizerr
\item $f(x,y) = \sqrt{7x^2y + 7xy^2}\enspace , \enspace \pdv{f}{y}$\quad \dotfill\; ???\;\dotfill \quad $\frac{7x^2-14xy}{\sqrt{7x^2y+7xy^2}}$
\item $A = [2, 3]$\quad \dotfill\; ???\;\dotfill \quad $\left( \frac{147}{2\,\sqrt{210}} \;,\; \frac{112}{2\,\sqrt{210}} \right)$
\item $g(x,y) = \sin\frac{x}{y} \enspace , \enspace D = [0, -3, ?]$\quad \dotfill\; ???\;\dotfill \quad $x+3z=0$
\item $\Delta x = {0.02} \enspace,\enspace \Delta y = {0.01}$\quad \dotfill\; ???\;\dotfill \quad $-0.06$
\item $h(x,y) = 2x^2+6y^2+2xy+2x$\quad \dotfill\; ???\;\dotfill \quad $\nicefrac{-6}{11}$
\item $y = x-4$\quad \dotfill\; ???\;\dotfill \quad $\left[\frac{27}{10}\;,\;\frac{-13}{10}\right]$
\end{enumerate}
\end{varwidth}
\end{center}
\end{minipage}
\begin{minipage}{0.20\linewidth}
\begin{center}
{\Huge\bfseries 1.} \\[2mm]
\includegraphics[height=37mm]{images/braille.png}
{\small Braillove písmeno}
\end{center}
\end{minipage}
\end{center}
\end{minipage}
&
\begin{minipage}[c][104.5mm][t]{0.5\linewidth}
\begin{center}
\vspace{7mm}
{\huge \LARGE Funkcia dvoch premenných, skupina \textit{Psi $\psi$} -\romannumeral2}\\[5mm]
\begin{minipage}{0.95\linewidth}
\begin{center}
{\small \begin{multicols}{2} \raggedright
\textbf{(a)} Vypočítajte zadanú parciálnu deriváciu funkcie \(f(x,y)\). \\[1mm]
\textbf{(b)} Určte gradient \(f(x,y)\) v bode \(A\). \\[1mm]
\textbf{(c)} Napíšte rovnicu dotykovej roviny ku grafu funkcie \(g(x,y)\) v dotykovom bode \(D\). \par \columnbreak
\textbf{(d)} Pomocou lineárnej aproximácie odhadnite funkčnú hodnotu funkcie \(g(x,y)\) v bode, ktorý je posunutý o vektor \((\Delta x, \Delta y)\) od bodu \(D_0\). \\[1mm]
\textbf{(e)} Vypočítajte \(h(x,y)\) v jej lokálnom extréme. \\[1mm]
\textbf{(f)} Nájdite bod viazeného lokálneho extrému funkcie \(h(x,y)\) na priamke. \\[1mm]
\end{multicols}
}
\end{center}
\end{minipage}
\\[4mm]
\begin{minipage}{0.79\linewidth}
\begin{center}
\begin{varwidth}{\linewidth}
\begin{enumerate}
\normalsizerr
\item $f(x,y) = \sqrt{x^2y + 5xy^2}\enspace , \enspace \pdv{f}{x}$\quad \dotfill\; ???\;\dotfill \quad $\frac{2xy-5y^2}{2\,\sqrt{x^2y+5xy^2}}$
\item $A = [3, 4]$\quad \dotfill\; ???\;\dotfill \quad $\left( \frac{104}{2\,\sqrt{276}} \;,\; \frac{129}{2\,\sqrt{276}} \right)$
\item $g(x,y) = \ln\frac{x}{y} \enspace , \enspace D=[e^{3},1,?]$\quad \dotfill\; ???\;\dotfill \quad $\frac{x}{e^4} - y - z + 4 = 0$
\item $\Delta x = {0.1} \enspace,\enspace \Delta y = {0.2}$\quad \dotfill\; ???\;\dotfill \quad $2.805$
\item $h(x,y) = -4x^2-3y^2-xy-3x$\quad \dotfill\; ???\;\dotfill \quad $\nicefrac{28}{47}$
\item $y = -5x-1$\quad \dotfill\; ???\;\dotfill \quad $\left[\frac{-8}{37}\;,\;\frac{3}{37}\right]$
\end{enumerate}
\end{varwidth}
\end{center}
\end{minipage}
\begin{minipage}{0.20\linewidth}
\begin{center}
{\Huge\bfseries 2.} \\[2mm]
\includegraphics[height=37mm]{images/braille.png}
{\small Braillove písmeno}
\end{center}
\end{minipage}
\end{center}
\end{minipage}
\\ \hdashline
\begin{minipage}[c][104.5mm][t]{0.5\linewidth}
\begin{center}
\vspace{7mm}
{\huge \LARGE Funkcia dvoch premenných, skupina \textit{Psi $\psi$} -\romannumeral3}\\[5mm]
\begin{minipage}{0.95\linewidth}
\begin{center}
{\small \begin{multicols}{2} \raggedright
\textbf{(a)} Vypočítajte zadanú parciálnu deriváciu funkcie \(f(x,y)\). \\[1mm]
\textbf{(b)} Určte gradient \(f(x,y)\) v bode \(A\). \\[1mm]
\textbf{(c)} Napíšte rovnicu dotykovej roviny ku grafu funkcie \(g(x,y)\) v dotykovom bode \(D\). \par \columnbreak
\textbf{(d)} Pomocou lineárnej aproximácie odhadnite funkčnú hodnotu funkcie \(g(x,y)\) v bode, ktorý je posunutý o vektor \((\Delta x, \Delta y)\) od bodu \(D_0\). \\[1mm]
\textbf{(e)} Vypočítajte \(h(x,y)\) v jej lokálnom extréme. \\[1mm]
\textbf{(f)} Nájdite bod viazeného lokálneho extrému funkcie \(h(x,y)\) na priamke. \\[1mm]
\end{multicols}
}
\end{center}
\end{minipage}
\\[4mm]
\begin{minipage}{0.79\linewidth}
\begin{center}
\begin{varwidth}{\linewidth}
\begin{enumerate}
\normalsizerr
\item $f(x,y) = \sqrt[4]{8x^2y + 4xy^2}\enspace , \enspace \pdv{f}{y}$\quad \dotfill\; ???\;\dotfill \quad $\frac{8x^2+8xy}{4\,\sqrt[4]{(8x^2y+4xy^2)^{3}}}$
\item $A = [1, 3]$\quad \dotfill\; ???\;\dotfill \quad $\left( \frac{84}{4\,\sqrt[4]{60^3}} \;,\; \frac{32}{4\,\sqrt[4]{60^3}} \right)$
\item $g(x,y) = \sin\frac{x}{y} \enspace , \enspace D = [0, -3, ?]$\quad \dotfill\; ???\;\dotfill \quad $x-3z=0$
\item $\Delta x = {0.01} \enspace,\enspace \Delta y = {0.01}$\quad \dotfill\; ???\;\dotfill \quad $-0.007$
\item $h(x,y) = 3x^2+5y^2-xy+4x$\quad \dotfill\; ???\;\dotfill \quad $\nicefrac{-80}{59}$
\item $y = x+1$\quad \dotfill\; ???\;\dotfill \quad $\left[\frac{-13}{14}\;,\;\frac{1}{14}\right]$
\end{enumerate}
\end{varwidth}
\end{center}
\end{minipage}
\begin{minipage}{0.20\linewidth}
\begin{center}
{\Huge\bfseries 3.} \\[2mm]
\includegraphics[height=37mm]{images/braille.png}
{\small Braillove písmeno}
\end{center}
\end{minipage}
\end{center}
\end{minipage}
&
\begin{minipage}[c][104.5mm][t]{0.5\linewidth}
\begin{center}
\vspace{7mm}
{\huge \LARGE Funkcia dvoch premenných, skupina \textit{Psi $\psi$} -\romannumeral4}\\[5mm]
\begin{minipage}{0.95\linewidth}
\begin{center}
{\small \begin{multicols}{2} \raggedright
\textbf{(a)} Vypočítajte zadanú parciálnu deriváciu funkcie \(f(x,y)\). \\[1mm]
\textbf{(b)} Určte gradient \(f(x,y)\) v bode \(A\). \\[1mm]
\textbf{(c)} Napíšte rovnicu dotykovej roviny ku grafu funkcie \(g(x,y)\) v dotykovom bode \(D\). \par \columnbreak
\textbf{(d)} Pomocou lineárnej aproximácie odhadnite funkčnú hodnotu funkcie \(g(x,y)\) v bode, ktorý je posunutý o vektor \((\Delta x, \Delta y)\) od bodu \(D_0\). \\[1mm]
\textbf{(e)} Vypočítajte \(h(x,y)\) v jej lokálnom extréme. \\[1mm]
\textbf{(f)} Nájdite bod viazeného lokálneho extrému funkcie \(h(x,y)\) na priamke. \\[1mm]
\end{multicols}
}
\end{center}
\end{minipage}
\\[4mm]
\begin{minipage}{0.79\linewidth}
\begin{center}
\begin{varwidth}{\linewidth}
\begin{enumerate}
\normalsizerr
\item $f(x,y) = \sqrt[3]{x^2y + 6xy^2}\enspace , \enspace \pdv{f}{x}$\quad \dotfill\; ???\;\dotfill \quad $\frac{2xy+6y^2}{3\,\sqrt[3]{(x^2y+6xy^2)^{2}}}$
\item $A = [1, 1]$\quad \dotfill\; ???\;\dotfill \quad $\left( \frac{8}{3\,\sqrt[3]{7^2}} \;,\; \frac{26}{3\,\sqrt[3]{7^2}} \right)$
\item $g(x,y) = \sin\frac{x}{y} \enspace , \enspace D = [0, -3, ?]$\quad \dotfill\; ???\;\dotfill \quad $x+3z=1$
\item $\Delta x = {0.02} \enspace,\enspace \Delta y = {-0.01}$\quad \dotfill\; ???\;\dotfill \quad $0.003$
\item $h(x,y) = -6x^2-4y^2-2xy-2x$\quad \dotfill\; ???\;\dotfill \quad $\nicefrac{17}{92}$
\item $y = -x-9$\quad \dotfill\; ???\;\dotfill \quad $\left[\frac{-7}{2}\;,\;\frac{-11}{3}\right]$
\end{enumerate}
\end{varwidth}
\end{center}
\end{minipage}
\begin{minipage}{0.20\linewidth}
\begin{center}
{\Huge\bfseries 4.} \\[2mm]
\includegraphics[height=37mm]{images/braille.png}
{\small Braillove písmeno}
\end{center}
\end{minipage}
\end{center}
\end{minipage}
%
\end{tabular}
\newpage
\thispagestyle{empty}
\begin{tabular}{c:c}
\begin{minipage}[c][104.5mm][t]{0.5\linewidth}
\begin{center}
\vspace{7mm}
{\huge \LARGE Funkcia dvoch premenných, skupina \textit{Omega $\omega$} -\romannumeral1}\\[5mm]
\begin{minipage}{0.95\linewidth}
\begin{center}
{\small \begin{multicols}{2} \raggedright
\textbf{(a)} Vypočítajte zadanú parciálnu deriváciu funkcie \(f(x,y)\). \\[1mm]
\textbf{(b)} Určte gradient \(f(x,y)\) v bode \(A\). \\[1mm]
\textbf{(c)} Napíšte rovnicu dotykovej roviny ku grafu funkcie \(g(x,y)\) v dotykovom bode \(D\). \par \columnbreak
\textbf{(d)} Pomocou lineárnej aproximácie odhadnite funkčnú hodnotu funkcie \(g(x,y)\) v bode, ktorý je posunutý o vektor \((\Delta x, \Delta y)\) od bodu \(D_0\). \\[1mm]
\textbf{(e)} Vypočítajte \(h(x,y)\) v jej lokálnom extréme. \\[1mm]
\textbf{(f)} Nájdite bod viazeného lokálneho extrému funkcie \(h(x,y)\) na priamke. \\[1mm]
\end{multicols}
}
\end{center}
\end{minipage}
\\[4mm]
\begin{minipage}{0.79\linewidth}
\begin{center}
\begin{varwidth}{\linewidth}
\begin{enumerate}
\normalsizerr
\item $f(x,y) = \sqrt{9x^2y + xy^2}\enspace , \enspace \pdv{f}{y}$\quad \dotfill\; ???\;\dotfill \quad $\frac{9x^2-2xy}{\sqrt{9x^2y+xy^2}}$
\item $A = [2, 1]$\quad \dotfill\; ???\;\dotfill \quad $\left( \frac{37}{2\,\sqrt{38}} \;,\; \frac{40}{2\,\sqrt{38}} \right)$
\item $g(x,y) = \ln\frac{x}{y} \enspace , \enspace D=[1,1,?]$\quad \dotfill\; ???\;\dotfill \quad $x + y + z + 0 = 0$
\item $\Delta x = {-0.02} \enspace,\enspace \Delta y = {-0.3}$\quad \dotfill\; ???\;\dotfill \quad $0.0$
\item $h(x,y) = -3x^2-4y^2-4xy+5x$\quad \dotfill\; ???\;\dotfill \quad $\nicefrac{25}{8}$
\item $y = 2x-1$\quad \dotfill\; ???\;\dotfill \quad $\left[\frac{25}{54}\;,\;\frac{-2}{27}\right]$
\end{enumerate}
\end{varwidth}
\end{center}
\end{minipage}
\begin{minipage}{0.20\linewidth}
\begin{center}
{\Huge\bfseries 1.} \\[2mm]
\includegraphics[height=37mm]{images/braille.png}
{\small Braillove písmeno}
\end{center}
\end{minipage}
\end{center}
\end{minipage}
&
\begin{minipage}[c][104.5mm][t]{0.5\linewidth}
\begin{center}
\vspace{7mm}
{\huge \LARGE Funkcia dvoch premenných, skupina \textit{Omega $\omega$} -\romannumeral2}\\[5mm]
\begin{minipage}{0.95\linewidth}
\begin{center}
{\small \begin{multicols}{2} \raggedright
\textbf{(a)} Vypočítajte zadanú parciálnu deriváciu funkcie \(f(x,y)\). \\[1mm]
\textbf{(b)} Určte gradient \(f(x,y)\) v bode \(A\). \\[1mm]
\textbf{(c)} Napíšte rovnicu dotykovej roviny ku grafu funkcie \(g(x,y)\) v dotykovom bode \(D\). \par \columnbreak
\textbf{(d)} Pomocou lineárnej aproximácie odhadnite funkčnú hodnotu funkcie \(g(x,y)\) v bode, ktorý je posunutý o vektor \((\Delta x, \Delta y)\) od bodu \(D_0\). \\[1mm]
\textbf{(e)} Vypočítajte \(h(x,y)\) v jej lokálnom extréme. \\[1mm]
\textbf{(f)} Nájdite bod viazeného lokálneho extrému funkcie \(h(x,y)\) na priamke. \\[1mm]
\end{multicols}
}
\end{center}
\end{minipage}
\\[4mm]
\begin{minipage}{0.79\linewidth}
\begin{center}
\begin{varwidth}{\linewidth}
\begin{enumerate}
\normalsizerr
\item $f(x,y) = \sqrt{4x^2y + 3xy^2}\enspace , \enspace \pdv{f}{y}$\quad \dotfill\; ???\;\dotfill \quad $\frac{4x^2-6xy}{2\,\sqrt{4x^2y+3xy^2}}$
\item $A = [3, 4]$\quad \dotfill\; ???\;\dotfill \quad $\left( \frac{144}{2\,\sqrt{288}} \;,\; \frac{108}{2\,\sqrt{288}} \right)$
\item $g(x,y) = \ln\frac{x}{y} \enspace , \enspace D=[1,1,?]$\quad \dotfill\; ???\;\dotfill \quad $x - y - z + 1 = 0$
\item $\Delta x = {-0.3} \enspace,\enspace \Delta y = {0.01}$\quad \dotfill\; ???\;\dotfill \quad $-0.31$
\item $h(x,y) = -2x^2-2y^2+xy+3x$\quad \dotfill\; ???\;\dotfill \quad $\nicefrac{19}{15}$
\item $y = 4x+1$\quad \dotfill\; ???\;\dotfill \quad $\left[\frac{-1}{5}\;,\;\frac{1}{5}\right]$
\end{enumerate}
\end{varwidth}
\end{center}
\end{minipage}
\begin{minipage}{0.20\linewidth}
\begin{center}
{\Huge\bfseries 2.} \\[2mm]
\includegraphics[height=37mm]{images/braille.png}
{\small Braillove písmeno}
\end{center}
\end{minipage}
\end{center}
\end{minipage}
\\ \hdashline
\begin{minipage}[c][104.5mm][t]{0.5\linewidth}
\begin{center}
\vspace{7mm}
{\huge \LARGE Funkcia dvoch premenných, skupina \textit{Omega $\omega$} -\romannumeral3}\\[5mm]
\begin{minipage}{0.95\linewidth}
\begin{center}
{\small \begin{multicols}{2} \raggedright
\textbf{(a)} Vypočítajte zadanú parciálnu deriváciu funkcie \(f(x,y)\). \\[1mm]
\textbf{(b)} Určte gradient \(f(x,y)\) v bode \(A\). \\[1mm]
\textbf{(c)} Napíšte rovnicu dotykovej roviny ku grafu funkcie \(g(x,y)\) v dotykovom bode \(D\). \par \columnbreak
\textbf{(d)} Pomocou lineárnej aproximácie odhadnite funkčnú hodnotu funkcie \(g(x,y)\) v bode, ktorý je posunutý o vektor \((\Delta x, \Delta y)\) od bodu \(D_0\). \\[1mm]
\textbf{(e)} Vypočítajte \(h(x,y)\) v jej lokálnom extréme. \\[1mm]
\textbf{(f)} Nájdite bod viazeného lokálneho extrému funkcie \(h(x,y)\) na priamke. \\[1mm]
\end{multicols}
}
\end{center}
\end{minipage}
\\[4mm]
\begin{minipage}{0.79\linewidth}
\begin{center}
\begin{varwidth}{\linewidth}
\begin{enumerate}
\normalsizerr
\item $f(x,y) = \sqrt[3]{4x^2y + xy^2}\enspace , \enspace \pdv{f}{y}$\quad \dotfill\; ???\;\dotfill \quad $\frac{4x^2+2xy}{3\,\sqrt[3]{(4x^2y+xy^2)^{2}}}$
\item $A = [4, 3]$\quad \dotfill\; ???\;\dotfill \quad $\left( \frac{105}{3\,\sqrt[3]{228^2}} \;,\; \frac{88}{3\,\sqrt[3]{228^2}} \right)$
\item $g(x,y) = \ln\frac{x}{y} \enspace , \enspace D=[e^{2},1,?]$\quad \dotfill\; ???\;\dotfill \quad $\frac{x}{e^3} - y - z + 3 = 0$
\item $\Delta x = {0.2} \enspace,\enspace \Delta y = {0.02}$\quad \dotfill\; ???\;\dotfill \quad $1.827$
\item $h(x,y) = -7x^2-y^2-2xy+6x$\quad \dotfill\; ???\;\dotfill \quad $\nicefrac{3}{2}$
\item $y = -x-7$\quad \dotfill\; ???\;\dotfill \quad $\left[\frac{1}{2}\;,\;\frac{-15}{2}\right]$
\end{enumerate}
\end{varwidth}
\end{center}
\end{minipage}
\begin{minipage}{0.20\linewidth}
\begin{center}
{\Huge\bfseries 3.} \\[2mm]
\includegraphics[height=37mm]{images/braille.png}
{\small Braillove písmeno}
\end{center}
\end{minipage}
\end{center}
\end{minipage}
&
\begin{minipage}[c][104.5mm][t]{0.5\linewidth}
\begin{center}
\vspace{7mm}
{\huge \LARGE Funkcia dvoch premenných, skupina \textit{Omega $\omega$} -\romannumeral4}\\[5mm]
\begin{minipage}{0.95\linewidth}
\begin{center}
{\small \begin{multicols}{2} \raggedright
\textbf{(a)} Vypočítajte zadanú parciálnu deriváciu funkcie \(f(x,y)\). \\[1mm]
\textbf{(b)} Určte gradient \(f(x,y)\) v bode \(A\). \\[1mm]
\textbf{(c)} Napíšte rovnicu dotykovej roviny ku grafu funkcie \(g(x,y)\) v dotykovom bode \(D\). \par \columnbreak
\textbf{(d)} Pomocou lineárnej aproximácie odhadnite funkčnú hodnotu funkcie \(g(x,y)\) v bode, ktorý je posunutý o vektor \((\Delta x, \Delta y)\) od bodu \(D_0\). \\[1mm]
\textbf{(e)} Vypočítajte \(h(x,y)\) v jej lokálnom extréme. \\[1mm]
\textbf{(f)} Nájdite bod viazeného lokálneho extrému funkcie \(h(x,y)\) na priamke. \\[1mm]
\end{multicols}
}
\end{center}
\end{minipage}
\\[4mm]
\begin{minipage}{0.79\linewidth}
\begin{center}
\begin{varwidth}{\linewidth}
\begin{enumerate}
\normalsizerr
\item $f(x,y) = \sqrt[4]{2x^2y + 6xy^2}\enspace , \enspace \pdv{f}{x}$\quad \dotfill\; ???\;\dotfill \quad $\frac{4xy+6y^2}{4\,\sqrt[4]{(2x^2y+6xy^2)^{3}}}$
\item $A = [2, 2]$\quad \dotfill\; ???\;\dotfill \quad $\left( \frac{40}{4\,\sqrt[4]{64^3}} \;,\; \frac{112}{4\,\sqrt[4]{64^3}} \right)$
\item $g(x,y) = \sin\frac{x}{y} \enspace , \enspace D = [0, -1, ?]$\quad \dotfill\; ???\;\dotfill \quad $x+z=1$
\item $\Delta x = {0.02} \enspace,\enspace \Delta y = {-0.02}$\quad \dotfill\; ???\;\dotfill \quad $0.02$
\item $h(x,y) = -2x^2-3y^2-2xy+5x$\quad \dotfill\; ???\;\dotfill \quad $\nicefrac{19}{5}$
\item $y = x+4$\quad \dotfill\; ???\;\dotfill \quad $\left[\frac{-27}{14}\;,\;\frac{29}{6}\right]$
\end{enumerate}
\end{varwidth}
\end{center}
\end{minipage}
\begin{minipage}{0.20\linewidth}
\begin{center}
{\Huge\bfseries 4.} \\[2mm]
\includegraphics[height=37mm]{images/braille.png}
{\small Braillove písmeno}
\end{center}
\end{minipage}
\end{center}
\end{minipage}
%
\end{tabular}
\newpage
\begin{landscape}
\newgeometry{total={194mm,285mm}, left=8mm, top=9mm}
\begin{center}
{\huge \LARGE Funkcia dvoch premenných (riešenia)}\\[4mm]
\begin{varwidth}{\linewidth}
\begin{center}
\scriptsize
\rule[1mm]{\linewidth}{0.5pt}
$\boxed{\bm{\alpha}} \quad \begin{aligned}
\romannumeral1 : \; &\textbf{M} 
 &&\mathrm{\textbf{(a) }} \frac{8x^2+18xy}{3\,\sqrt[3]{(8x^2y+9xy^2)^{2}}}\,\text{\ding{51}}
 &&\mathrm{\textbf{(b) }} \left( \frac{25}{3\,\sqrt[3]{17^2}} \;,\; \frac{26}{3\,\sqrt[3]{17^2}} \right)\,\text{\ding{55}}
 &&\mathrm{\textbf{(c) }} x+3z=0\,\text{\ding{51}}
 &&\mathrm{\textbf{(d) }} -0.007\,\text{\ding{55}}
 &&\mathrm{\textbf{(e) }} \nicefrac{-5}{44}\,\text{\ding{55}}
 &&\mathrm{\textbf{(f) }} \left[\frac{41}{8}\;,\;\frac{-15}{8}\right]\,\text{\ding{51}}
\\[-0.4mm]
\romannumeral2 : \; &\textbf{A} 
 &&\mathrm{\textbf{(a) }} \frac{4xy+9y^2}{4\,\sqrt[4]{(2x^2y+9xy^2)^{3}}}\,\text{\ding{51}}
 &&\mathrm{\textbf{(b) }} \left( \frac{117}{4\,\sqrt[4]{297^3}} \;,\; \frac{180}{4\,\sqrt[4]{297^3}} \right)\,\text{\ding{55}}
 &&\mathrm{\textbf{(c) }} \frac{x}{e} - y - z + 1 = 0\,\text{\ding{55}}
 &&\mathrm{\textbf{(d) }} 0.976\,\text{\ding{55}}
 &&\mathrm{\textbf{(e) }} \nicefrac{4}{7}\,\text{\ding{55}}
 &&\mathrm{\textbf{(f) }} \left[\frac{-20}{11}\;,\;\frac{7}{11}\right]\,\text{\ding{55}}
\\[-0.4mm]
\romannumeral3 : \; &\textbf{P} 
 &&\mathrm{\textbf{(a) }} \frac{8xy+4y^2}{4\,\sqrt[4]{(4x^2y+4xy^2)^{3}}}\,\text{\ding{51}}
 &&\mathrm{\textbf{(b) }} \left( \frac{48}{4\,\sqrt[4]{64^3}} \;,\; \frac{48}{4\,\sqrt[4]{64^3}} \right)\,\text{\ding{51}}
 &&\mathrm{\textbf{(c) }} x+z=0\,\text{\ding{51}}
 &&\mathrm{\textbf{(d) }} -0.3\,\text{\ding{55}}
 &&\mathrm{\textbf{(e) }} \nicefrac{-32}{15}\,\text{\ding{55}}
 &&\mathrm{\textbf{(f) }} \left[\frac{-7}{6}\;,\;\frac{4}{3}\right]\,\text{\ding{51}}
\\[-0.4mm]
\romannumeral4 : \; &\textbf{A} 
 &&\mathrm{\textbf{(a) }} \frac{7x^2+6xy}{3\,\sqrt[3]{(7x^2y+3xy^2)^{2}}}\,\text{\ding{51}}
 &&\mathrm{\textbf{(b) }} \left( \frac{59}{3\,\sqrt[3]{124^2}} \;,\; \frac{136}{3\,\sqrt[3]{124^2}} \right)\,\text{\ding{55}}
 &&\mathrm{\textbf{(c) }} \frac{x}{e^3} - y - z + 3 = 0\,\text{\ding{55}}
 &&\mathrm{\textbf{(d) }} 2.685\,\text{\ding{55}}
 &&\mathrm{\textbf{(e) }} \nicefrac{4}{71}\,\text{\ding{55}}
 &&\mathrm{\textbf{(f) }} \left[\frac{4}{3}\;,\;\frac{5}{3}\right]\,\text{\ding{55}}
\end{aligned} $
\\[2mm]
\rule[1mm]{\linewidth}{0.5pt}
$\boxed{\bm{\beta}} \quad \begin{aligned}
\romannumeral1 : \; &\textbf{V} 
 &&\mathrm{\textbf{(a) }} \frac{x^2+10xy}{4\,\sqrt[4]{(x^2y+5xy^2)^{3}}}\,\text{\ding{51}}
 &&\mathrm{\textbf{(b) }} \left( \frac{69}{4\,\sqrt[4]{228^3}} \;,\; \frac{136}{4\,\sqrt[4]{228^3}} \right)\,\text{\ding{51}}
 &&\mathrm{\textbf{(c) }} \frac{x}{e^3} - y - z + 3 = 0\,\text{\ding{51}}
 &&\mathrm{\textbf{(d) }} 3.02\,\text{\ding{51}}
 &&\mathrm{\textbf{(e) }} \nicefrac{4}{63}\,\text{\ding{55}}
 &&\mathrm{\textbf{(f) }} \left[\frac{4}{7}\;,\;\frac{-3}{7}\right]\,\text{\ding{55}}
\\[-0.4mm]
\romannumeral2 : \; &\textbf{Á} 
 &&\mathrm{\textbf{(a) }} \frac{14xy+2y^2}{3\,\sqrt[3]{(7x^2y+2xy^2)^{2}}}\,\text{\ding{51}}
 &&\mathrm{\textbf{(b) }} \left( \frac{102}{3\,\sqrt[3]{120^2}} \;,\; \frac{52}{3\,\sqrt[3]{120^2}} \right)\,\text{\ding{55}}
 &&\mathrm{\textbf{(c) }} x+z=0\,\text{\ding{55}}
 &&\mathrm{\textbf{(d) }} -0.03\,\text{\ding{51}}
 &&\mathrm{\textbf{(e) }} \nicefrac{75}{8}\,\text{\ding{55}}
 &&\mathrm{\textbf{(f) }} \left[\frac{-47}{114}\;,\;\frac{20}{57}\right]\,\text{\ding{55}}
\\[-0.4mm]
\romannumeral3 : \; &\textbf{H} 
 &&\mathrm{\textbf{(a) }} \frac{5x^2+2xy}{2\,\sqrt{5x^2y+xy^2}}\,\text{\ding{51}}
 &&\mathrm{\textbf{(b) }} \left( \frac{11}{2\,\sqrt{6}} \;,\; \frac{7}{2\,\sqrt{6}} \right)\,\text{\ding{51}}
 &&\mathrm{\textbf{(c) }} \frac{x}{e^3} - y - z + 3 = 0\,\text{\ding{55}}
 &&\mathrm{\textbf{(d) }} 3.105\,\text{\ding{55}}
 &&\mathrm{\textbf{(e) }} \nicefrac{-3}{44}\,\text{\ding{51}}
 &&\mathrm{\textbf{(f) }} \left[\frac{-61}{74}\;,\;\frac{-39}{74}\right]\,\text{\ding{55}}
\\[-0.4mm]
\romannumeral4 : \; &\textbf{A} 
 &&\mathrm{\textbf{(a) }} \frac{6x^2+16xy}{2\,\sqrt{6x^2y+8xy^2}}\,\text{\ding{51}}
 &&\mathrm{\textbf{(b) }} \left( \frac{180}{2\,\sqrt{378}} \;,\; \frac{198}{2\,\sqrt{378}} \right)\,\text{\ding{55}}
 &&\mathrm{\textbf{(c) }} x-2z=0\,\text{\ding{55}}
 &&\mathrm{\textbf{(d) }} 0.15\,\text{\ding{55}}
 &&\mathrm{\textbf{(e) }} \nicefrac{-6}{47}\,\text{\ding{55}}
 &&\mathrm{\textbf{(f) }} \left[\frac{-59}{74}\;,\;\frac{15}{37}\right]\,\text{\ding{55}}
\end{aligned} $
\\[2mm]
\rule[1mm]{\linewidth}{0.5pt}
$\boxed{\bm{\gamma}} \quad \begin{aligned}
\romannumeral1 : \; &\textbf{V} 
 &&\mathrm{\textbf{(a) }} \frac{14xy+6y^2}{2\,\sqrt{7x^2y+6xy^2}}\,\text{\ding{51}}
 &&\mathrm{\textbf{(b) }} \left( \frac{138}{2\,\sqrt{192}} \;,\; \frac{100}{2\,\sqrt{192}} \right)\,\text{\ding{51}}
 &&\mathrm{\textbf{(c) }} x-3z=0\,\text{\ding{51}}
 &&\mathrm{\textbf{(d) }} -0.01\,\text{\ding{51}}
 &&\mathrm{\textbf{(e) }} \nicefrac{64}{15}\,\text{\ding{55}}
 &&\mathrm{\textbf{(f) }} \left[\frac{28}{15}\;,\;\frac{-4}{15}\right]\,\text{\ding{55}}
\\[-0.4mm]
\romannumeral2 : \; &\textbf{A} 
 &&\mathrm{\textbf{(a) }} \frac{4x^2+2xy}{4\,\sqrt[4]{(4x^2y+xy^2)^{3}}}\,\text{\ding{51}}
 &&\mathrm{\textbf{(b) }} \left( \frac{25}{4\,\sqrt[4]{39^3}} \;,\; \frac{42}{4\,\sqrt[4]{39^3}} \right)\,\text{\ding{55}}
 &&\mathrm{\textbf{(c) }} x-2z=0\,\text{\ding{55}}
 &&\mathrm{\textbf{(d) }} -0.15\,\text{\ding{55}}
 &&\mathrm{\textbf{(e) }} \nicefrac{-25}{3}\,\text{\ding{55}}
 &&\mathrm{\textbf{(f) }} \left[1\;,\;-6\right]\,\text{\ding{55}}
\\[-0.4mm]
\romannumeral3 : \; &\textbf{Ň} 
 &&\mathrm{\textbf{(a) }} \frac{3x^2+10xy}{4\,\sqrt[4]{(3x^2y+5xy^2)^{3}}}\,\text{\ding{51}}
 &&\mathrm{\textbf{(b) }} \left( \frac{56}{4\,\sqrt[4]{114^3}} \;,\; \frac{87}{4\,\sqrt[4]{114^3}} \right)\,\text{\ding{51}}
 &&\mathrm{\textbf{(c) }} x+z=0\,\text{\ding{55}}
 &&\mathrm{\textbf{(d) }} 0.03\,\text{\ding{51}}
 &&\mathrm{\textbf{(e) }} \nicefrac{80}{31}\,\text{\ding{55}}
 &&\mathrm{\textbf{(f) }} \left[\frac{-33}{140}\;,\;\frac{2}{35}\right]\,\text{\ding{51}}
\\[-0.4mm]
\romannumeral4 : \; &\textbf{A} 
 &&\mathrm{\textbf{(a) }} \frac{3x^2+10xy}{2\,\sqrt{3x^2y+5xy^2}}\,\text{\ding{51}}
 &&\mathrm{\textbf{(b) }} \left( \frac{117}{2\,\sqrt{324}} \;,\; \frac{168}{2\,\sqrt{324}} \right)\,\text{\ding{55}}
 &&\mathrm{\textbf{(c) }} x-z=0\,\text{\ding{55}}
 &&\mathrm{\textbf{(d) }} -0.1\,\text{\ding{55}}
 &&\mathrm{\textbf{(e) }} \nicefrac{-3}{2}\,\text{\ding{55}}
 &&\mathrm{\textbf{(f) }} \left[\frac{-12}{19}\;,\;\frac{-14}{19}\right]\,\text{\ding{55}}
\end{aligned} $
\\[2mm]
\rule[1mm]{\linewidth}{0.5pt}
$\boxed{\bm{\delta}} \quad \begin{aligned}
\romannumeral1 : \; &\textbf{E} 
 &&\mathrm{\textbf{(a) }} \frac{3x^2+8xy}{2\,\sqrt{3x^2y+4xy^2}}\,\text{\ding{51}}
 &&\mathrm{\textbf{(b) }} \left( \frac{22}{2\,\sqrt{39}} \;,\; \frac{51}{2\,\sqrt{39}} \right)\,\text{\ding{55}}
 &&\mathrm{\textbf{(c) }} x+z=0\,\text{\ding{55}}
 &&\mathrm{\textbf{(d) }} -0.03\,\text{\ding{55}}
 &&\mathrm{\textbf{(e) }} \nicefrac{-27}{23}\,\text{\ding{51}}
 &&\mathrm{\textbf{(f) }} \left[\frac{-49}{32}\;,\;\frac{-79}{32}\right]\,\text{\ding{55}}
\\[-0.4mm]
\romannumeral2 : \; &\textbf{P} 
 &&\mathrm{\textbf{(a) }} \frac{14xy+7y^2}{2\,\sqrt{7x^2y+7xy^2}}\,\text{\ding{51}}
 &&\mathrm{\textbf{(b) }} \left( \frac{63}{2\,\sqrt{140}} \;,\; \frac{168}{2\,\sqrt{140}} \right)\,\text{\ding{51}}
 &&\mathrm{\textbf{(c) }} \frac{x}{e} - y - z + 1 = 0\,\text{\ding{51}}
 &&\mathrm{\textbf{(d) }} 0.92\,\text{\ding{55}}
 &&\mathrm{\textbf{(e) }} \nicefrac{-36}{11}\,\text{\ding{55}}
 &&\mathrm{\textbf{(f) }} \left[\frac{41}{30}\;,\;\frac{-9}{10}\right]\,\text{\ding{51}}
\\[-0.4mm]
\romannumeral3 : \; &\textbf{O} 
 &&\mathrm{\textbf{(a) }} \frac{12xy+5y^2}{3\,\sqrt[3]{(6x^2y+5xy^2)^{2}}}\,\text{\ding{51}}
 &&\mathrm{\textbf{(b) }} \left( \frac{116}{3\,\sqrt[3]{272^2}} \;,\; \frac{176}{3\,\sqrt[3]{272^2}} \right)\,\text{\ding{55}}
 &&\mathrm{\textbf{(c) }} x+z=0\,\text{\ding{51}}
 &&\mathrm{\textbf{(d) }} -0.3\,\text{\ding{55}}
 &&\mathrm{\textbf{(e) }} \nicefrac{-100}{31}\,\text{\ding{51}}
 &&\mathrm{\textbf{(f) }} \left[\frac{95}{32}\;,\;\frac{1}{16}\right]\,\text{\ding{55}}
\\[-0.4mm]
\romannumeral4 : \; &\textbf{S} 
 &&\mathrm{\textbf{(a) }} \frac{7x^2+4xy}{2\,\sqrt{7x^2y+2xy^2}}\,\text{\ding{55}}
 &&\mathrm{\textbf{(b) }} \left( \frac{120}{2\,\sqrt{256}} \;,\; \frac{144}{2\,\sqrt{256}} \right)\,\text{\ding{51}}
 &&\mathrm{\textbf{(c) }} \frac{x}{e^4} - y - z + 4 = 0\,\text{\ding{51}}
 &&\mathrm{\textbf{(d) }} 3.695\,\text{\ding{55}}
 &&\mathrm{\textbf{(e) }} \nicefrac{100}{87}\,\text{\ding{55}}
 &&\mathrm{\textbf{(f) }} \left[\frac{3}{2}\;,\;\frac{-5}{2}\right]\,\text{\ding{51}}
\end{aligned} $
\\[2mm]
\rule[1mm]{\linewidth}{0.5pt}
$\boxed{\bm{\epsilon}} \quad \begin{aligned}
\romannumeral1 : \; &\textbf{D} 
 &&\mathrm{\textbf{(a) }} \frac{9x^2+18xy}{3\,\sqrt[3]{(9x^2y+9xy^2)^{2}}}\,\text{\ding{51}}
 &&\mathrm{\textbf{(b) }} \left( \frac{72}{3\,\sqrt[3]{54^2}} \;,\; \frac{45}{3\,\sqrt[3]{54^2}} \right)\,\text{\ding{55}}
 &&\mathrm{\textbf{(c) }} x - y - z + 0 = 0\,\text{\ding{55}}
 &&\mathrm{\textbf{(d) }} 0.0\,\text{\ding{55}}
 &&\mathrm{\textbf{(e) }} \nicefrac{-63}{80}\,\text{\ding{51}}
 &&\mathrm{\textbf{(f) }} \left[\frac{-47}{144}\;,\;\frac{-1}{48}\right]\,\text{\ding{51}}
\\[-0.4mm]
\romannumeral2 : \; &\textbf{R} 
 &&\mathrm{\textbf{(a) }} \frac{x^2+2xy}{2\,\sqrt{x^2y+xy^2}}\,\text{\ding{51}}
 &&\mathrm{\textbf{(b) }} \left( \frac{32}{2\,\sqrt{48}} \;,\; \frac{20}{2\,\sqrt{48}} \right)\,\text{\ding{51}}
 &&\mathrm{\textbf{(c) }} x+2z=0\,\text{\ding{51}}
 &&\mathrm{\textbf{(d) }} -0.015\,\text{\ding{55}}
 &&\mathrm{\textbf{(e) }} \nicefrac{-12}{5}\,\text{\ding{51}}
 &&\mathrm{\textbf{(f) }} \left[\frac{-13}{4}\;,\;\frac{3}{4}\right]\,\text{\ding{55}}
\\[-0.4mm]
\romannumeral3 : \; &\textbf{A} 
 &&\mathrm{\textbf{(a) }} \frac{x^2+8xy}{4\,\sqrt[4]{(x^2y+4xy^2)^{3}}}\,\text{\ding{51}}
 &&\mathrm{\textbf{(b) }} \left( \frac{54}{4\,\sqrt[4]{135^3}} \;,\; \frac{81}{4\,\sqrt[4]{135^3}} \right)\,\text{\ding{55}}
 &&\mathrm{\textbf{(c) }} \frac{x}{e} - y - z + 1 = 0\,\text{\ding{55}}
 &&\mathrm{\textbf{(d) }} 0.984\,\text{\ding{55}}
 &&\mathrm{\textbf{(e) }} \nicefrac{-1}{7}\,\text{\ding{55}}
 &&\mathrm{\textbf{(f) }} \left[\frac{-37}{8}\;,\;\frac{-11}{8}\right]\,\text{\ding{55}}
\\[-0.4mm]
\romannumeral4 : \; &\textbf{K} 
 &&\mathrm{\textbf{(a) }} \frac{6xy+4y^2}{3\,\sqrt[3]{(3x^2y+4xy^2)^{2}}}\,\text{\ding{51}}
 &&\mathrm{\textbf{(b) }} \left( \frac{88}{3\,\sqrt[3]{76^2}} \;,\; \frac{35}{3\,\sqrt[3]{76^2}} \right)\,\text{\ding{55}}
 &&\mathrm{\textbf{(c) }} \frac{x}{e} - y - z + 1 = 0\,\text{\ding{51}}
 &&\mathrm{\textbf{(d) }} 0.933\,\text{\ding{55}}
 &&\mathrm{\textbf{(e) }} \nicefrac{-1}{5}\,\text{\ding{55}}
 &&\mathrm{\textbf{(f) }} \left[\frac{-11}{38}\;,\;\frac{5}{38}\right]\,\text{\ding{55}}
\end{aligned} $
\\[2mm]
\rule[1mm]{\linewidth}{0.5pt}
$\boxed{\bm{\zeta}} \quad \begin{aligned}
\romannumeral1 : \; &\textbf{M} 
 &&\mathrm{\textbf{(a) }} \frac{6xy+7y^2}{4\,\sqrt[4]{(3x^2y+7xy^2)^{3}}}\,\text{\ding{51}}
 &&\mathrm{\textbf{(b) }} \left( \frac{31}{4\,\sqrt[4]{76^3}} \;,\; \frac{104}{4\,\sqrt[4]{76^3}} \right)\,\text{\ding{55}}
 &&\mathrm{\textbf{(c) }} x+3z=0\,\text{\ding{51}}
 &&\mathrm{\textbf{(d) }} -0.1\,\text{\ding{55}}
 &&\mathrm{\textbf{(e) }} \nicefrac{4}{3}\,\text{\ding{55}}
 &&\mathrm{\textbf{(f) }} \left[\frac{-13}{62}\;,\;\frac{5}{31}\right]\,\text{\ding{51}}
\\[-0.4mm]
\romannumeral2 : \; &\textbf{R} 
 &&\mathrm{\textbf{(a) }} \frac{6x^2+4xy}{2\,\sqrt{6x^2y+2xy^2}}\,\text{\ding{51}}
 &&\mathrm{\textbf{(b) }} \left( \frac{14}{2\,\sqrt{8}} \;,\; \frac{10}{2\,\sqrt{8}} \right)\,\text{\ding{51}}
 &&\mathrm{\textbf{(c) }} x+z=0\,\text{\ding{51}}
 &&\mathrm{\textbf{(d) }} -0.02\,\text{\ding{55}}
 &&\mathrm{\textbf{(e) }} \nicefrac{-9}{5}\,\text{\ding{51}}
 &&\mathrm{\textbf{(f) }} \left[\frac{3}{2}\;,\;\frac{-1}{2}\right]\,\text{\ding{55}}
\\[-0.4mm]
\romannumeral3 : \; &\textbf{A} 
 &&\mathrm{\textbf{(a) }} \frac{4x^2+6xy}{4\,\sqrt[4]{(4x^2y+3xy^2)^{3}}}\,\text{\ding{51}}
 &&\mathrm{\textbf{(b) }} \left( \frac{11}{4\,\sqrt[4]{7^3}} \;,\; \frac{10}{4\,\sqrt[4]{7^3}} \right)\,\text{\ding{55}}
 &&\mathrm{\textbf{(c) }} x+3z=0\,\text{\ding{55}}
 &&\mathrm{\textbf{(d) }} -0.007\,\text{\ding{55}}
 &&\mathrm{\textbf{(e) }} \nicefrac{-4}{7}\,\text{\ding{55}}
 &&\mathrm{\textbf{(f) }} \left[\frac{-5}{2}\;,\;\frac{11}{2}\right]\,\text{\ding{55}}
\\[-0.4mm]
\romannumeral4 : \; &\textbf{K} 
 &&\mathrm{\textbf{(a) }} \frac{6xy+6y^2}{4\,\sqrt[4]{(3x^2y+6xy^2)^{3}}}\,\text{\ding{51}}
 &&\mathrm{\textbf{(b) }} \left( \frac{30}{4\,\sqrt[4]{72^3}} \;,\; \frac{96}{4\,\sqrt[4]{72^3}} \right)\,\text{\ding{55}}
 &&\mathrm{\textbf{(c) }} \frac{x}{e} - y - z + 1 = 0\,\text{\ding{51}}
 &&\mathrm{\textbf{(d) }} 1.067\,\text{\ding{55}}
 &&\mathrm{\textbf{(e) }} \nicefrac{-8}{7}\,\text{\ding{55}}
 &&\mathrm{\textbf{(f) }} \left[\frac{-1}{8}\;,\;\frac{-7}{8}\right]\,\text{\ding{55}}
\end{aligned} $
\\[2mm]
\rule[1mm]{\linewidth}{0.5pt}
$\boxed{\bm{\eta}} \quad \begin{aligned}
\romannumeral1 : \; &\textbf{N} 
 &&\mathrm{\textbf{(a) }} \frac{3x^2+12xy}{2\,\sqrt{3x^2y+6xy^2}}\,\text{\ding{51}}
 &&\mathrm{\textbf{(b) }} \left( \frac{192}{2\,\sqrt{576}} \;,\; \frac{240}{2\,\sqrt{576}} \right)\,\text{\ding{55}}
 &&\mathrm{\textbf{(c) }} \frac{x}{e^4} - y - z + 4 = 0\,\text{\ding{51}}
 &&\mathrm{\textbf{(d) }} 4.195\,\text{\ding{55}}
 &&\mathrm{\textbf{(e) }} \nicefrac{-16}{23}\,\text{\ding{51}}
 &&\mathrm{\textbf{(f) }} \left[\frac{-2}{3}\;,\;\frac{4}{3}\right]\,\text{\ding{51}}
\\[-0.4mm]
\romannumeral2 : \; &\textbf{O} 
 &&\mathrm{\textbf{(a) }} \frac{5x^2+10xy}{2\,\sqrt{5x^2y+5xy^2}}\,\text{\ding{51}}
 &&\mathrm{\textbf{(b) }} \left( \frac{80}{2\,\sqrt{150}} \;,\; \frac{105}{2\,\sqrt{150}} \right)\,\text{\ding{55}}
 &&\mathrm{\textbf{(c) }} x+3z=0\,\text{\ding{51}}
 &&\mathrm{\textbf{(d) }} 0.01\,\text{\ding{55}}
 &&\mathrm{\textbf{(e) }} \nicefrac{-18}{23}\,\text{\ding{51}}
 &&\mathrm{\textbf{(f) }} \left[-1\;,\;-1\right]\,\text{\ding{55}}
\\[-0.4mm]
\romannumeral3 : \; &\textbf{H} 
 &&\mathrm{\textbf{(a) }} \frac{14xy+3y^2}{3\,\sqrt[3]{(7x^2y+3xy^2)^{2}}}\,\text{\ding{51}}
 &&\mathrm{\textbf{(b) }} \left( \frac{272}{3\,\sqrt[3]{640^2}} \;,\; \frac{208}{3\,\sqrt[3]{640^2}} \right)\,\text{\ding{51}}
 &&\mathrm{\textbf{(c) }} x+3z=0\,\text{\ding{55}}
 &&\mathrm{\textbf{(d) }} 0.067\,\text{\ding{55}}
 &&\mathrm{\textbf{(e) }} \nicefrac{5}{111}\,\text{\ding{51}}
 &&\mathrm{\textbf{(f) }} \left[\frac{-11}{32}\;,\;\frac{5}{16}\right]\,\text{\ding{55}}
\\[-0.4mm]
\romannumeral4 : \; &\textbf{A} 
 &&\mathrm{\textbf{(a) }} \frac{12xy+y^2}{3\,\sqrt[3]{(6x^2y+xy^2)^{2}}}\,\text{\ding{51}}
 &&\mathrm{\textbf{(b) }} \left( \frac{13}{3\,\sqrt[3]{7^2}} \;,\; \frac{8}{3\,\sqrt[3]{7^2}} \right)\,\text{\ding{55}}
 &&\mathrm{\textbf{(c) }} x-z=0\,\text{\ding{55}}
 &&\mathrm{\textbf{(d) }} -0.03\,\text{\ding{55}}
 &&\mathrm{\textbf{(e) }} \nicefrac{2}{11}\,\text{\ding{55}}
 &&\mathrm{\textbf{(f) }} \left[2\;,\;3\right]\,\text{\ding{55}}
\end{aligned} $
\\[2mm]
\rule[1mm]{\linewidth}{0.5pt}
\end{center}\end{varwidth}\end{center}\clearpage\begin{center}{\huge \LARGE Funkcia dvoch premenných (riešenia)}\\[4mm]\begin{varwidth}{\linewidth}\begin{center}
\scriptsize
\rule[1mm]{\linewidth}{0.5pt}
$\boxed{\bm{\theta}} \quad \begin{aligned}
\romannumeral1 : \; &\textbf{V} 
 &&\mathrm{\textbf{(a) }} \frac{4xy+6y^2}{4\,\sqrt[4]{(2x^2y+6xy^2)^{3}}}\,\text{\ding{51}}
 &&\mathrm{\textbf{(b) }} \left( \frac{56}{4\,\sqrt[4]{160^3}} \;,\; \frac{128}{4\,\sqrt[4]{160^3}} \right)\,\text{\ding{51}}
 &&\mathrm{\textbf{(c) }} \frac{x}{e} - y - z + 1 = 0\,\text{\ding{51}}
 &&\mathrm{\textbf{(d) }} 0.763\,\text{\ding{51}}
 &&\mathrm{\textbf{(e) }} \nicefrac{2}{21}\,\text{\ding{55}}
 &&\mathrm{\textbf{(f) }} \left[\frac{-37}{12}\;,\;\frac{11}{12}\right]\,\text{\ding{55}}
\\[-0.4mm]
\romannumeral2 : \; &\textbf{L} 
 &&\mathrm{\textbf{(a) }} \frac{4x^2+12xy}{4\,\sqrt[4]{(4x^2y+6xy^2)^{3}}}\,\text{\ding{51}}
 &&\mathrm{\textbf{(b) }} \left( \frac{40}{4\,\sqrt[4]{32^3}} \;,\; \frac{28}{4\,\sqrt[4]{32^3}} \right)\,\text{\ding{51}}
 &&\mathrm{\textbf{(c) }} \frac{x}{e^3} - y - z + 3 = 0\,\text{\ding{51}}
 &&\mathrm{\textbf{(d) }} 2.965\,\text{\ding{55}}
 &&\mathrm{\textbf{(e) }} \nicefrac{4}{11}\,\text{\ding{55}}
 &&\mathrm{\textbf{(f) }} \left[\frac{11}{6}\;,\;\frac{-16}{3}\right]\,\text{\ding{55}}
\\[-0.4mm]
\romannumeral3 : \; &\textbf{E} 
 &&\mathrm{\textbf{(a) }} \frac{3x^2+10xy}{2\,\sqrt{3x^2y+5xy^2}}\,\text{\ding{51}}
 &&\mathrm{\textbf{(b) }} \left( \frac{23}{2\,\sqrt{42}} \;,\; \frac{57}{2\,\sqrt{42}} \right)\,\text{\ding{55}}
 &&\mathrm{\textbf{(c) }} x+2z=0\,\text{\ding{55}}
 &&\mathrm{\textbf{(d) }} -0.01\,\text{\ding{55}}
 &&\mathrm{\textbf{(e) }} \nicefrac{-5}{4}\,\text{\ding{51}}
 &&\mathrm{\textbf{(f) }} \left[\frac{31}{4}\;,\;\frac{1}{4}\right]\,\text{\ding{55}}
\\[-0.4mm]
\romannumeral4 : \; &\textbf{K} 
 &&\mathrm{\textbf{(a) }} \frac{2xy+3y^2}{4\,\sqrt[4]{(x^2y+3xy^2)^{3}}}\,\text{\ding{51}}
 &&\mathrm{\textbf{(b) }} \left( \frac{80}{4\,\sqrt[4]{256^3}} \;,\; \frac{112}{4\,\sqrt[4]{256^3}} \right)\,\text{\ding{55}}
 &&\mathrm{\textbf{(c) }} x+z=0\,\text{\ding{51}}
 &&\mathrm{\textbf{(d) }} -0.03\,\text{\ding{55}}
 &&\mathrm{\textbf{(e) }} \nicefrac{-4}{17}\,\text{\ding{55}}
 &&\mathrm{\textbf{(f) }} \left[\frac{43}{46}\;,\;\frac{26}{23}\right]\,\text{\ding{55}}
\end{aligned} $
\\[2mm]
\rule[1mm]{\linewidth}{0.5pt}
$\boxed{\bm{\iota}} \quad \begin{aligned}
\romannumeral1 : \; &\textbf{E} 
 &&\mathrm{\textbf{(a) }} \frac{3x^2+4xy}{2\,\sqrt{3x^2y+2xy^2}}\,\text{\ding{51}}
 &&\mathrm{\textbf{(b) }} \left( \frac{20}{2\,\sqrt{14}} \;,\; \frac{11}{2\,\sqrt{14}} \right)\,\text{\ding{55}}
 &&\mathrm{\textbf{(c) }} x - y - z + 0 = 0\,\text{\ding{55}}
 &&\mathrm{\textbf{(d) }} -0.3\,\text{\ding{55}}
 &&\mathrm{\textbf{(e) }} \nicefrac{12}{5}\,\text{\ding{51}}
 &&\mathrm{\textbf{(f) }} \left[\frac{-13}{9}\;,\;\frac{10}{9}\right]\,\text{\ding{55}}
\\[-0.4mm]
\romannumeral2 : \; &\textbf{R} 
 &&\mathrm{\textbf{(a) }} \frac{12xy+5y^2}{2\,\sqrt{6x^2y+5xy^2}}\,\text{\ding{51}}
 &&\mathrm{\textbf{(b) }} \left( \frac{81}{2\,\sqrt{63}} \;,\; \frac{36}{2\,\sqrt{63}} \right)\,\text{\ding{51}}
 &&\mathrm{\textbf{(c) }} \frac{x}{e^4} - y - z + 4 = 0\,\text{\ding{51}}
 &&\mathrm{\textbf{(d) }} 3.796\,\text{\ding{55}}
 &&\mathrm{\textbf{(e) }} \nicefrac{3}{5}\,\text{\ding{51}}
 &&\mathrm{\textbf{(f) }} \left[\frac{-57}{20}\;,\;\frac{63}{20}\right]\,\text{\ding{55}}
\\[-0.4mm]
\romannumeral3 : \; &\textbf{I} 
 &&\mathrm{\textbf{(a) }} \frac{4xy+5y^2}{3\,\sqrt[3]{(2x^2y+5xy^2)^{2}}}\,\text{\ding{55}}
 &&\mathrm{\textbf{(b) }} \left( \frac{96}{3\,\sqrt[3]{88^2}} \;,\; \frac{42}{3\,\sqrt[3]{88^2}} \right)\,\text{\ding{51}}
 &&\mathrm{\textbf{(c) }} x+z=0\,\text{\ding{55}}
 &&\mathrm{\textbf{(d) }} 0.03\,\text{\ding{55}}
 &&\mathrm{\textbf{(e) }} \nicefrac{-28}{47}\,\text{\ding{55}}
 &&\mathrm{\textbf{(f) }} \left[\frac{29}{8}\;,\;\frac{11}{8}\right]\,\text{\ding{51}}
\\[-0.4mm]
\romannumeral4 : \; &\textbf{K} 
 &&\mathrm{\textbf{(a) }} \frac{2xy+7y^2}{3\,\sqrt[3]{(x^2y+7xy^2)^{2}}}\,\text{\ding{51}}
 &&\mathrm{\textbf{(b) }} \left( \frac{40}{3\,\sqrt[3]{102^2}} \;,\; \frac{93}{3\,\sqrt[3]{102^2}} \right)\,\text{\ding{55}}
 &&\mathrm{\textbf{(c) }} x+2z=0\,\text{\ding{51}}
 &&\mathrm{\textbf{(d) }} -0.1\,\text{\ding{55}}
 &&\mathrm{\textbf{(e) }} \nicefrac{9}{5}\,\text{\ding{55}}
 &&\mathrm{\textbf{(f) }} \left[\frac{9}{46}\;,\;\frac{-1}{46}\right]\,\text{\ding{55}}
\end{aligned} $
\\[2mm]
\rule[1mm]{\linewidth}{0.5pt}
$\boxed{\bm{\kappa}} \quad \begin{aligned}
\romannumeral1 : \; &\textbf{K} 
 &&\mathrm{\textbf{(a) }} \frac{2xy+5y^2}{2\,\sqrt{x^2y+5xy^2}}\,\text{\ding{51}}
 &&\mathrm{\textbf{(b) }} \left( \frac{36}{2\,\sqrt{112}} \;,\; \frac{96}{2\,\sqrt{112}} \right)\,\text{\ding{55}}
 &&\mathrm{\textbf{(c) }} x+2z=0\,\text{\ding{51}}
 &&\mathrm{\textbf{(d) }} 0.05\,\text{\ding{55}}
 &&\mathrm{\textbf{(e) }} \nicefrac{16}{23}\,\text{\ding{55}}
 &&\mathrm{\textbf{(f) }} \left[\frac{23}{58}\;,\;\frac{11}{58}\right]\,\text{\ding{55}}
\\[-0.4mm]
\romannumeral2 : \; &\textbf{O} 
 &&\mathrm{\textbf{(a) }} \frac{6x^2+6xy}{2\,\sqrt{6x^2y+3xy^2}}\,\text{\ding{51}}
 &&\mathrm{\textbf{(b) }} \left( \frac{240}{2\,\sqrt{576}} \;,\; \frac{192}{2\,\sqrt{576}} \right)\,\text{\ding{55}}
 &&\mathrm{\textbf{(c) }} x+z=0\,\text{\ding{51}}
 &&\mathrm{\textbf{(d) }} 0.3\,\text{\ding{55}}
 &&\mathrm{\textbf{(e) }} \nicefrac{-9}{11}\,\text{\ding{51}}
 &&\mathrm{\textbf{(f) }} \left[\frac{5}{6}\;,\;\frac{2}{3}\right]\,\text{\ding{55}}
\\[-0.4mm]
\romannumeral3 : \; &\textbf{S} 
 &&\mathrm{\textbf{(a) }} \frac{4x^2+6xy}{2\,\sqrt{4x^2y+3xy^2}}\,\text{\ding{55}}
 &&\mathrm{\textbf{(b) }} \left( \frac{27}{2\,\sqrt{45}} \;,\; \frac{54}{2\,\sqrt{45}} \right)\,\text{\ding{51}}
 &&\mathrm{\textbf{(c) }} x-3z=0\,\text{\ding{51}}
 &&\mathrm{\textbf{(d) }} 0.007\,\text{\ding{55}}
 &&\mathrm{\textbf{(e) }} \nicefrac{20}{111}\,\text{\ding{55}}
 &&\mathrm{\textbf{(f) }} \left[\frac{45}{32}\;,\;\frac{19}{16}\right]\,\text{\ding{51}}
\\[-0.4mm]
\romannumeral4 : \; &\textbf{Ť} 
 &&\mathrm{\textbf{(a) }} \frac{3x^2+6xy}{4\,\sqrt[4]{(3x^2y+3xy^2)^{3}}}\,\text{\ding{51}}
 &&\mathrm{\textbf{(b) }} \left( \frac{48}{4\,\sqrt[4]{90^3}} \;,\; \frac{63}{4\,\sqrt[4]{90^3}} \right)\,\text{\ding{51}}
 &&\mathrm{\textbf{(c) }} x+3z=0\,\text{\ding{55}}
 &&\mathrm{\textbf{(d) }} -0.033\,\text{\ding{51}}
 &&\mathrm{\textbf{(e) }} \nicefrac{-32}{31}\,\text{\ding{51}}
 &&\mathrm{\textbf{(f) }} \left[\frac{-1}{2}\;,\;\frac{3}{2}\right]\,\text{\ding{55}}
\end{aligned} $
\\[2mm]
\rule[1mm]{\linewidth}{0.5pt}
$\boxed{\bm{\lambda}} \quad \begin{aligned}
\romannumeral1 : \; &\textbf{I} 
 &&\mathrm{\textbf{(a) }} \frac{x^2+4xy}{3\,\sqrt[3]{(x^2y+2xy^2)^{2}}}\,\text{\ding{55}}
 &&\mathrm{\textbf{(b) }} \left( \frac{16}{3\,\sqrt[3]{24^2}} \;,\; \frac{20}{3\,\sqrt[3]{24^2}} \right)\,\text{\ding{51}}
 &&\mathrm{\textbf{(c) }} x - y - z + 0 = 0\,\text{\ding{55}}
 &&\mathrm{\textbf{(d) }} -0.3\,\text{\ding{55}}
 &&\mathrm{\textbf{(e) }} \nicefrac{18}{7}\,\text{\ding{55}}
 &&\mathrm{\textbf{(f) }} \left[\frac{9}{8}\;,\;\frac{3}{4}\right]\,\text{\ding{51}}
\\[-0.4mm]
\romannumeral2 : \; &\textbf{V} 
 &&\mathrm{\textbf{(a) }} \frac{2x^2+14xy}{3\,\sqrt[3]{(2x^2y+7xy^2)^{2}}}\,\text{\ding{51}}
 &&\mathrm{\textbf{(b) }} \left( \frac{111}{3\,\sqrt[3]{348^2}} \;,\; \frac{200}{3\,\sqrt[3]{348^2}} \right)\,\text{\ding{51}}
 &&\mathrm{\textbf{(c) }} x-3z=0\,\text{\ding{51}}
 &&\mathrm{\textbf{(d) }} -0.067\,\text{\ding{51}}
 &&\mathrm{\textbf{(e) }} \nicefrac{4}{7}\,\text{\ding{55}}
 &&\mathrm{\textbf{(f) }} \left[\frac{9}{28}\;,\;\frac{-5}{14}\right]\,\text{\ding{55}}
\\[-0.4mm]
\romannumeral3 : \; &\textbf{A} 
 &&\mathrm{\textbf{(a) }} \frac{14xy+5y^2}{3\,\sqrt[3]{(7x^2y+5xy^2)^{2}}}\,\text{\ding{51}}
 &&\mathrm{\textbf{(b) }} \left( \frac{33}{3\,\sqrt[3]{38^2}} \;,\; \frac{48}{3\,\sqrt[3]{38^2}} \right)\,\text{\ding{55}}
 &&\mathrm{\textbf{(c) }} x-z=0\,\text{\ding{55}}
 &&\mathrm{\textbf{(d) }} 0.2\,\text{\ding{55}}
 &&\mathrm{\textbf{(e) }} \nicefrac{4}{3}\,\text{\ding{55}}
 &&\mathrm{\textbf{(f) }} \left[\frac{6}{11}\;,\;\frac{16}{11}\right]\,\text{\ding{55}}
\\[-0.4mm]
\romannumeral4 : \; &\textbf{N} 
 &&\mathrm{\textbf{(a) }} \frac{2x^2+12xy}{2\,\sqrt{2x^2y+6xy^2}}\,\text{\ding{51}}
 &&\mathrm{\textbf{(b) }} \left( \frac{18}{2\,\sqrt{36}} \;,\; \frac{54}{2\,\sqrt{36}} \right)\,\text{\ding{55}}
 &&\mathrm{\textbf{(c) }} x - y - z + 0 = 0\,\text{\ding{51}}
 &&\mathrm{\textbf{(d) }} 0.31\,\text{\ding{55}}
 &&\mathrm{\textbf{(e) }} \nicefrac{-1}{24}\,\text{\ding{51}}
 &&\mathrm{\textbf{(f) }} \left[\frac{-21}{62}\;,\;\frac{10}{31}\right]\,\text{\ding{51}}
\end{aligned} $
\\[2mm]
\rule[1mm]{\linewidth}{0.5pt}
$\boxed{\bm{\mu}} \quad \begin{aligned}
\romannumeral1 : \; &\textbf{P} 
 &&\mathrm{\textbf{(a) }} \frac{14xy+2y^2}{2\,\sqrt{7x^2y+2xy^2}}\,\text{\ding{51}}
 &&\mathrm{\textbf{(b) }} \left( \frac{120}{2\,\sqrt{256}} \;,\; \frac{144}{2\,\sqrt{256}} \right)\,\text{\ding{51}}
 &&\mathrm{\textbf{(c) }} x+z=0\,\text{\ding{51}}
 &&\mathrm{\textbf{(d) }} -0.2\,\text{\ding{55}}
 &&\mathrm{\textbf{(e) }} \nicefrac{9}{10}\,\text{\ding{55}}
 &&\mathrm{\textbf{(f) }} \left[\frac{-13}{18}\;,\;\frac{5}{6}\right]\,\text{\ding{51}}
\\[-0.4mm]
\romannumeral2 : \; &\textbf{I} 
 &&\mathrm{\textbf{(a) }} \frac{12xy+y^2}{4\,\sqrt[4]{(6x^2y+xy^2)^{3}}}\,\text{\ding{55}}
 &&\mathrm{\textbf{(b) }} \left( \frac{25}{4\,\sqrt[4]{26^3}} \;,\; \frac{28}{4\,\sqrt[4]{26^3}} \right)\,\text{\ding{51}}
 &&\mathrm{\textbf{(c) }} \frac{x}{e^5} - y - z + 5 = 0\,\text{\ding{55}}
 &&\mathrm{\textbf{(d) }} 4.969\,\text{\ding{55}}
 &&\mathrm{\textbf{(e) }} \nicefrac{9}{19}\,\text{\ding{55}}
 &&\mathrm{\textbf{(f) }} \left[\frac{9}{14}\;,\;\frac{-5}{14}\right]\,\text{\ding{51}}
\\[-0.4mm]
\romannumeral3 : \; &\textbf{V} 
 &&\mathrm{\textbf{(a) }} \frac{2xy+4y^2}{4\,\sqrt[4]{(x^2y+4xy^2)^{3}}}\,\text{\ding{51}}
 &&\mathrm{\textbf{(b) }} \left( \frac{24}{4\,\sqrt[4]{40^3}} \;,\; \frac{36}{4\,\sqrt[4]{40^3}} \right)\,\text{\ding{51}}
 &&\mathrm{\textbf{(c) }} x-z=0\,\text{\ding{51}}
 &&\mathrm{\textbf{(d) }} 0.01\,\text{\ding{51}}
 &&\mathrm{\textbf{(e) }} \nicefrac{75}{8}\,\text{\ding{55}}
 &&\mathrm{\textbf{(f) }} \left[\frac{-5}{6}\;,\;\frac{2}{3}\right]\,\text{\ding{55}}
\\[-0.4mm]
\romannumeral4 : \; &\textbf{O} 
 &&\mathrm{\textbf{(a) }} \frac{5x^2+8xy}{2\,\sqrt{5x^2y+4xy^2}}\,\text{\ding{51}}
 &&\mathrm{\textbf{(b) }} \left( \frac{96}{2\,\sqrt{224}} \;,\; \frac{144}{2\,\sqrt{224}} \right)\,\text{\ding{55}}
 &&\mathrm{\textbf{(c) }} x+2z=0\,\text{\ding{51}}
 &&\mathrm{\textbf{(d) }} -0.15\,\text{\ding{55}}
 &&\mathrm{\textbf{(e) }} \nicefrac{-5}{4}\,\text{\ding{51}}
 &&\mathrm{\textbf{(f) }} \left[\frac{67}{106}\;,\;\frac{11}{106}\right]\,\text{\ding{55}}
\end{aligned} $
\\[2mm]
\rule[1mm]{\linewidth}{0.5pt}
$\boxed{\bm{\nu}} \quad \begin{aligned}
\romannumeral1 : \; &\textbf{P} 
 &&\mathrm{\textbf{(a) }} \frac{10xy+2y^2}{2\,\sqrt{5x^2y+2xy^2}}\,\text{\ding{51}}
 &&\mathrm{\textbf{(b) }} \left( \frac{88}{2\,\sqrt{192}} \;,\; \frac{112}{2\,\sqrt{192}} \right)\,\text{\ding{51}}
 &&\mathrm{\textbf{(c) }} x-3z=0\,\text{\ding{51}}
 &&\mathrm{\textbf{(d) }} 0.01\,\text{\ding{55}}
 &&\mathrm{\textbf{(e) }} \nicefrac{15}{8}\,\text{\ding{55}}
 &&\mathrm{\textbf{(f) }} \left[\frac{-31}{20}\;,\;\frac{-9}{10}\right]\,\text{\ding{51}}
\\[-0.4mm]
\romannumeral2 : \; &\textbf{R} 
 &&\mathrm{\textbf{(a) }} \frac{12xy+4y^2}{4\,\sqrt[4]{(6x^2y+4xy^2)^{3}}}\,\text{\ding{51}}
 &&\mathrm{\textbf{(b) }} \left( \frac{180}{4\,\sqrt[4]{432^3}} \;,\; \frac{192}{4\,\sqrt[4]{432^3}} \right)\,\text{\ding{51}}
 &&\mathrm{\textbf{(c) }} x+2z=0\,\text{\ding{51}}
 &&\mathrm{\textbf{(d) }} 0.005\,\text{\ding{55}}
 &&\mathrm{\textbf{(e) }} \nicefrac{-16}{63}\,\text{\ding{51}}
 &&\mathrm{\textbf{(f) }} \left[\frac{47}{18}\;,\;\frac{-43}{18}\right]\,\text{\ding{55}}
\\[-0.4mm]
\romannumeral3 : \; &\textbf{S} 
 &&\mathrm{\textbf{(a) }} \frac{4xy+2y^2}{2\,\sqrt{2x^2y+2xy^2}}\,\text{\ding{55}}
 &&\mathrm{\textbf{(b) }} \left( \frac{10}{2\,\sqrt{12}} \;,\; \frac{16}{2\,\sqrt{12}} \right)\,\text{\ding{51}}
 &&\mathrm{\textbf{(c) }} \frac{x}{e^2} - y - z + 2 = 0\,\text{\ding{51}}
 &&\mathrm{\textbf{(d) }} 2.259\,\text{\ding{55}}
 &&\mathrm{\textbf{(e) }} \nicefrac{3}{5}\,\text{\ding{55}}
 &&\mathrm{\textbf{(f) }} \left[\frac{1}{4}\;,\;0\right]\,\text{\ding{51}}
\\[-0.4mm]
\romannumeral4 : \; &\textbf{T} 
 &&\mathrm{\textbf{(a) }} \frac{4xy+6y^2}{4\,\sqrt[4]{(2x^2y+6xy^2)^{3}}}\,\text{\ding{55}}
 &&\mathrm{\textbf{(b) }} \left( \frac{144}{4\,\sqrt[4]{360^3}} \;,\; \frac{162}{4\,\sqrt[4]{360^3}} \right)\,\text{\ding{51}}
 &&\mathrm{\textbf{(c) }} \frac{x}{e^3} - y - z + 3 = 0\,\text{\ding{51}}
 &&\mathrm{\textbf{(d) }} 2.969\,\text{\ding{55}}
 &&\mathrm{\textbf{(e) }} \nicefrac{20}{91}\,\text{\ding{51}}
 &&\mathrm{\textbf{(f) }} \left[\frac{-20}{7}\;,\;\frac{-22}{7}\right]\,\text{\ding{51}}
\end{aligned} $
\\[2mm]
\rule[1mm]{\linewidth}{0.5pt}
$\boxed{\bm{\xi}} \quad \begin{aligned}
\romannumeral1 : \; &\textbf{R} 
 &&\mathrm{\textbf{(a) }} \frac{12xy+8y^2}{3\,\sqrt[3]{(6x^2y+8xy^2)^{2}}}\,\text{\ding{51}}
 &&\mathrm{\textbf{(b) }} \left( \frac{108}{3\,\sqrt[3]{90^2}} \;,\; \frac{54}{3\,\sqrt[3]{90^2}} \right)\,\text{\ding{51}}
 &&\mathrm{\textbf{(c) }} \frac{x}{e^2} - y - z + 2 = 0\,\text{\ding{51}}
 &&\mathrm{\textbf{(d) }} 2.159\,\text{\ding{55}}
 &&\mathrm{\textbf{(e) }} \nicefrac{-32}{7}\,\text{\ding{51}}
 &&\mathrm{\textbf{(f) }} \left[\frac{33}{212}\;,\;\frac{-19}{212}\right]\,\text{\ding{55}}
\\[-0.4mm]
\romannumeral2 : \; &\textbf{O} 
 &&\mathrm{\textbf{(a) }} \frac{4xy+6y^2}{3\,\sqrt[3]{(2x^2y+6xy^2)^{2}}}\,\text{\ding{51}}
 &&\mathrm{\textbf{(b) }} \left( \frac{128}{3\,\sqrt[3]{224^2}} \;,\; \frac{104}{3\,\sqrt[3]{224^2}} \right)\,\text{\ding{55}}
 &&\mathrm{\textbf{(c) }} x+3z=0\,\text{\ding{51}}
 &&\mathrm{\textbf{(d) }} -0.003\,\text{\ding{55}}
 &&\mathrm{\textbf{(e) }} \nicefrac{54}{23}\,\text{\ding{51}}
 &&\mathrm{\textbf{(f) }} \left[\frac{-47}{12}\;,\;\frac{-1}{12}\right]\,\text{\ding{55}}
\\[-0.4mm]
\romannumeral3 : \; &\textbf{P} 
 &&\mathrm{\textbf{(a) }} \frac{2x^2+4xy}{3\,\sqrt[3]{(2x^2y+2xy^2)^{2}}}\,\text{\ding{51}}
 &&\mathrm{\textbf{(b) }} \left( \frac{18}{3\,\sqrt[3]{40^2}} \;,\; \frac{48}{3\,\sqrt[3]{40^2}} \right)\,\text{\ding{51}}
 &&\mathrm{\textbf{(c) }} \frac{x}{e^5} - y - z + 5 = 0\,\text{\ding{51}}
 &&\mathrm{\textbf{(d) }} 5.009\,\text{\ding{55}}
 &&\mathrm{\textbf{(e) }} \nicefrac{1}{2}\,\text{\ding{55}}
 &&\mathrm{\textbf{(f) }} \left[-2\;,\;-1\right]\,\text{\ding{51}}
\\[-0.4mm]
\romannumeral4 : \; &\textbf{A} 
 &&\mathrm{\textbf{(a) }} \frac{5x^2+2xy}{2\,\sqrt{5x^2y+xy^2}}\,\text{\ding{51}}
 &&\mathrm{\textbf{(b) }} \left( \frac{176}{2\,\sqrt{384}} \;,\; \frac{112}{2\,\sqrt{384}} \right)\,\text{\ding{55}}
 &&\mathrm{\textbf{(c) }} x+2z=0\,\text{\ding{55}}
 &&\mathrm{\textbf{(d) }} 0.1\,\text{\ding{55}}
 &&\mathrm{\textbf{(e) }} \nicefrac{-50}{7}\,\text{\ding{55}}
 &&\mathrm{\textbf{(f) }} \left[\frac{-11}{37}\;,\;\frac{4}{37}\right]\,\text{\ding{55}}
\end{aligned} $
\\[2mm]
\rule[1mm]{\linewidth}{0.5pt}
\end{center}\end{varwidth}\end{center}\clearpage\begin{center}{\huge \LARGE Funkcia dvoch premenných (riešenia)}\\[4mm]\begin{varwidth}{\linewidth}\begin{center}
\scriptsize
\rule[1mm]{\linewidth}{0.5pt}
$\boxed{\bm{\omicron}} \quad \begin{aligned}
\romannumeral1 : \; &\textbf{Ž} 
 &&\mathrm{\textbf{(a) }} \frac{4x^2+4xy}{4\,\sqrt[4]{(4x^2y+2xy^2)^{3}}}\,\text{\ding{55}}
 &&\mathrm{\textbf{(b) }} \left( \frac{56}{4\,\sqrt[4]{96^3}} \;,\; \frac{60}{4\,\sqrt[4]{96^3}} \right)\,\text{\ding{51}}
 &&\mathrm{\textbf{(c) }} x-z=0\,\text{\ding{51}}
 &&\mathrm{\textbf{(d) }} 0.03\,\text{\ding{51}}
 &&\mathrm{\textbf{(e) }} \nicefrac{-27}{10}\,\text{\ding{55}}
 &&\mathrm{\textbf{(f) }} \left[\frac{-41}{42}\;,\;\frac{-1}{21}\right]\,\text{\ding{51}}
\\[-0.4mm]
\romannumeral2 : \; &\textbf{A} 
 &&\mathrm{\textbf{(a) }} \frac{12xy+7y^2}{2\,\sqrt{6x^2y+7xy^2}}\,\text{\ding{51}}
 &&\mathrm{\textbf{(b) }} \left( \frac{76}{2\,\sqrt{104}} \;,\; \frac{80}{2\,\sqrt{104}} \right)\,\text{\ding{55}}
 &&\mathrm{\textbf{(c) }} x+z=0\,\text{\ding{55}}
 &&\mathrm{\textbf{(d) }} 0.01\,\text{\ding{55}}
 &&\mathrm{\textbf{(e) }} \nicefrac{6}{143}\,\text{\ding{55}}
 &&\mathrm{\textbf{(f) }} \left[\frac{-20}{13}\;,\;\frac{19}{13}\right]\,\text{\ding{55}}
\\[-0.4mm]
\romannumeral3 : \; &\textbf{B} 
 &&\mathrm{\textbf{(a) }} \frac{4xy+4y^2}{3\,\sqrt[3]{(2x^2y+4xy^2)^{2}}}\,\text{\ding{51}}
 &&\mathrm{\textbf{(b) }} \left( \frac{24}{3\,\sqrt[3]{20^2}} \;,\; \frac{18}{3\,\sqrt[3]{20^2}} \right)\,\text{\ding{51}}
 &&\mathrm{\textbf{(c) }} x+3z=0\,\text{\ding{55}}
 &&\mathrm{\textbf{(d) }} 0.033\,\text{\ding{55}}
 &&\mathrm{\textbf{(e) }} \nicefrac{100}{63}\,\text{\ding{55}}
 &&\mathrm{\textbf{(f) }} \left[\frac{1}{7}\;,\;\frac{6}{7}\right]\,\text{\ding{55}}
\\[-0.4mm]
\romannumeral4 : \; &\textbf{A} 
 &&\mathrm{\textbf{(a) }} \frac{7x^2+14xy}{3\,\sqrt[3]{(7x^2y+7xy^2)^{2}}}\,\text{\ding{51}}
 &&\mathrm{\textbf{(b) }} \left( \frac{49}{3\,\sqrt[3]{84^2}} \;,\; \frac{105}{3\,\sqrt[3]{84^2}} \right)\,\text{\ding{55}}
 &&\mathrm{\textbf{(c) }} \frac{x}{e^4} - y - z + 4 = 0\,\text{\ding{55}}
 &&\mathrm{\textbf{(d) }} 3.99\,\text{\ding{55}}
 &&\mathrm{\textbf{(e) }} \nicefrac{1}{4}\,\text{\ding{55}}
 &&\mathrm{\textbf{(f) }} \left[\frac{-1}{2}\;,\;\frac{1}{2}\right]\,\text{\ding{55}}
\end{aligned} $
\\[2mm]
\rule[1mm]{\linewidth}{0.5pt}
$\boxed{\bm{\pi}} \quad \begin{aligned}
\romannumeral1 : \; &\textbf{W} 
 &&\mathrm{\textbf{(a) }} \frac{18xy+8y^2}{2\,\sqrt{9x^2y+8xy^2}}\,\text{\ding{51}}
 &&\mathrm{\textbf{(b) }} \left( \frac{44}{2\,\sqrt{52}} \;,\; \frac{68}{2\,\sqrt{52}} \right)\,\text{\ding{51}}
 &&\mathrm{\textbf{(c) }} x - y - z + 0 = 0\,\text{\ding{51}}
 &&\mathrm{\textbf{(d) }} -0.4\,\text{\ding{51}}
 &&\mathrm{\textbf{(e) }} \nicefrac{3}{4}\,\text{\ding{51}}
 &&\mathrm{\textbf{(f) }} \left[\frac{1}{2}\;,\;\frac{-7}{2}\right]\,\text{\ding{55}}
\\[-0.4mm]
\romannumeral2 : \; &\textbf{I} 
 &&\mathrm{\textbf{(a) }} \frac{4xy+y^2}{2\,\sqrt{2x^2y+xy^2}}\,\text{\ding{55}}
 &&\mathrm{\textbf{(b) }} \left( \frac{80}{2\,\sqrt{192}} \;,\; \frac{64}{2\,\sqrt{192}} \right)\,\text{\ding{51}}
 &&\mathrm{\textbf{(c) }} x+3z=0\,\text{\ding{55}}
 &&\mathrm{\textbf{(d) }} -0.007\,\text{\ding{55}}
 &&\mathrm{\textbf{(e) }} -1\,\text{\ding{55}}
 &&\mathrm{\textbf{(f) }} \left[\frac{3}{7}\;,\;\frac{2}{7}\right]\,\text{\ding{51}}
\\[-0.4mm]
\romannumeral3 : \; &\textbf{F} 
 &&\mathrm{\textbf{(a) }} \frac{2xy+2y^2}{4\,\sqrt[4]{(x^2y+2xy^2)^{3}}}\,\text{\ding{51}}
 &&\mathrm{\textbf{(b) }} \left( \frac{4}{4\,\sqrt[4]{3^3}} \;,\; \frac{5}{4\,\sqrt[4]{3^3}} \right)\,\text{\ding{51}}
 &&\mathrm{\textbf{(c) }} x+3z=0\,\text{\ding{55}}
 &&\mathrm{\textbf{(d) }} 0.033\,\text{\ding{55}}
 &&\mathrm{\textbf{(e) }} \nicefrac{-1}{3}\,\text{\ding{55}}
 &&\mathrm{\textbf{(f) }} \left[\frac{32}{21}\;,\;\frac{17}{7}\right]\,\text{\ding{51}}
\\[-0.4mm]
\romannumeral4 : \; &\textbf{I} 
 &&\mathrm{\textbf{(a) }} \frac{16xy+y^2}{4\,\sqrt[4]{(8x^2y+xy^2)^{3}}}\,\text{\ding{55}}
 &&\mathrm{\textbf{(b) }} \left( \frac{49}{4\,\sqrt[4]{75^3}} \;,\; \frac{78}{4\,\sqrt[4]{75^3}} \right)\,\text{\ding{51}}
 &&\mathrm{\textbf{(c) }} x - y - z + 0 = 0\,\text{\ding{55}}
 &&\mathrm{\textbf{(d) }} 0.21\,\text{\ding{55}}
 &&\mathrm{\textbf{(e) }} \nicefrac{20}{19}\,\text{\ding{55}}
 &&\mathrm{\textbf{(f) }} \left[\frac{12}{7}\;,\;\frac{-2}{7}\right]\,\text{\ding{51}}
\end{aligned} $
\\[2mm]
\rule[1mm]{\linewidth}{0.5pt}
$\boxed{\bm{\rho}} \quad \begin{aligned}
\romannumeral1 : \; &\textbf{H} 
 &&\mathrm{\textbf{(a) }} \frac{16xy+y^2}{3\,\sqrt[3]{(8x^2y+xy^2)^{2}}}\,\text{\ding{51}}
 &&\mathrm{\textbf{(b) }} \left( \frac{36}{3\,\sqrt[3]{20^2}} \;,\; \frac{12}{3\,\sqrt[3]{20^2}} \right)\,\text{\ding{51}}
 &&\mathrm{\textbf{(c) }} \frac{x}{e^4} - y - z + 4 = 0\,\text{\ding{55}}
 &&\mathrm{\textbf{(d) }} 4.098\,\text{\ding{55}}
 &&\mathrm{\textbf{(e) }} \nicefrac{48}{11}\,\text{\ding{51}}
 &&\mathrm{\textbf{(f) }} \left[\frac{-17}{10}\;,\;\frac{13}{10}\right]\,\text{\ding{55}}
\\[-0.4mm]
\romannumeral2 : \; &\textbf{R} 
 &&\mathrm{\textbf{(a) }} \frac{2x^2+6xy}{3\,\sqrt[3]{(2x^2y+3xy^2)^{2}}}\,\text{\ding{51}}
 &&\mathrm{\textbf{(b) }} \left( \frac{51}{3\,\sqrt[3]{78^2}} \;,\; \frac{44}{3\,\sqrt[3]{78^2}} \right)\,\text{\ding{51}}
 &&\mathrm{\textbf{(c) }} x+2z=0\,\text{\ding{51}}
 &&\mathrm{\textbf{(d) }} 0.005\,\text{\ding{55}}
 &&\mathrm{\textbf{(e) }} \nicefrac{1}{23}\,\text{\ding{51}}
 &&\mathrm{\textbf{(f) }} \left[\frac{-5}{4}\;,\;\frac{23}{4}\right]\,\text{\ding{55}}
\\[-0.4mm]
\romannumeral3 : \; &\textbf{A} 
 &&\mathrm{\textbf{(a) }} \frac{x^2+4xy}{4\,\sqrt[4]{(x^2y+2xy^2)^{3}}}\,\text{\ding{51}}
 &&\mathrm{\textbf{(b) }} \left( \frac{24}{4\,\sqrt[4]{21^3}} \;,\; \frac{13}{4\,\sqrt[4]{21^3}} \right)\,\text{\ding{55}}
 &&\mathrm{\textbf{(c) }} \frac{x}{e} - y - z + 1 = 0\,\text{\ding{55}}
 &&\mathrm{\textbf{(d) }} 1.084\,\text{\ding{55}}
 &&\mathrm{\textbf{(e) }} \nicefrac{25}{16}\,\text{\ding{55}}
 &&\mathrm{\textbf{(f) }} \left[\frac{5}{8}\;,\;\frac{11}{8}\right]\,\text{\ding{55}}
\\[-0.4mm]
\romannumeral4 : \; &\textbf{D} 
 &&\mathrm{\textbf{(a) }} \frac{6xy+2y^2}{3\,\sqrt[3]{(3x^2y+2xy^2)^{2}}}\,\text{\ding{51}}
 &&\mathrm{\textbf{(b) }} \left( \frac{32}{3\,\sqrt[3]{40^2}} \;,\; \frac{28}{3\,\sqrt[3]{40^2}} \right)\,\text{\ding{55}}
 &&\mathrm{\textbf{(c) }} x - y - z + 0 = 0\,\text{\ding{55}}
 &&\mathrm{\textbf{(d) }} -0.28\,\text{\ding{55}}
 &&\mathrm{\textbf{(e) }} \nicefrac{-4}{31}\,\text{\ding{51}}
 &&\mathrm{\textbf{(f) }} \left[\frac{1}{2}\;,\;0\right]\,\text{\ding{51}}
\end{aligned} $
\\[2mm]
\rule[1mm]{\linewidth}{0.5pt}
$\boxed{\bm{\sigma}} \quad \begin{aligned}
\romannumeral1 : \; &\textbf{L} 
 &&\mathrm{\textbf{(a) }} \frac{14xy+2y^2}{2\,\sqrt{7x^2y+2xy^2}}\,\text{\ding{51}}
 &&\mathrm{\textbf{(b) }} \left( \frac{144}{2\,\sqrt{176}} \;,\; \frac{60}{2\,\sqrt{176}} \right)\,\text{\ding{51}}
 &&\mathrm{\textbf{(c) }} \frac{x}{e} - y - z + 1 = 0\,\text{\ding{51}}
 &&\mathrm{\textbf{(d) }} 0.966\,\text{\ding{55}}
 &&\mathrm{\textbf{(e) }} \nicefrac{2}{3}\,\text{\ding{55}}
 &&\mathrm{\textbf{(f) }} \left[\frac{7}{15}\;,\;\frac{-17}{15}\right]\,\text{\ding{55}}
\\[-0.4mm]
\romannumeral2 : \; &\textbf{A} 
 &&\mathrm{\textbf{(a) }} \frac{2xy+8y^2}{2\,\sqrt{x^2y+8xy^2}}\,\text{\ding{51}}
 &&\mathrm{\textbf{(b) }} \left( \frac{84}{2\,\sqrt{156}} \;,\; \frac{100}{2\,\sqrt{156}} \right)\,\text{\ding{55}}
 &&\mathrm{\textbf{(c) }} \frac{x}{e^2} - y - z + 2 = 0\,\text{\ding{55}}
 &&\mathrm{\textbf{(d) }} 2.097\,\text{\ding{55}}
 &&\mathrm{\textbf{(e) }} \nicefrac{18}{5}\,\text{\ding{55}}
 &&\mathrm{\textbf{(f) }} \left[\frac{1}{3}\;,\;\frac{11}{3}\right]\,\text{\ding{55}}
\\[-0.4mm]
\romannumeral3 : \; &\textbf{M} 
 &&\mathrm{\textbf{(a) }} \frac{6xy+5y^2}{2\,\sqrt{3x^2y+5xy^2}}\,\text{\ding{51}}
 &&\mathrm{\textbf{(b) }} \left( \frac{11}{2\,\sqrt{8}} \;,\; \frac{13}{2\,\sqrt{8}} \right)\,\text{\ding{55}}
 &&\mathrm{\textbf{(c) }} x+2z=0\,\text{\ding{51}}
 &&\mathrm{\textbf{(d) }} 0.1\,\text{\ding{55}}
 &&\mathrm{\textbf{(e) }} \nicefrac{-4}{3}\,\text{\ding{55}}
 &&\mathrm{\textbf{(f) }} \left[\frac{-15}{86}\;,\;\frac{13}{43}\right]\,\text{\ding{51}}
\\[-0.4mm]
\romannumeral4 : \; &\textbf{A} 
 &&\mathrm{\textbf{(a) }} \frac{8x^2+4xy}{2\,\sqrt{8x^2y+2xy^2}}\,\text{\ding{51}}
 &&\mathrm{\textbf{(b) }} \left( \frac{114}{2\,\sqrt{132}} \;,\; \frac{56}{2\,\sqrt{132}} \right)\,\text{\ding{55}}
 &&\mathrm{\textbf{(c) }} x+z=0\,\text{\ding{55}}
 &&\mathrm{\textbf{(d) }} 0.3\,\text{\ding{55}}
 &&\mathrm{\textbf{(e) }} \nicefrac{8}{7}\,\text{\ding{55}}
 &&\mathrm{\textbf{(f) }} \left[\frac{10}{23}\;,\;\frac{-6}{23}\right]\,\text{\ding{55}}
\end{aligned} $
\\[2mm]
\rule[1mm]{\linewidth}{0.5pt}
$\boxed{\bm{\tau}} \quad \begin{aligned}
\romannumeral1 : \; &\textbf{Ž} 
 &&\mathrm{\textbf{(a) }} \frac{x^2+4xy}{3\,\sqrt[3]{(x^2y+2xy^2)^{2}}}\,\text{\ding{55}}
 &&\mathrm{\textbf{(b) }} \left( \frac{10}{3\,\sqrt[3]{24^2}} \;,\; \frac{32}{3\,\sqrt[3]{24^2}} \right)\,\text{\ding{51}}
 &&\mathrm{\textbf{(c) }} x-2z=0\,\text{\ding{51}}
 &&\mathrm{\textbf{(d) }} -0.005\,\text{\ding{51}}
 &&\mathrm{\textbf{(e) }} \nicefrac{-16}{23}\,\text{\ding{55}}
 &&\mathrm{\textbf{(f) }} \left[\frac{17}{48}\;,\;\frac{-1}{16}\right]\,\text{\ding{51}}
\\[-0.4mm]
\romannumeral2 : \; &\textbf{A} 
 &&\mathrm{\textbf{(a) }} \frac{12xy+5y^2}{3\,\sqrt[3]{(6x^2y+5xy^2)^{2}}}\,\text{\ding{51}}
 &&\mathrm{\textbf{(b) }} \left( \frac{44}{3\,\sqrt[3]{32^2}} \;,\; \frac{26}{3\,\sqrt[3]{32^2}} \right)\,\text{\ding{55}}
 &&\mathrm{\textbf{(c) }} x-3z=0\,\text{\ding{55}}
 &&\mathrm{\textbf{(d) }} -0.007\,\text{\ding{55}}
 &&\mathrm{\textbf{(e) }} \nicefrac{25}{16}\,\text{\ding{55}}
 &&\mathrm{\textbf{(f) }} \left[\frac{17}{16}\;,\;\frac{-3}{16}\right]\,\text{\ding{55}}
\\[-0.4mm]
\romannumeral3 : \; &\textbf{B} 
 &&\mathrm{\textbf{(a) }} \frac{6xy+4y^2}{4\,\sqrt[4]{(3x^2y+4xy^2)^{3}}}\,\text{\ding{51}}
 &&\mathrm{\textbf{(b) }} \left( \frac{16}{4\,\sqrt[4]{20^3}} \;,\; \frac{28}{4\,\sqrt[4]{20^3}} \right)\,\text{\ding{51}}
 &&\mathrm{\textbf{(c) }} x-2z=0\,\text{\ding{55}}
 &&\mathrm{\textbf{(d) }} -0.005\,\text{\ding{55}}
 &&\mathrm{\textbf{(e) }} \nicefrac{1}{13}\,\text{\ding{55}}
 &&\mathrm{\textbf{(f) }} \left[\frac{8}{47}\;,\;\frac{7}{47}\right]\,\text{\ding{55}}
\\[-0.4mm]
\romannumeral4 : \; &\textbf{A} 
 &&\mathrm{\textbf{(a) }} \frac{18xy+2y^2}{3\,\sqrt[3]{(9x^2y+2xy^2)^{2}}}\,\text{\ding{51}}
 &&\mathrm{\textbf{(b) }} \left( \frac{80}{3\,\sqrt[3]{88^2}} \;,\; \frac{52}{3\,\sqrt[3]{88^2}} \right)\,\text{\ding{55}}
 &&\mathrm{\textbf{(c) }} x-3z=0\,\text{\ding{55}}
 &&\mathrm{\textbf{(d) }} -0.003\,\text{\ding{55}}
 &&\mathrm{\textbf{(e) }} \nicefrac{-12}{5}\,\text{\ding{55}}
 &&\mathrm{\textbf{(f) }} \left[\frac{-22}{35}\;,\;\frac{-4}{35}\right]\,\text{\ding{55}}
\end{aligned} $
\\[2mm]
\rule[1mm]{\linewidth}{0.5pt}
$\boxed{\bm{\upsilon}} \quad \begin{aligned}
\romannumeral1 : \; &\textbf{V} 
 &&\mathrm{\textbf{(a) }} \frac{2x^2+12xy}{4\,\sqrt[4]{(2x^2y+6xy^2)^{3}}}\,\text{\ding{51}}
 &&\mathrm{\textbf{(b) }} \left( \frac{18}{4\,\sqrt[4]{36^3}} \;,\; \frac{54}{4\,\sqrt[4]{36^3}} \right)\,\text{\ding{51}}
 &&\mathrm{\textbf{(c) }} x-2z=0\,\text{\ding{51}}
 &&\mathrm{\textbf{(d) }} -0.01\,\text{\ding{51}}
 &&\mathrm{\textbf{(e) }} \nicefrac{-1}{23}\,\text{\ding{55}}
 &&\mathrm{\textbf{(f) }} \left[\frac{19}{52}\;,\;\frac{7}{26}\right]\,\text{\ding{55}}
\\[-0.4mm]
\romannumeral2 : \; &\textbf{Z} 
 &&\mathrm{\textbf{(a) }} \frac{6xy+2y^2}{4\,\sqrt[4]{(3x^2y+2xy^2)^{3}}}\,\text{\ding{51}}
 &&\mathrm{\textbf{(b) }} \left( \frac{90}{4\,\sqrt[4]{216^3}} \;,\; \frac{96}{4\,\sqrt[4]{216^3}} \right)\,\text{\ding{55}}
 &&\mathrm{\textbf{(c) }} \frac{x}{e^4} - y - z + 4 = 0\,\text{\ding{51}}
 &&\mathrm{\textbf{(d) }} 3.902\,\text{\ding{51}}
 &&\mathrm{\textbf{(e) }} \nicefrac{-7}{96}\,\text{\ding{51}}
 &&\mathrm{\textbf{(f) }} \left[\frac{73}{30}\;,\;\frac{-47}{30}\right]\,\text{\ding{55}}
\\[-0.4mm]
\romannumeral3 : \; &\textbf{O} 
 &&\mathrm{\textbf{(a) }} \frac{6x^2+2xy}{2\,\sqrt{6x^2y+xy^2}}\,\text{\ding{51}}
 &&\mathrm{\textbf{(b) }} \left( \frac{25}{2\,\sqrt{26}} \;,\; \frac{28}{2\,\sqrt{26}} \right)\,\text{\ding{55}}
 &&\mathrm{\textbf{(c) }} x-3z=0\,\text{\ding{51}}
 &&\mathrm{\textbf{(d) }} -0.007\,\text{\ding{55}}
 &&\mathrm{\textbf{(e) }} \nicefrac{28}{83}\,\text{\ding{51}}
 &&\mathrm{\textbf{(f) }} \left[\frac{37}{18}\;,\;\frac{-17}{18}\right]\,\text{\ding{55}}
\\[-0.4mm]
\romannumeral4 : \; &\textbf{R} 
 &&\mathrm{\textbf{(a) }} \frac{2x^2+6xy}{3\,\sqrt[3]{(2x^2y+3xy^2)^{2}}}\,\text{\ding{51}}
 &&\mathrm{\textbf{(b) }} \left( \frac{15}{3\,\sqrt[3]{27^2}} \;,\; \frac{36}{3\,\sqrt[3]{27^2}} \right)\,\text{\ding{51}}
 &&\mathrm{\textbf{(c) }} x+z=0\,\text{\ding{51}}
 &&\mathrm{\textbf{(d) }} 0.3\,\text{\ding{55}}
 &&\mathrm{\textbf{(e) }} \nicefrac{5}{91}\,\text{\ding{51}}
 &&\mathrm{\textbf{(f) }} \left[\frac{-17}{59}\;,\;\frac{-8}{59}\right]\,\text{\ding{55}}
\end{aligned} $
\\[2mm]
\rule[1mm]{\linewidth}{0.5pt}
$\boxed{\bm{\phi}} \quad \begin{aligned}
\romannumeral1 : \; &\textbf{S} 
 &&\mathrm{\textbf{(a) }} \frac{2xy+y^2}{4\,\sqrt[4]{(x^2y+xy^2)^{3}}}\,\text{\ding{55}}
 &&\mathrm{\textbf{(b) }} \left( \frac{7}{4\,\sqrt[4]{12^3}} \;,\; \frac{15}{4\,\sqrt[4]{12^3}} \right)\,\text{\ding{51}}
 &&\mathrm{\textbf{(c) }} x+3z=0\,\text{\ding{51}}
 &&\mathrm{\textbf{(d) }} -0.007\,\text{\ding{55}}
 &&\mathrm{\textbf{(e) }} \nicefrac{-7}{38}\,\text{\ding{55}}
 &&\mathrm{\textbf{(f) }} \left[\frac{-19}{17}\;,\;\frac{15}{17}\right]\,\text{\ding{51}}
\\[-0.4mm]
\romannumeral2 : \; &\textbf{E} 
 &&\mathrm{\textbf{(a) }} \frac{9x^2+8xy}{2\,\sqrt{9x^2y+4xy^2}}\,\text{\ding{51}}
 &&\mathrm{\textbf{(b) }} \left( \frac{88}{2\,\sqrt{104}} \;,\; \frac{68}{2\,\sqrt{104}} \right)\,\text{\ding{55}}
 &&\mathrm{\textbf{(c) }} x+2z=0\,\text{\ding{55}}
 &&\mathrm{\textbf{(d) }} 0.05\,\text{\ding{55}}
 &&\mathrm{\textbf{(e) }} \nicefrac{-32}{7}\,\text{\ding{51}}
 &&\mathrm{\textbf{(f) }} \left[\frac{15}{32}\;,\;\frac{1}{16}\right]\,\text{\ding{55}}
\\[-0.4mm]
\romannumeral3 : \; &\textbf{N} 
 &&\mathrm{\textbf{(a) }} \frac{12xy+2y^2}{4\,\sqrt[4]{(6x^2y+2xy^2)^{3}}}\,\text{\ding{51}}
 &&\mathrm{\textbf{(b) }} \left( \frac{14}{4\,\sqrt[4]{8^3}} \;,\; \frac{10}{4\,\sqrt[4]{8^3}} \right)\,\text{\ding{55}}
 &&\mathrm{\textbf{(c) }} x+2z=0\,\text{\ding{51}}
 &&\mathrm{\textbf{(d) }} -0.15\,\text{\ding{55}}
 &&\mathrm{\textbf{(e) }} \nicefrac{81}{16}\,\text{\ding{51}}
 &&\mathrm{\textbf{(f) }} \left[\frac{-15}{26}\;,\;\frac{2}{13}\right]\,\text{\ding{51}}
\\[-0.4mm]
\romannumeral4 : \; &\textbf{O} 
 &&\mathrm{\textbf{(a) }} \frac{3x^2+2xy}{4\,\sqrt[4]{(3x^2y+xy^2)^{3}}}\,\text{\ding{51}}
 &&\mathrm{\textbf{(b) }} \left( \frac{40}{4\,\sqrt[4]{66^3}} \;,\; \frac{39}{4\,\sqrt[4]{66^3}} \right)\,\text{\ding{55}}
 &&\mathrm{\textbf{(c) }} \frac{x}{e^3} - y - z + 3 = 0\,\text{\ding{51}}
 &&\mathrm{\textbf{(d) }} 2.98\,\text{\ding{55}}
 &&\mathrm{\textbf{(e) }} \nicefrac{-24}{11}\,\text{\ding{51}}
 &&\mathrm{\textbf{(f) }} \left[\frac{-23}{10}\;,\;\frac{7}{10}\right]\,\text{\ding{55}}
\end{aligned} $
\\[2mm]
\rule[1mm]{\linewidth}{0.5pt}
\end{center}\end{varwidth}\end{center}\clearpage\begin{center}{\huge \LARGE Funkcia dvoch premenných (riešenia)}\\[4mm]\begin{varwidth}{\linewidth}\begin{center}
\scriptsize
\rule[1mm]{\linewidth}{0.5pt}
$\boxed{\bm{\chi}} \quad \begin{aligned}
\romannumeral1 : \; &\textbf{V} 
 &&\mathrm{\textbf{(a) }} \frac{12xy+7y^2}{4\,\sqrt[4]{(6x^2y+7xy^2)^{3}}}\,\text{\ding{51}}
 &&\mathrm{\textbf{(b) }} \left( \frac{171}{4\,\sqrt[4]{351^3}} \;,\; \frac{180}{4\,\sqrt[4]{351^3}} \right)\,\text{\ding{51}}
 &&\mathrm{\textbf{(c) }} \frac{x}{e^5} - y - z + 5 = 0\,\text{\ding{51}}
 &&\mathrm{\textbf{(d) }} 4.701\,\text{\ding{51}}
 &&\mathrm{\textbf{(e) }} \nicefrac{63}{136}\,\text{\ding{55}}
 &&\mathrm{\textbf{(f) }} \left[\frac{27}{20}\;,\;\frac{-13}{20}\right]\,\text{\ding{55}}
\\[-0.4mm]
\romannumeral2 : \; &\textbf{L} 
 &&\mathrm{\textbf{(a) }} \frac{4xy+4y^2}{4\,\sqrt[4]{(2x^2y+4xy^2)^{3}}}\,\text{\ding{51}}
 &&\mathrm{\textbf{(b) }} \left( \frac{8}{4\,\sqrt[4]{6^3}} \;,\; \frac{10}{4\,\sqrt[4]{6^3}} \right)\,\text{\ding{51}}
 &&\mathrm{\textbf{(c) }} x-2z=0\,\text{\ding{51}}
 &&\mathrm{\textbf{(d) }} 0.005\,\text{\ding{55}}
 &&\mathrm{\textbf{(e) }} \nicefrac{12}{11}\,\text{\ding{55}}
 &&\mathrm{\textbf{(f) }} \left[\frac{13}{22}\;,\;\frac{9}{22}\right]\,\text{\ding{55}}
\\[-0.4mm]
\romannumeral3 : \; &\textbf{C} 
 &&\mathrm{\textbf{(a) }} \frac{12xy+5y^2}{4\,\sqrt[4]{(6x^2y+5xy^2)^{3}}}\,\text{\ding{51}}
 &&\mathrm{\textbf{(b) }} \left( \frac{128}{4\,\sqrt[4]{104^3}} \;,\; \frac{46}{4\,\sqrt[4]{104^3}} \right)\,\text{\ding{55}}
 &&\mathrm{\textbf{(c) }} x - y - z + 0 = 0\,\text{\ding{55}}
 &&\mathrm{\textbf{(d) }} 0.23\,\text{\ding{55}}
 &&\mathrm{\textbf{(e) }} \nicefrac{24}{71}\,\text{\ding{55}}
 &&\mathrm{\textbf{(f) }} \left[-4\;,\;2\right]\,\text{\ding{51}}
\\[-0.4mm]
\romannumeral4 : \; &\textbf{I} 
 &&\mathrm{\textbf{(a) }} \frac{4xy+2y^2}{4\,\sqrt[4]{(2x^2y+2xy^2)^{3}}}\,\text{\ding{55}}
 &&\mathrm{\textbf{(b) }} \left( \frac{18}{4\,\sqrt[4]{40^3}} \;,\; \frac{48}{4\,\sqrt[4]{40^3}} \right)\,\text{\ding{51}}
 &&\mathrm{\textbf{(c) }} \frac{x}{e} - y - z + 1 = 0\,\text{\ding{55}}
 &&\mathrm{\textbf{(d) }} 0.91\,\text{\ding{55}}
 &&\mathrm{\textbf{(e) }} \nicefrac{-6}{13}\,\text{\ding{55}}
 &&\mathrm{\textbf{(f) }} \left[\frac{7}{16}\;,\;\frac{25}{16}\right]\,\text{\ding{51}}
\end{aligned} $
\\[2mm]
\rule[1mm]{\linewidth}{0.5pt}
$\boxed{\bm{\psi}} \quad \begin{aligned}
\romannumeral1 : \; &\textbf{T} 
 &&\mathrm{\textbf{(a) }} \frac{7x^2+14xy}{2\,\sqrt{7x^2y+7xy^2}}\,\text{\ding{55}}
 &&\mathrm{\textbf{(b) }} \left( \frac{147}{2\,\sqrt{210}} \;,\; \frac{112}{2\,\sqrt{210}} \right)\,\text{\ding{51}}
 &&\mathrm{\textbf{(c) }} x+3z=0\,\text{\ding{51}}
 &&\mathrm{\textbf{(d) }} -0.007\,\text{\ding{55}}
 &&\mathrm{\textbf{(e) }} \nicefrac{-6}{11}\,\text{\ding{51}}
 &&\mathrm{\textbf{(f) }} \left[\frac{27}{10}\;,\;\frac{-13}{10}\right]\,\text{\ding{51}}
\\[-0.4mm]
\romannumeral2 : \; &\textbf{Ó} 
 &&\mathrm{\textbf{(a) }} \frac{2xy+5y^2}{2\,\sqrt{x^2y+5xy^2}}\,\text{\ding{55}}
 &&\mathrm{\textbf{(b) }} \left( \frac{104}{2\,\sqrt{276}} \;,\; \frac{129}{2\,\sqrt{276}} \right)\,\text{\ding{51}}
 &&\mathrm{\textbf{(c) }} \frac{x}{e^3} - y - z + 3 = 0\,\text{\ding{55}}
 &&\mathrm{\textbf{(d) }} 2.805\,\text{\ding{51}}
 &&\mathrm{\textbf{(e) }} \nicefrac{27}{47}\,\text{\ding{55}}
 &&\mathrm{\textbf{(f) }} \left[\frac{-8}{37}\;,\;\frac{3}{37}\right]\,\text{\ding{51}}
\\[-0.4mm]
\romannumeral3 : \; &\textbf{G} 
 &&\mathrm{\textbf{(a) }} \frac{8x^2+8xy}{4\,\sqrt[4]{(8x^2y+4xy^2)^{3}}}\,\text{\ding{51}}
 &&\mathrm{\textbf{(b) }} \left( \frac{84}{4\,\sqrt[4]{60^3}} \;,\; \frac{32}{4\,\sqrt[4]{60^3}} \right)\,\text{\ding{51}}
 &&\mathrm{\textbf{(c) }} x+3z=0\,\text{\ding{55}}
 &&\mathrm{\textbf{(d) }} -0.003\,\text{\ding{55}}
 &&\mathrm{\textbf{(e) }} \nicefrac{-80}{59}\,\text{\ding{51}}
 &&\mathrm{\textbf{(f) }} \left[\frac{-13}{14}\;,\;\frac{1}{14}\right]\,\text{\ding{51}}
\\[-0.4mm]
\romannumeral4 : \; &\textbf{A} 
 &&\mathrm{\textbf{(a) }} \frac{2xy+6y^2}{3\,\sqrt[3]{(x^2y+6xy^2)^{2}}}\,\text{\ding{51}}
 &&\mathrm{\textbf{(b) }} \left( \frac{8}{3\,\sqrt[3]{7^2}} \;,\; \frac{13}{3\,\sqrt[3]{7^2}} \right)\,\text{\ding{55}}
 &&\mathrm{\textbf{(c) }} x+3z=0\,\text{\ding{55}}
 &&\mathrm{\textbf{(d) }} -0.007\,\text{\ding{55}}
 &&\mathrm{\textbf{(e) }} \nicefrac{4}{23}\,\text{\ding{55}}
 &&\mathrm{\textbf{(f) }} \left[\frac{-7}{2}\;,\;\frac{-11}{2}\right]\,\text{\ding{55}}
\end{aligned} $
\\[2mm]
\rule[1mm]{\linewidth}{0.5pt}
$\boxed{\bm{\omega}} \quad \begin{aligned}
\romannumeral1 : \; &\textbf{J} 
 &&\mathrm{\textbf{(a) }} \frac{9x^2+2xy}{2\,\sqrt{9x^2y+xy^2}}\,\text{\ding{55}}
 &&\mathrm{\textbf{(b) }} \left( \frac{37}{2\,\sqrt{38}} \;,\; \frac{40}{2\,\sqrt{38}} \right)\,\text{\ding{51}}
 &&\mathrm{\textbf{(c) }} x - y - z + 0 = 0\,\text{\ding{55}}
 &&\mathrm{\textbf{(d) }} 0.28\,\text{\ding{55}}
 &&\mathrm{\textbf{(e) }} \nicefrac{25}{8}\,\text{\ding{51}}
 &&\mathrm{\textbf{(f) }} \left[\frac{25}{54}\;,\;\frac{-2}{27}\right]\,\text{\ding{51}}
\\[-0.4mm]
\romannumeral2 : \; &\textbf{Ó} 
 &&\mathrm{\textbf{(a) }} \frac{4x^2+6xy}{2\,\sqrt{4x^2y+3xy^2}}\,\text{\ding{55}}
 &&\mathrm{\textbf{(b) }} \left( \frac{144}{2\,\sqrt{288}} \;,\; \frac{108}{2\,\sqrt{288}} \right)\,\text{\ding{51}}
 &&\mathrm{\textbf{(c) }} x - y - z + 0 = 0\,\text{\ding{55}}
 &&\mathrm{\textbf{(d) }} -0.31\,\text{\ding{51}}
 &&\mathrm{\textbf{(e) }} \nicefrac{6}{5}\,\text{\ding{55}}
 &&\mathrm{\textbf{(f) }} \left[\frac{-1}{5}\;,\;\frac{1}{5}\right]\,\text{\ding{51}}
\\[-0.4mm]
\romannumeral3 : \; &\textbf{G} 
 &&\mathrm{\textbf{(a) }} \frac{4x^2+2xy}{3\,\sqrt[3]{(4x^2y+xy^2)^{2}}}\,\text{\ding{51}}
 &&\mathrm{\textbf{(b) }} \left( \frac{105}{3\,\sqrt[3]{228^2}} \;,\; \frac{88}{3\,\sqrt[3]{228^2}} \right)\,\text{\ding{51}}
 &&\mathrm{\textbf{(c) }} \frac{x}{e^2} - y - z + 2 = 0\,\text{\ding{55}}
 &&\mathrm{\textbf{(d) }} 2.007\,\text{\ding{55}}
 &&\mathrm{\textbf{(e) }} \nicefrac{3}{2}\,\text{\ding{51}}
 &&\mathrm{\textbf{(f) }} \left[\frac{1}{2}\;,\;\frac{-15}{2}\right]\,\text{\ding{51}}
\\[-0.4mm]
\romannumeral4 : \; &\textbf{A} 
 &&\mathrm{\textbf{(a) }} \frac{4xy+6y^2}{4\,\sqrt[4]{(2x^2y+6xy^2)^{3}}}\,\text{\ding{51}}
 &&\mathrm{\textbf{(b) }} \left( \frac{40}{4\,\sqrt[4]{64^3}} \;,\; \frac{56}{4\,\sqrt[4]{64^3}} \right)\,\text{\ding{55}}
 &&\mathrm{\textbf{(c) }} x+z=0\,\text{\ding{55}}
 &&\mathrm{\textbf{(d) }} -0.02\,\text{\ding{55}}
 &&\mathrm{\textbf{(e) }} \nicefrac{15}{4}\,\text{\ding{55}}
 &&\mathrm{\textbf{(f) }} \left[\frac{-27}{14}\;,\;\frac{29}{14}\right]\,\text{\ding{55}}
\end{aligned} $
\\[1mm]
\rule[-1mm]{\linewidth}{0.5pt}
\end{center}
\end{varwidth}
\end{center}
\end{landscape}
\end{document}
